\documentclass[11pt]{article}

\usepackage[a4paper,margin=1in]{geometry}
\usepackage{amsmath,amssymb,amsthm,mathtools}
\usepackage{booktabs,array,longtable}
\usepackage{enumitem}
\usepackage{microtype}
\usepackage{xcolor}
\usepackage{hyperref}
\usepackage[nameinlink,capitalise,noabbrev]{cleveref}

\hypersetup{
  colorlinks=true,
  linkcolor=blue!55!black,
  citecolor=blue!55!black,
  urlcolor=blue!55!black
}

\newtheorem{theorem}{Theorem}[section]
\newtheorem{proposition}[theorem]{Proposition}
\newtheorem{lemma}[theorem]{Lemma}
\newtheorem{corollary}[theorem]{Corollary}
\newtheorem{definition}[theorem]{Definition}
\newtheorem{remark}[theorem]{Remark}

% The environments above share the theorem counter, so cleveref would name
% them all "Theorem"; these aliases restore the correct names.
\usepackage{etoolbox}
\AtBeginEnvironment{proposition}{\crefalias{theorem}{proposition}}
\AtBeginEnvironment{lemma}{\crefalias{theorem}{lemma}}
\AtBeginEnvironment{corollary}{\crefalias{theorem}{corollary}}
\AtBeginEnvironment{definition}{\crefalias{theorem}{definition}}
\AtBeginEnvironment{remark}{\crefalias{theorem}{remark}}

\numberwithin{equation}{section}

\newcommand{\one}{\mathbf 1}
\newcommand{\R}{\mathbb R}
\newcommand{\ip}[2]{\left\langle #1,#2\right\rangle}
\newcommand{\norm}[1]{\left\lVert #1\right\rVert}
\newcommand{\Tr}{\operatorname{Tr}}
\newcommand{\ord}{\operatorname{ord}}
\newcommand{\pos}[1]{\left(#1\right)_{+}}
\newcommand{\cH}{\mathcal H}
\newcommand{\dd}{\,\mathrm d}
\newcommand{\E}{\mathbb E}
\newcommand{\GammaLaw}{\operatorname{Gamma}}
\newcommand{\BetaLaw}{\operatorname{Beta}}
\newcommand{\supp}{\operatorname{supp}}

\newcommand*{\tk}[1]{\textcolor{blue}{\textbf{[TK: #1]}}}

\title{A Complete Proof of the Goodman--Style Lower Bound\\
for Odd-Cycle Densities in Graphons}
\author{Consolidated working note for H.~Yang and collaborators}
\date{July 2026}


\begin{document}
\maketitle

\begin{abstract}
  Let $W:[0,1]^2\to[0,1]$ be a graphon of edge density $p$, and let
  $m\ge3$ be odd.  We prove the conjectured Goodman--style inequality
  \[
    t(C_m,W)\ge p^m-p(1-p)^{m-1}
  \]
  for every $p\in[0,1]$.  The proof is divided at $p=2/3$, or equivalently
  at the complement density $q=1/3$.

  For $p\ge2/3$, a two-sided Schur-complement identity cancels the unknown
  odd moments of the compressed complement operator.  Expansion in complete
  homogeneous symmetric polynomials is then diagonalized by a Dirichlet
  average, converted to a beta integral, and smoothed by a Laplace transform
  into expectations of a shifted gamma random variable.  The remaining
  positivity follows from a uniform odd-to-even gamma moment inequality.

  For $1/2<p<2/3$, the Hilbert--Schmidt bound allows at most one compression
  eigenvalue above $q=1-p$.  Isolating this frontier eigenvalue and retaining
  both its direct coupling and the coupling forced into the safe spectral
  subspace leads to a one-dimensional Huber-type envelope.  Two canonical
  dual witnesses settle its quadratic and linear branches.  The only finite
  residue is a pair of exact univariate Bernstein certificates at $m=9$.
  The cycles $C_3,C_5,C_7$ are also given separate elementary proofs.

  A final section extracts the reusable mechanisms: one-step Jacobi
  coefficient stripping, first-return and cyclic excursion expansions,
  Dirichlet probabilistic polarization, beta--gamma/Laplace positivity
  transfer, the Laguerre differential equation hidden in the gamma moment
  lemma, and the one-dangerous-mode/Huber principle.  It also explains why the
  same threshold $q=1/3$ emerges independently from the two sides.
\end{abstract}

\tableofcontents

\section{The theorem, notation, and proof map}\label{sec:intro}

For a finite graph $F$ and a graphon $W$, write
\[
  t(F,W)=\int_{[0,1]^{V(F)}}
  \prod_{ij\in E(F)}W(x_i,x_j)
  \prod_{i\in V(F)}\dd x_i.
\]
In particular,
\[
  p=t(K_2,W)=\int_{[0,1]^2}W(x,y)\dd x\dd y
\]
is the edge density.  We prove the following theorem.

\begin{theorem}[Odd-cycle Goodman bound]\label{thm:main}
  Let $W$ be a graphon of edge density $p\in[0,1]$.  For every odd integer
  $m\ge3$,
  \begin{equation}\label{eq:main-goal}
    t(C_m,W)\ge p^m-p(1-p)^{m-1}.
  \end{equation}
\end{theorem}

Put
\[
  U=1-W,\qquad q=t(K_2,U)=1-p.
\]
If $p\le1/2$, then
\[
  p^m-pq^{m-1}=p\bigl(p^{m-1}-q^{m-1}\bigr)\le0,
\]
whereas $t(C_m,W)\ge0$.  Thus only $p>1/2$, or $q<1/2$, needs proof.
The nontrivial range splits into
\[
  \underbrace{0\le q\le\frac{1}{3}}_{p\ge2/3}
  \qquad\text{and}\qquad
  \underbrace{\frac{1}{3}<q<\frac{1}{2}}_{1/2<p<2/3}.
\]

\begin{corollary} \label{cor:main}
  For every graphon $W$ of edge density $p = 1 - 1 / k$ for some integer $k \ge 2$, and for every odd integer $m \ge 3$, the balanced complete $k$-partite graphon 
  \[
    W_k(x, y) = \begin{cases}
      1 & \text{if } \lfloor kx \rfloor \neq \lfloor ky \rfloor,
      \\
      0 & \text{otherwise}.
    \end{cases}
  \]
  is the minimiser of $t(C_m, W)$ among all graphons of edge density $p$.
\end{corollary}

\subsection{Proof map}

Every part of the proof works with the same few quantities, so we introduce
them first.  Let $T_U$ be the integral operator with kernel $U=1-W$, and
split it against the constant function $\one$.  The constant direction
contributes the scalar $q$, and the remainder is the mean-zero function
\[
  g=T_U\one-q\one
  =p-d_W,\qquad d_W(x)=\int_0^1W(x,y)\dd y,
\]
that is, the mean $p$ minus the degree function of $W$.  We also restrict
$T_U$ to the mean-zero subspace $\one^\perp$, and write $A$ for the
resulting self-adjoint operator, so that $\ip f{Ah}=\ip f{T_Uh}$ for all
$f,h\perp\one$.

Throughout the paper, $\lambda_1,\lambda_2,\ldots$ denote the eigenvalues of
$A$, and $P_i$ the orthogonal projection onto the eigenspace of
$\lambda_i$.  We write
\[
  \mu=\sum_i\norm{P_ig}^2\delta_{\lambda_i}
\]
for the associated measure on the spectrum of $A$, so that
$\int\lambda^j\dd\mu(\lambda)=\ip g{A^jg}$ for every $j\ge0$ and
$\mu(\R)=\norm g^2$.  Two elementary lemmas bound these eigenvalues.  The
first,
\[
  |\lambda_i|\le\frac{1}{2}\quad\text{for every }i \qquad(\Cref{lem:half-norm}),
\]
is used from Part~2 on; the second,
\[
  \sum_i\lambda_i^2+2\norm g^2=\Tr(A^2)+2\norm g^2\le pq \qquad(\Cref{lem:plt23-HS-bound}),
\]
only in Part~3.

\begin{enumerate}[label=\textbf{Part \arabic*.},leftmargin=3.4em]
  \item \Cref{sec:short-cycles} considers $m=3,5,7$. We prove these cases so that when we consider the general odd $m$ in the later parts, we can assume $m\ge9$.
  The case $m=3$ is Goodman's original bound, in graphon form.  For
  $m=5$ and $m=7$, the inclusion--exclusion principle applied to
  $\prod_{i=1}^m(1-U_i)$ expands $t(C_m,W)$ over the subsets of the
  $m$ edges.  A proper subset selects a disjoint union of paths, and
  therefore contributes a product of the path densities
  \[
    x_j=t(P_j,U)=\ip\one{T_U^j\one},
  \]
  while the full subset contributes $t(C_m,U)$, for which
  $t(C_m,U)\le x_{m-1}$ because $0\le U\le1$.  Each $x_j$ can in turn
  be written as a polynomial in $q$ and the moments
  \[
    s_j=\ip g{A^jg}=\int\lambda^j\dd\mu(\lambda).
  \]
  For $m=5$, the quantity to be shown nonnegative is a quadratic form
  in $g$ and $T_Ug$, which \Cref{subsec:C5} writes as a sum of two
  squares.  For $m=7$ it becomes
  \[
    6s_1^2+\int R_{q,r}(\lambda)\dd\mu(\lambda),
    \qquad r=\norm g^2,
  \]
  where $R_{q,r}$ is a quartic in $\lambda$; \Cref{subsec:C7} writes it
  as an sum of squares, nonnegative for every real $\lambda$.

  \item \Cref{sec:common-operator,sec:pge23} prove all odd lengths
  when $q\le1/3$.  We first show that the logarithm of a determinant
  generates the cycle densities: for any kernel $K$,
  \[
    -\log\det(I-zT_K)=\sum_{j\ge1}\frac{\Tr(T_K^j)}{j}z^j,
    \qquad
    \Tr(T_K^j)=t(C_j,K)\quad(j\ge2).
  \]  
  We apply this to $U$ and to $W$, and factorize the two determinants by
  Schur complementation against the constant direction.  This produces two
  further generating functions $L_+(z)$ and $L_-(z)$, whose form is
  recorded in \eqref{eq:pge23-Lminus}--\eqref{eq:pge23-Lplus}, and
  \begin{align*}
    t(C_m,U)&=\phantom{-}\Tr(A^m)+q^m+m[z^m]L_+(z),\\
    t(C_m,W)&=-\Tr(A^m)+p^m+m[z^m]L_-(z).
  \end{align*}
  The two traces have opposite signs because the mean-zero part of $T_U$
  is $A$, that of $T_W=T_1-T_U$ is $-A$, and $m$ is odd.  Since neither
  lemma above controls the odd power sum $\Tr(A^m)=\sum_i\lambda_i^m$, we
  add the two identities to remove it (\Cref{lem:pge23-two-sided-identity}):
  \[
    t(C_m,W)+t(C_m,U)=p^m+q^m+S_m,
    \qquad
    S_m=m[z^m]\bigl(L_-(z)+L_+(z)\bigr).
  \]
  Bounding $t(C_m,U)$ by the path density $x_{m-1}$, everything reduces to
  the nonnegativity of
  \[
    \Phi_m=q^{m-1}+S_m-x_{m-1}.
  \]

  Expanding $L_\pm$ and $x_{m-1}$ gives $\Phi_m$ as an integral against
  $\dd\mu^{\otimes r}$ of a symmetric polynomial
  (\Cref{prop:pge23-expansion}): with $n=m-2r$,
  \[
    \Phi_m=\sum_{r=1}^{(m-1)/2}\int\!\cdots\!\int
    P_{m,r}(q;\lambda_1,\ldots,\lambda_r)
    \prod_{i=1}^r\dd\mu(\lambda_i),
  \]
  \[
    P_{m,r}=\frac {m}{r}\Bigl[
    h_n(p^{\times r},-\lambda_1,\ldots,-\lambda_r)
    +h_n(q^{\times r},\lambda_1,\ldots,\lambda_r)\Bigr]
    -h_{n-1}(q^{\times(r+1)},\lambda_1,\ldots,\lambda_r),
  \]
  where $h_d$ is the complete homogeneous symmetric polynomial of degree
  $d$ and $a^{\times k}$ means $a$ repeated $k$ times.  It therefore
  suffices to prove $P_{m,r}\ge0$ on $[-1/2,1/2]^r$.

  That polynomial is then simplified three times until it can be shown nonnegative with elementary methods.
  First, using the fact that a complete homogeneous
  polynomial is an expectation of a power of a random convex combination,
  so with $\Theta$ uniform on the simplex $\Delta^r$ and
  $\widetilde P_{m,r}(q,\ell)=P_{m,r}(q;\ell,\ldots,\ell)$
  (\Cref{lem:pge23-dirichlet}),
  \[
    P_{m,r}(q;\lambda_1,\ldots,\lambda_r)
    =\E\,\widetilde P_{m,r}
    \Bigl(q,\sum_{i=1}^r\Theta_i\lambda_i\Bigr).
  \]
  Every such convex combination again lies in $[-1/2,1/2]$, so the $r$
  variables collapse to one and it suffices to prove
  $\widetilde P_{m,r}(q,\ell)\ge0$ for $\ell\in[-1/2,1/2]$.  Second, the
  repeated variable turns the diagonal polynomial into a beta integral,
  which after a projective change of variable reads, for $\ell>0$
  (\Cref{lem:pge23-beta}),
  \begin{align*}
    \widetilde P_{m,r}(q,\ell) &=\frac{\Gamma(m)}{r\Gamma(n)\Gamma(r)^2}\,\ell^{n+r}
    \int_0^\infty\frac{s^{r-1}}{(\ell+s)^m}
    \rho_{n,m}(q+s)\dd s,
    \\    
    \rho_{n,m}(u)&=\frac {m}{n}\bigl(u^n+(1-u)^n\bigr)-u^{n-1}.
  \end{align*}
  The case $\ell\le0$ is proved separately in
  \Cref{lem:pge23-ell-negative}, where the argument of $\rho_{n,m}$ stays
  at most $q\le1/2$ and $\rho_{n,m}$ is nonnegative there.  As
  $\rho_{n,m}$ is not
  nonnegative at every real number, the integrand cannot be handled
  pointwise, and we estimate the average instead.  Third, the Laplace
  representation of
  $(\ell+s)^{-m}$ replaces that average by an expectation over a gamma
  variable $Y\sim\GammaLaw(r,1)$ (\Cref{lem:pge23-laplace-gamma}):
  \[
    \widetilde P_{m,r}(q,\ell)
    =\frac{\ell^{n+r}}{r\Gamma(n)\Gamma(r)}
    \int_0^\infty t^{m-r-1}e^{-\ell t}\,
    \E\,\rho_{n,m}\Bigl(q+\frac {Y}{t}\Bigr)\dd t.
  \]
  As all the remaining terms are nonnegative, it is enough to prove that
  $\E\rho_{n,m}(q+xY)\ge0$ for every $x\ge0$.  Centering at $u=1/2$ and
  using that $n$ is odd, this reduces to a single inequality comparing
  consecutive shifted gamma moments,
  \[
    3j\,\E(zY-1)^{2j-1}\le(r+j)\,\E(zY-1)^{2j}
    \qquad(z\ge0),
  \]
  proved in \Cref{lem:pge23-gamma-moment} from a Stein identity and the
  differential equation it induces.  The constant $3$ is where the
  assumption $p\ge2/3$ is used.

  \item \Cref{sec:plt23}
  proves all odd $m\ge9$ when $1/2<p<2/3$.  The argument in this part is a
  spectral analysis of $A$.  We use two results, the identity for $W$ from
  Part~2,
  \[
    t(C_m,W)=p^m-\Tr(A^m)+m[z^m]L_-(z),
  \]
  and the nonnegativity of every coefficient of $L_-$; \Cref{lem:plt23-one-sided-identity}
  states both.  Nothing cancels $\Tr(A^m)$ this time, so it must be bounded
  directly.

  We first consider the eigenvalues of $A$ larger than $q$.  If there were
  two of them, then $\sum_i\lambda_i^2>2q^2$, whereas
  $\sum_i\lambda_i^2\le pq$ by \Cref{lem:plt23-HS-bound}, a contradiction because $2q^2 > pq$ when $q>1/3$.
  Hence at most one eigenvalue exceeds $q$.  
  If none does the theorem is immediate (\Cref{lem:plt23-eigenvalues-above-q}), 
  so we may assume that exactly one does and write $\Lambda$ for it.  
  The same bound gives
  $|\lambda_i|\le M:=\sqrt{pq-\Lambda^2}$ for every other eigenvalue.  
  One can also derive an upper bound on $\Lambda$ by using its eigenfunction $\phi$: 
  since $U\ge0$, replacing $\phi$ by $|\phi|$ can only increase
  $\ip\phi{T_U\phi}$, and expanding $|\phi|$ along $\one$ and $\phi$ gives
  \[
    \norm g^2\ge\frac{\Lambda(\Lambda-q)^2}{2(1-2\Lambda)},
  \]
  which with $2\norm g^2\le pq-\Lambda^2$, again from \Cref{lem:plt23-HS-bound},
  gives $\Lambda^2+q\Lambda-q\le0$ (\Cref{lem:plt23-forced-variance}).

  Decomposing $g=\gamma\phi+g_s$ into its component along $\phi$ and its
  component in the remaining spectral subspace, the two nonconstant terms
  of the identity can be bounded.  For $\Tr(A^m)$, the eigenvalues other
  than $\Lambda$ satisfy $\sum\lambda_i^m\le M^{m-2}\sum\lambda_i^2$, so
  \Cref{lem:plt23-HS-bound} gives
  \[
    \Tr(A^m)\le\Lambda^m+M^m-2M^{m-2}\bigl(\gamma^2+\norm{g_s}^2\bigr).
  \]
  For the generating function, keeping only the first term of
  $-\log(1-X)=\sum_{k\ge1}X^k/k$ discards
  nonnegative coefficients, and leaves
  \[
    m[z^m]L_-(z)\ge m\int k_m\dd\mu
    \ge mk_m(\Lambda)\gamma^2+mk_m(M)\norm{g_s}^2,
    \qquad
    k_m(\lambda)=\frac{p^{m-1}-\lambda^{m-1}}{p+\lambda}.
  \]
  Together they give (\Cref{prop:plt23-splitting-lower-bound})
  \[
    t(C_m,W)-\bigl(p^m-pq^{m-1}\bigr)
    \ge-R_m+\Lambda_m\gamma^2+M_m\norm{g_s}^2,
  \]
  where $R_m=\Lambda^m+M^m-pq^{m-1}$, $\Lambda_m=2M^{m-2}+mk_m(\Lambda)$ and
  $M_m=2M^{m-2}+mk_m(M)$, with $0<\Lambda_m<M_m$ because $k_m$ decreases and
  $M<\Lambda$.  It therefore suffices to prove
  $\Lambda_m\gamma^2+M_m\norm{g_s}^2\ge R_m$.

  To prove this lower bound, we first derive bounds on $\gamma$ and
  $\norm{g_s}$ from the eigenfunction $\phi$.  Decompose
  \[
    |\phi|=a_\phi\one+b\phi+k,
    \qquad a_\phi=\int_0^1|\phi|,
    \qquad k\perp\one,\phi,
  \]
  and set $K=\norm k^2$, so that $a_\phi^2+b^2+K=1$.  One can then show
  (\Cref{lem:plt23-gamma-upper-bound,lem:plt23-gs-lower-bound}) that
  \[
    \gamma\le\frac{a_\phi^2-2\Lambda(1+b)}{2a_\phi},
    \qquad
    b\gamma+\ip{g_s}k\ge
    \frac{(\Lambda-q)a_\phi^2+(\Lambda-M)K}{2a_\phi},
  \]
  and the Cauchy--Schwarz inequality $\ip{g_s}k\le\norm{g_s}\sqrt K$ turns
  the second inequality into a lower bound on $\norm{g_s}^2$.  We then solve the constrained minimization of
  $\Lambda_m\gamma^2+M_m\norm{g_s}^2$ under these two bounds.  Minimizing first
  in $\gamma$, and then over the three numbers $a_\phi,b,K$ that describe
  $\phi$, gives (\Cref{thm:plt23-huber-elim})
  \[
    \Lambda_m\gamma^2+M_m\norm{g_s}^2\ge \Omega_m\psi(\xi,\rho),
    \qquad
    \psi(\xi,\rho)=\min_{0\le v\le1}
    \bigl\{\rho v^2+\pos{\xi-v+v^2}\bigr\},
  \]
  where $\Omega_m$, $\xi$, and $\rho$ depend on $q$, $\Lambda$,
  and $m$.  The theorem in this region is thereby reduced to the scalar
  inequality $R_m\le \Omega_m\psi(\xi,\rho)$.  Since
  $\pos x=\max_{0\le\lambda\le1}\lambda x$, the definition of $\psi$ is a
  nested optimization,
  \[
    \psi(\xi,\rho)=\min_{0\le v\le1}\ \max_{0\le\lambda\le1}
    \bigl\{\rho v^2+\lambda(\xi-v+v^2)\bigr\}.
  \]
  We interchange the minimum and the maximum by the strong duality of this
  problem and carry out the inner minimization in $v$, which gives
  (\Cref{prop:plt23-huber-dual})
  \[
    \psi(\xi,\rho)=\max_{0\le\lambda\le1}
    \left\{\lambda\xi-\frac{\lambda^2}{4(\rho+\lambda)}\right\}.
  \]
  Two choices of $\lambda$ give the lower bounds
  (\Cref{cor:plt23-two-witnesses}),
  \begin{align*}
    \psi(\xi,\rho)&\ge\rho\xi^2\,\frac{1+4\xi}{1+2\xi}
    &&(\lambda=2\rho\xi,\text{ admissible when }2\rho\xi\le1),\\
    \psi(\xi,\rho)&\ge\xi-\frac{1}{4(1+\rho)}
    &&(\lambda=1),
  \end{align*}
  one for each of the regimes $2\rho\xi\le1$ and $2\rho\xi>1$.

  The goal $R_m\le \Omega_m\psi(\xi,\rho)$ involves only $q$, $\Lambda$, and $m$.
  From $p=1-q$ and $M^2=pq-\Lambda^2$, the two ratios
  \[
    v=\frac {p}{\Lambda}-1,
    \qquad
    \zeta=\frac{M^2}{\Lambda(\Lambda-q)}
  \]
  determine $q$, $p$, $\Lambda$, and $M$, so we may replace $q$ and $\Lambda$
  by $v$ and $\zeta$.  For every $m$, the pairs $(v,\zeta)$ to be considered
  are those satisfying
  \[
    0<v<1,
    \qquad
    \zeta(\zeta+v)\ge1
  \]
  (\Cref{subsec:plt23-chart}).  What has to be proved is that
  $R_m\le \Omega_m\psi(\xi,\rho)$ holds in that region for every odd $m\ge9$.

  This is done by the following case analysis on the region.
  \begin{itemize}[leftmargin=1.6em,topsep=2pt,itemsep=1pt]
    \item The regime $2\rho\xi\le1$ (\Cref{prop:plt23-quadratic-branch}):
    \begin{itemize}[leftmargin=1.4em,topsep=1pt,itemsep=0pt]
      \item $v\le3/5$ and $m\ge11$;
      \item $v\ge3/5$ and $m\ge11$;
      \item $m=9$.
    \end{itemize}
    \item The regime $2\rho\xi>1$ (\Cref{subsec:plt23-linear}):
    \begin{itemize}[leftmargin=1.4em,topsep=1pt,itemsep=0pt]
      \item $v\ge5/8$, for every odd $m\ge9$
      (\Cref{lem:plt23-linear-large-v});
      \item $v\le5/8$ and $\zeta\ge m-2$
      (\Cref{lem:plt23-linear-high-zeta});
      \item $\zeta\le m-2$ and $m\ge11$
      (\Cref{lem:plt23-linear-low-zeta,lem:plt23-J-growth});
      \item $m=9$ with $\zeta\le7$ and $v\le5/8$
      (\Cref{lem:plt23-N7-small-v,lem:plt23-N7-middle-v}).
    \end{itemize}
  \end{itemize}
  These can be proved by elementary analysis, except the last case, $m=9$, where we should show two polynomials are nonnegative on an interval.
  We expand the polynomials in the Bernstein basis of that interval, where each element of the basis is nonnegative there.
  To verify the nonnegativity, one can check that the coefficients of the expansion are nonnegative, which is listed in \Cref{app:bernstein}.
\end{enumerate}

We prove \Cref{thm:main} in \Cref{sec:assembly}, combining the parts above.
Separately from the proof, \Cref{sec:transferable} discusses the possible
applicability of the methods used here to other problems in extremal
combinatorics.

\section{The elementary cases \texorpdfstring{$m=3,5,7$}{m=3,5,7}}\label{sec:short-cycles}

Throughout this section, $W$ is any graphon, $p=t(K_2,W)$,
\[
  U=1-W,\qquad q=t(K_2,U)=1-p,
\]
and $T=T_U$ is the integral operator of $U$.  When $q>1/2$, the
right side of \eqref{eq:main-goal} is nonpositive, so the result is trivial.
We may therefore assume
\begin{equation}\label{eq:q-short}
  0\le q\le\frac{1}{2}.
\end{equation}
The arguments below actually prove the three cases at every edge density.

\subsection{The triangle: Goodman's one-line argument}\label{subsec:C3}

The following is the graphon form of Goodman's classical counting
argument~\cite{Goodman1959}.

\begin{proposition}\label{prop:C3}
  For every graphon $W$ with edge density $p=t(K_2,W)$,
  \[
    t(C_3,W)\ge p^3-p(1-p)^2.
  \]
\end{proposition}

\begin{proof}
  The target polynomial is
  \[
    p^3-p(1-p)^2=2p^2-p.
  \]
  Let $d_W(x)=\int_0^1W(x,y)\dd y$.  Since
  $ab\ge a+b-1$ for $a,b\in[0,1]$,
  \[
    \int_0^1W(x,z)W(y,z)\dd z\ge d_W(x)+d_W(y)-1.
  \]
  Multiplying by $W(x,y)\ge0$, integrating in $(x,y)$, and applying
  Jensen's inequality gives
  \begin{align*}
    t(C_3,W)
    &\ge \int W(x,y)\bigl(d_W(x)+d_W(y)-1\bigr)\dd x\dd y\\
    &=2\int_0^1d_W(x)^2\dd x-p
    \ge2p^2-p,
  \end{align*}
  as desired.
\end{proof}

\subsection{Path moments in the complement}\label{subsec:path-moments}

Write
\begin{equation}\label{eq:T1-decomp}
  T\one=q\one+g,
  \qquad g\perp\one.
\end{equation}
This decomposition has the explicit pointwise form
\[
  g(x)=\int_0^1U(x,y)\dd y-q
  =\bigl(1-d_W(x)\bigr)-(1-p)
  =p-d_W(x).
\]
Thus $g$ is the negative of the deviation of the $W$-degree from its
mean $p$.
For $j\ge1$, let $P_j$ denote the path with $j$ edges and put
\[
  x_j=t(P_j,U)=\ip{\one}{T^j\one}.
\]
Thus $x_1=q$.  The following first moments will be used for $C_5$.
Note that $T$ is a symmetric operator, so $\ip{f}{Tg}=\ip{Tf}{g}$ for all $f,g$.

\begin{lemma}\label{lem:path-to-4}
  With \eqref{eq:T1-decomp},
  \begin{align*}
    x_2&=q^2+\norm g^2,\\
    x_3&=q^3+2q\norm g^2+\ip g{Tg},\\
    x_4&=q^4+3q^2\norm g^2+2q\ip g{Tg}+\norm{Tg}^2.
  \end{align*}
\end{lemma}

\begin{proof}
  The first identity is 
  \[
    x_2 = \langle T\one, T\one\rangle = \langle q\one + g, q\one + g\rangle = q^2 \norm{\one}^2 + \norm{g}^2 = q^2 + \norm{g}^2.
  \]
  Next,
  $T^2\one=T(q \one + g) = q^2\one + qg + Tg$, while
  \[
    \ip{\one}{Tg}=\ip{T\one}{g}= \ip{q\one + g}{g}=\norm g^2.
  \]
  Expanding $x_3$ and $x_4$ similarly gives the other two identities:
  \begin{align*}
    x_3 &= \ip{T^2\one}{T\one} = \ip{q^2\one + qg + Tg}{q\one + g} \\
    {} &= \ip{q^2\one}{q\one} + \ip{q^2\one}{g} + \ip{qg}{q\one} + \ip{qg}{g} + \ip{Tg}{q\one} + \ip{Tg}{g} \\
    {} &= q^3 + q \norm{g}^2 + q \ip{g}{T \one} + \ip{Tg}{g} \\
    {} &= q^3 + q \norm{g}^2 + q \ip{g}{q \one + g} + \ip{Tg}{g} \\
    {} &= q^3 + 2q \norm{g}^2 + \ip{g}{Tg}, 
    \\ 
    x_4 &= \ip{T^2 \one}{T^2 \one} = \ip{q^2\one + qg + Tg}{q^2\one + qg + Tg} \\
    {} &= \ip{q^2\one}{q^2\one} + 2 \ip{q^2\one}{qg} + 2 \ip{q^2\one}{Tg} + \ip{qg}{qg} + 2 \ip{qg}{Tg} + \ip{Tg}{Tg} \\
    {} &= q^4 + 2q^2 \ip{T \one}{g} + q^2 \norm{g}^2 + 2q \ip{g}{Tg} + \norm{Tg}^2 \\
    {} &= q^4 + 2q^2 \ip{q \one + g}{g} + q^2 \norm{g}^2 + 2q \ip{g}{Tg} + \norm{Tg}^2 \\
    {} &= q^4 + 3q^2 \norm{g}^2 + 2q \ip{g}{Tg} + \norm{Tg}^2.
  \end{align*}
\end{proof}

For $C_7$, let $H_0=\one^\perp$ and
\[
  A=P_{H_0}TP_{H_0}.
\]
Concretely, $H_0$ is the space of mean-zero functions, and its orthogonal
projection is the centering operation
\[
  (P_{H_0}f)(x)=f(x)-\int_0^1f(y)\dd y.
\]
Thus $P_{H_0}$ removes the constant part of $f$, and $A$ is the restriction
of $T$ to $H_0$, in the sense that $\ip f{Ah}=\ip f{Th}$ for all
$f,h\in H_0$.  Define
\begin{equation}\label{eq:sj-def}
  s_j=\ip g{A^jg}\qquad(j\ge0).
\end{equation}
The kernel of $T$ is square-integrable, so $A$ is a compact self-adjoint
operator and $H_0$ has an orthonormal basis of its eigenfunctions.  Write
$\lambda_i$ for the distinct eigenvalues of $A$ and $P_i$ for the orthogonal
projection onto the eigenspace of $\lambda_i$, so that $\sum_iP_i$ is the
identity on $H_0$ and $A^j=\sum_i\lambda_i^jP_i$.  Then
\[
  s_j=\sum_i\lambda_i^j\norm{P_ig}^2,
\]
which is the $j$-th moment of the finite positive measure
\begin{equation}\label{eq:mu-short}
  \mu=\sum_i\norm{P_ig}^2\,\delta_{\lambda_i}.
\end{equation}
Rewriting $s_j$ in terms of this measure,
\begin{equation}\label{eq:sj-moment}
  s_j=\int\lambda^j\dd\mu(\lambda),
  \qquad s_0=\norm g^2=\mu(\R).
\end{equation}

We now express the path densities $x_j$ through $q$ and the moments $s_j$.
Split $T^n\one$ into its constant and mean-zero parts.  The constant part is
$\ip\one{T^n\one}\one=x_n\one$, so
\[
  T^n\one=x_n\one+h_n,\qquad h_n\perp\one,
\]
where $x_0=1$ and $h_0=0$.  Now \eqref{eq:T1-decomp} and the definition of
$A$ give the recursive formulas
\begin{equation}\label{eq:path-recursion}
  x_{n+1}=qx_n+\ip g{h_n},
  \qquad h_{n+1}=x_ng+Ah_n.
\end{equation}
For clarity, iterating the second recursion first gives
\begin{align*}
  h_1={}&g,\\
  h_2={}&qg+Ag,\\
  h_3={}&x_2g+qAg+A^2g,\\
  h_4={}&x_3g+x_2Ag+qA^2g+A^3g,\\
  h_5={}&x_4g+x_3Ag+x_2A^2g+qA^3g+A^4g,\\
  h_6={}&x_5g+x_4Ag+x_3A^2g+x_2A^3g
  +qA^4g+A^5g.
\end{align*}
More generally,
\[
  h_n=\sum_{j=0}^{n-1}x_{n-1-j}A^jg.
\]
If we take the inner product with $g$ and use the definition
\eqref{eq:sj-def} of $s_j$,
\[
  \ip g{h_n}
  =\sum_{j=0}^{n-1}x_{n-1-j}\ip g{A^jg}
  =\sum_{j=0}^{n-1}x_{n-1-j}s_j.
\]
Substituting this into the first recursion of \eqref{eq:path-recursion}
eliminates $h_n$, and leaves a recursion for the path densities alone,
\begin{equation}\label{eq:path-scalar-recursion}
  x_{n+1}=qx_n+\sum_{j=0}^{n-1}x_{n-1-j}s_j
  \qquad(n\ge1).
\end{equation}
We expand the next few terms that are used in the $C_5$ and $C_7$ calculations:
\begin{align}
  x_2={}& q^2+s_0, \notag\\
  x_3={}& q^3+2qs_0+s_1, \notag\\
  x_4={}& q^4+3q^2s_0+2qs_1+s_0^2+s_2, \notag\\
  x_5={}& q^5+4q^3s_0+3q^2s_1+3qs_0^2+2qs_2+2s_0s_1+s_3, \notag\\
  x_6={}& q^6+5q^4s_0+4q^3s_1+6q^2s_0^2+3q^2s_2
  +6qs_0s_1+2qs_3 \notag\\
  &\qquad{}+s_0^3+2s_0s_2+s_1^2+s_4. \label{eq:path-2-6}
\end{align}

\subsection{The pentagon: one completion of a square}\label{subsec:C5}

\begin{proposition}\label{prop:C5}
  For every graphon $W$ with edge density $p=t(K_2,W)$,
  \[
    t(C_5,W)\ge p^5-p(1-p)^4.
  \]
\end{proposition}

\begin{proof}
  Assume \eqref{eq:q-short}.  Label the integration variables cyclically by
  $z_1,\ldots,z_5$, put $z_6=z_1$, and abbreviate
  \[
    U_i=U(z_i,z_{i+1})\qquad(1\le i\le5).
  \]
  Since $W=1-U$, the integrand defining $t(C_5,W)$ is
  $\prod_{i=1}^5(1-U_i)$.  Ordinary inclusion--exclusion gives the exact
  pointwise expansion
  \begin{align*}
    \prod_{i=1}^5(1-U_i)
    ={}&1-\sum_{i=1}^5U_i
    +\sum_{1\le i<j\le5}U_iU_j
    -\sum_{1\le i<j<k\le5}U_iU_jU_k\\
    &+\sum_{i=1}^5\prod_{j\ne i}U_j
    -\prod_{i=1}^5U_i.
  \end{align*}
  We now integrate this identity over $(z_1,\ldots,z_5)$.  Each single-edge
  term has integral $q=x_1$, so the first sum contributes $5q$.  Among the
  $\binom52=10$ pairs of edges, five pairs are adjacent around the cycle;
  each forms a two-edge path and hence has integral $x_2$.  The other five
  pairs are disjoint edges.  Their variables separate, so each has integral
  $x_1^2=q^2$.  Therefore
  \[
    \int\sum_{i<j}U_iU_j=5x_2+5q^2.
  \]

  Similarly, among the $\binom53=10$ triples, five consist of three
  consecutive edges and form a path with three edges, contributing $x_3$
  each.  Each of the other five consists of a two-edge path together with one
  disjoint edge.  Homomorphism density is multiplicative over disjoint unions,
  so each of these contributes $x_2x_1=qx_2$.  Thus
  \[
    \int\sum_{i<j<k}U_iU_jU_k=5x_3+5qx_2.
  \]
  There are five ways to select four edges; deleting any one edge from the
  cycle leaves a four-edge path, so every corresponding integral is $x_4$.
  Finally, selecting all five edges gives $c_5:=t(C_5,U)$.  Combining these
  contributions with their alternating signs yields
  \begin{align*}
    t(C_5,1-U)
    ={}&1-5q+5x_2+5q^2-5x_3-5qx_2+5x_4-c_5.
  \end{align*}
  Since $0\le U\le1$, deleting one edge from the cycle integrand gives
  $c_5\le x_4$.  Therefore
  \begin{equation}\label{eq:C5-IE}
    t(C_5,W)\ge1-5q+5q^2+5(1-q)x_2-5x_3+4x_4.
  \end{equation}
  Subtracting the lower bound
  \[
    (1-q)^5-(1-q)q^4=1-5q+10q^2-10q^3+4q^4
  \]
  from \eqref{eq:C5-IE} and using \Cref{lem:path-to-4} gives
  \begin{align}
    \Phi_5={}&4\norm{Tg}^2+(8q-5)\ip g{Tg}
    +(12q^2-15q+5)\norm g^2 \notag\\
    ={}&4\left\|Tg+\frac{8q-5}{8}g\right\|^2
    +\left(8q^2-10q+\frac{55}{16}\right)\norm g^2. \label{eq:Phi5}
  \end{align}
  The quadratic coefficient $8q^2-10q+\frac{55}{16}$ has discriminant
  \[
    100-4\cdot8\cdot\frac{55}{16}=-10<0
  \]
  and positive leading coefficient, so it is positive for every $q$.  Thus
  $\Phi_5\ge0$, proving the claim.
\end{proof}

\subsection{The heptagon: a univariate sum of squares}\label{subsec:C7}

\begin{proposition}\label{prop:C7}
  For every graphon $W$ with edge density $p=t(K_2,W)$,
  \[
    t(C_7,W)\ge p^7-p(1-p)^6.
  \]
\end{proposition}

\begin{proof}
  Write $U=1-W$ and $q=1-p=t(K_2,U)$, and let $x_j=t(P_j,U)$ be the path
  densities of \eqref{eq:sj-def} above, so that $x_1=q$.  If $p<1/2$, the
  right side of the claim is nonpositive and there is nothing to prove, so
  assume $p\ge1/2$, that is $q\le1/2$.  As in the proof for $C_5$, label the
  integration variables cyclically by $z_1,\ldots,z_7$, put $z_8=z_1$, and
  write
  \[
    U_i=U(z_i,z_{i+1})\qquad(1\le i\le7).
  \]
  Expanding the seven factors $1-U_i$ gives the exact inclusion--exclusion
  identity
  \begin{equation}\label{eq:C7-subset-expansion}
    t(C_7,W)
    =\sum_{S\subseteq\{1,\ldots,7\}}(-1)^{|S|}
    \int_{[0,1]^7}\prod_{i\in S}U_i\prod_{j=1}^7\dd z_j.
  \end{equation}
  A proper subset $S$ selects a disjoint union of paths around the cycle, and
  its variables separate between the components, so its term is the
  corresponding product of path densities; the full subset gives
  $c_7=t(C_7,U)$.  Collecting the subsets of each size, \eqref{eq:C7-subset-expansion}
  reads
  \begin{align*}
    t(C_7,W)={}&1\\
    &-7x_1\\
    &+7x_2+14x_1^2\\
    &-7x_3-21x_1x_2-7x_1^3\\
    &+7x_4+14x_1x_3+7x_2^2+7x_1^2x_2\\
    &-7x_5-7x_1x_4-7x_2x_3\\
    &+7x_6\\
    &-c_7.
  \end{align*}
  The seven subsets of size one each contribute $x_1$; of the $21$ subsets of
  size two, seven consist of two adjacent edges and contribute $x_2$ while the
  other fourteen contribute $x_1^2$; and so on, the coefficients in each line
  summing to $\binom7k$.  Because $0\le U_i\le1$, deleting one factor from the
  integrand of $c_7$
  can only increase it; the remaining six factors form a six-edge path.
  Hence
  \[
    c_7=\int\prod_{i=1}^7U_i
    \le\int\prod_{i=1}^6U_i=x_6.
  \]
  Replacing $-c_7$ by $-x_6$ in the preceding identity and collecting the
  coefficients of the same path densities gives
  \begin{align}
    t(C_7,W)\ge{}&6x_6-7x_5+7(1-q)x_4+7(2q-1)x_3 \notag\\
    &+7(q^2-3q+1)x_2+7x_2^2-7x_2x_3 \notag\\
    &+1-7q+14q^2-7q^3. \label{eq:C7-IE}
  \end{align}
  Since
  \[
    (1-q)^7-(1-q)q^6
    =1-7q+21q^2-35q^3+35q^4-21q^5+6q^6,
  \]
  substitution of \eqref{eq:path-2-6} reduces the desired inequality to the nonnegativity of the polynomial
  \begin{align}
    \Phi_7={}&6s_4+(12q-7)s_3+(18q^2-21q+7)s_2 \notag\\
    &+(24q^3-42q^2+28q-7)s_1 \notag\\
    &+(30q^4-70q^3+70q^2-35q+7)s_0 \notag\\
    &+12s_0s_2+(36q-21)s_0s_1
    +(36q^2-42q+14)s_0^2+6s_0^3+6s_1^2. \label{eq:Phi7-expanded}
  \end{align}
  Write $r$ for the total mass $s_0=\mu(\R)$.  Keeping one factor $r$ in each
  product and expressing every remaining $s_j$ by \eqref{eq:sj-moment} as an
  integral against $\mu$, the polynomial \eqref{eq:Phi7-expanded} becomes
  \begin{equation}\label{eq:Phi7-integral}
    \Phi_7=6s_1^2+\int R_{q,r}(\lambda)\dd\mu(\lambda),
  \end{equation}
  where
  \[
    R_{q,r}(\lambda)=P_q(\lambda)+rB_q(\lambda)+6r^2,
  \]
  \begin{align*}
    P_q(\lambda)={}&6\lambda^4+(12q-7)\lambda^3
    +(18q^2-21q+7)\lambda^2\\
    &+(24q^3-42q^2+28q-7)\lambda
    +30q^4-70q^3+70q^2-35q+7,\\
    B_q(\lambda)={}&12\lambda^2+(36q-21)\lambda+36q^2-42q+14.
  \end{align*}
  We show that both $P_q$ and $B_q$ are nonnegative.

  The minimum of $B_q$ as a quadratic in $\lambda$ occurs at
  $\lambda_*=(7-12q)/8$, and
  \[
    B_q(\lambda_*)=\frac{144q^2-168q+77}{16}\ge\frac{29}{16},
  \]
  because
  \[
    144q^2-168q+48=144\left(q-\frac{1}{2}\right)
    \left(q-\frac{2}{3}\right)\ge0
  \]
  on $[0,1/2]$.

  For $P_q$, set
  \begin{align*}
    D(q)&=288q^2-336q+119,\\
    C(q)&=24q^3-42q^2+28q-7,\\
    N(q)&=5184q^6-18144q^5+28602q^4-25802q^3\\
    &\qquad+13874q^2-4165q+539.
  \end{align*}
  Direct expansion gives the exact identity
  \begin{align}
    P_q(\lambda)
    ={}&6\left(\lambda^2+\left(q-\frac{7}{12}\right)\lambda\right)^2 \notag\\
    &+\frac{D(q)}{24}
    \left(\lambda+\frac{12C(q)}{D(q)}\right)^2
    +\frac{N(q)}{D(q)}. \label{eq:Pq-SOS}
  \end{align}
  Writing $y=1/2-q\ge0$,
  \[
    D(q)=288y^2+48y+23>0
  \]
  and
  \begin{align*}
    8N(q)={}&41472y^6+20736y^5+21456y^4+7984y^3\\
    &+2032y^2+316y+11>0.
  \end{align*}
  Thus $P_q(\lambda)\ge0$ for every real $\lambda$.  Since
  $r\ge0$, $R_{q,r}(\lambda)\ge0$, and \eqref{eq:Phi7-integral} yields
  $\Phi_7\ge0$.
\end{proof}

The identity \eqref{eq:Pq-SOS} is an explicit sum-of-squares certificate for
$P_q$.  For background on positivity of one-variable polynomials and on such
certificates, see Powers--Reznick~\cite{PowersReznick2000}.

\iffalse
\begin{remark}[Pointwise nonnegativity versus operator positivity]
  The $C_5$ and $C_7$ arguments do not require the kernel operator $T_U$
  to be positive semidefinite.  Such positivity would mean that
  \[
    \ip f{T_Uf}
    =\int_{[0,1]^2}U(x,y)f(x)f(y)\dd x\dd y\ge0
    \qquad\text{for every }f\in L^2([0,1]).
  \]
  Although $U(x,y)\ge0$ pointwise, this condition need not hold, because
  $f(x)f(y)$ may be negative.  For example, let $U$ be the balanced
  complete bipartite graphon, equal to $1$ between the two halves of
  $[0,1]$ and to $0$ within each half.  If $f$ equals $1$ on the first
  half and $-1$ on the second, then $T_Uf=-\tfrac{1}{2}f$, and hence
  $\ip f{T_Uf}=-\tfrac{1}{2}\norm f^2<0$.

  What the proofs use instead is self-adjointness, which follows from the
  symmetry of $U$.  The spectral theorem then associates with $g$ a
  nonnegative measure $\mu$ satisfying $s_j=\int\lambda^j\dd\mu(\lambda)$.
  Here nonnegative refers to the weights of the measure; its support may still
  contain negative eigenvalues $\lambda$.  This is why the proof above shows
  that the relevant polynomials are nonnegative for every real $\lambda$,
  rather than only for $\lambda\ge0$.  No positive-semidefiniteness assumption
  on $T_U$ is needed.
\end{remark}

\fi

\section{Reductions and notation used in both regions}\label{sec:common-operator}

\subsection{Step-graphon reduction}\label{subsec:step-reduction}

As described in the proof map above, the cycle densities appear as the
coefficients of $-\log\det(I-zT_W)$.  To avoid taking determinants of
operators on an infinite-dimensional space, we reduce the theorem to step
graphons, whose operators have finite rank.  The approximation is the
dyadic conditional expectation of \cite[Appendix~A.3.6]{Lovasz2012}, which
has the property we need here: the approximating graphons have edge density
exactly $p$.

\begin{lemma}[Finite-rank reduction]\label{lem:step-reduction}
  For each fixed odd $m\ge3$, it is enough to prove \eqref{eq:main-goal} for
  symmetric step graphons.
\end{lemma}

\begin{proof}
  For $n\ge1$, divide $[0,1]$ into the $2^n$ dyadic
  intervals
  \[
    I_{n,a}=\left[\frac{a}{2^n},\frac{a+1}{2^n}\right),
    \qquad 0\le a<2^n,
  \]
  with the convention that the last interval contains the endpoint
  $1$.  Let $\mathcal A_n$ be the finite sigma-algebra consisting of all
  unions of these intervals.  Thus the atoms of the product sigma-algebra
  $\mathcal A_n\otimes\mathcal A_n$ are precisely the dyadic rectangles
  $I_{n,a}\times I_{n,b}$.  The partitions refine as $n$ increases, and
  dyadic intervals generate the Borel sigma-algebra.

  On each such rectangle, replace $W$ by its average.  More explicitly, set
  \begin{equation}\label{eq:dyadic-step-approximation}
    w^{(n)}_{ab}
    =2^{2n}\int_{I_{n,a}\times I_{n,b}}W(x,y)\dd x\dd y,
    \qquad
    W_n(x,y)=w^{(n)}_{ab}
    \quad\text{if }(x,y)\in I_{n,a}\times I_{n,b}.
  \end{equation}
  This cell-averaging operation is the standard conditional 
  expectation operation with respect to the sigma algebra
  $\mathcal A_n\otimes\mathcal A_n$, and is 
  denoted by the notation
  $\E[\,-\, \mid\mathcal A_n\otimes\mathcal A_n]$. So,
  $W_n = \E[W\mid\mathcal A_n\otimes\mathcal A_n] $. Since $0\le W\le1$,
  every cell average lies in $[0,1]$.  Moreover, symmetry of $W$ gives
  $w^{(n)}_{ab}=w^{(n)}_{ba}$.  Hence $W_n$ is a symmetric step graphon.

  There is also a completely finite interpretation.  The matrix
  $(w^{(n)}_{ab})_{0\le a,b<2^n}$ is the weighted adjacency matrix of a graph
  on $2^n$ equally weighted vertices (loops are allowed), and
  \[
    t(C_m,W_n)
    =\frac{1}{2^{nm}}
    \sum_{a_1,\ldots,a_m=0}^{2^n-1}
    w^{(n)}_{a_1a_2}w^{(n)}_{a_2a_3}\cdots
    w^{(n)}_{a_ma_1}.
  \]
  The same finite structure can be seen directly at the operator level.  For
  any function $f$, if $x\in I_{n,a}$, then
  \[
    (T_{W_n}f)(x)
    =\int_0^1W_n(x,y)f(y)\dd y
    =\sum_{b=0}^{2^n-1}w^{(n)}_{ab}
    \int_{I_{n,b}}f(y)\dd y.
  \]
  Thus the output $T_{W_n}f$ is constant on each interval $I_{n,a}$, no
  matter how complicated the input $f$ is.  Every such output is therefore a
  linear combination of the $2^n$ indicator functions
  $\one_{I_{n,0}},\ldots,\one_{I_{n,2^n-1}}$.  The set of all possible
  outputs is called the \emph{range} of the operator, and its dimension is the
  \emph{rank}.  Since the range is contained in a space of dimension $2^n$,
  the rank is at most $2^n$.  This explains the term ``finite-rank
  reduction.''

  Two facts remain to be checked: the edge density is unchanged, and the cycle
  density converges.  For the first, summing the averages in
  \eqref{eq:dyadic-step-approximation} over all cells gives
  \[
    \int_{[0,1]^2}W_n(x,y)\dd x\dd y
    =\sum_{a,b}|I_{n,a}|\,|I_{n,b}|\,w^{(n)}_{ab}
    =\int_{[0,1]^2}W(x,y)\dd x\dd y=p.
  \]
  Thus $t(K_2,W_n)=p$ for every $n$, not merely in the limit.

  Next we verify that $W_n\to W$ in $L^2([0,1]^2)$.  Dyadic rectangle
  step functions are dense in this space.  Therefore, given $\varepsilon>0$,
  we may choose an integer $N$ and a function $S$, constant on every
  rectangle $I_{N,a}\times I_{N,b}$, such that
  $\norm{W-S}_2<\varepsilon$.  For $n\ge N$, the function $S$ is also
  constant on every level-$n$ rectangle.  Averaging is an $L^2$-contraction
  (this is Jensen's inequality on each rectangle), so
  \[
    \norm{W_n-S}_2
    =\norm{\E[W-S\mid\mathcal A_n\otimes\mathcal A_n]}_2
    \le \norm{W-S}_2.
  \]
  Consequently
  \[
    \norm{W_n-W}_2
    \le \norm{W_n-S}_2+\norm{S-W}_2
    <2\varepsilon,
  \]
  which proves the asserted convergence.

  Finally, cycle densities are continuous under this convergence.  Put $x_{m+1}=x_1$, and for a fixed
  $(x_1,\ldots,x_m)$, write
  \[
    a_i=W_n(x_i,x_{i+1}),\qquad b_i=W(x_i,x_{i+1}).
  \]
  The telescoping identity
  \[
    \prod_{i=1}^m a_i-\prod_{i=1}^m b_i
    =\sum_{j=1}^m
    \left(\prod_{i<j}a_i\right)(a_j-b_j)
    \left(\prod_{i>j}b_i\right)
  \]
  and the bounds $0\le a_i,b_i\le1$ show, after integration, that
  \begin{align*}
    \bigl|t(C_m,W_n)-t(C_m,W)\bigr|
    &\le \sum_{j=1}^m
    \int_{[0,1]^m}
    |W_n(x_j,x_{j+1})-W(x_j,x_{j+1})|
    \prod_{i=1}^m\dd x_i \\
    &=m\norm{W_n-W}_1
    \le m\norm{W_n-W}_2
    \longrightarrow0.
  \end{align*}
  Here the equality in the second line follows because all variables other
  than $x_j,x_{j+1}$ integrate out to $1$.

  Now suppose \eqref{eq:main-goal} has been proved for every symmetric step
  graphon.  Applying it to $W_n$, whose edge density is exactly $p$, gives
  \[
    t(C_m,W_n)\ge p^m-p(1-p)^{m-1}.
  \]
  Letting $n\to\infty$ proves the same inequality for $W$, as required.
\end{proof}

\subsection{Reading the power-series
calculations}\label{subsec:analytic-reading}

Several later arguments introduce an auxiliary variable $z$ and use the
notation $[z^m]F(z)$; the exponent of $z$ is the length of a walk or cycle.
We do not assign a value to $z$; every $F$ below is a formal power series,
that is, a sequence of real coefficients $(c_j)_{j\ge0}$ written as
$F=\sum_{j\ge0}c_jz^j$, and $[z^m]F=c_m$.  We say that two such series are
equal when all their coefficients are equal.

Four operations on these series are used below, and in each of them every
coefficient of the result is computed from finitely many coefficients of the
arguments.  Series are added coefficientwise.  They are multiplied by the
Cauchy product,
\[
  [z^m](FG)=\sum_{i+j=m}\bigl([z^i]F\bigr)\bigl([z^j]G\bigr).
\]
A series with $c_0\ne0$ has an inverse, obtained by solving the equations
for the coefficients of $FF^{-1}=1$ recursively, so we may divide by such a
series.  Finally, if we write $\ord F$ for the order of $F$, that is the
least $j$ with $c_j\ne0$, then the series
\[
  -\log(1-Y)=\sum_{k\ge1}\frac{Y^k}{k}
\]
is defined whenever $\ord Y\ge1$: in that case $\ord Y^k\ge k$, so only the
terms with $k\le m$ contribute to $[z^m]$.

By \Cref{lem:step-reduction} we will work with a step graphon, whose operator
is a finite matrix; this is made explicit in \Cref{sec:common-operator}.  We
build our series from such matrices, using the two expansions in the next
lemma.

\begin{lemma}[Series expansions used below]\label{lem:series-expansions}
  Let $X$ be a real symmetric $N\times N$ matrix.  As formal power series
  in $z$:
  \begin{enumerate}[label=\textup{(\arabic*)},leftmargin=2.6em]
    \item \emph{(Resolvent.)}  The matrix $I-zX$ is invertible over $\R[[z]]$
    and
    \[
      (I-zX)^{-1}=\sum_{j\ge0}z^jX^j,
    \]
    so that $[z^j]\ip u{(I-zX)^{-1}v}=\ip u{X^jv}$ for fixed vectors $u,v$.
    \item \emph{(Log-determinant.)}  The polynomial $\det(I-zX)$ has constant
    term $1$, and
    \[
      -\log\det(I-zX)=\sum_{j\ge1}\frac{\Tr(X^j)}{j}z^j,
    \]
    so that $m\,[z^m]\bigl(-\log\det(I-zX)\bigr)=\Tr(X^m)$.
  \end{enumerate}
\end{lemma}

\begin{proof}
  (1)  The constant term of $I-zX$ is $I$, so the matrix is invertible over
  $\R[[z]]$; multiplying $\sum_{j\ge0}z^jX^j$ by $I-zX$ cancels all positive
  powers of $z$ and leaves $I$.

  (2)  Both sides are unchanged under $X\mapsto Q^{\!\top}XQ$ with $Q$
  orthogonal, so we may take $X$ diagonal with entries $\lambda_i$.  Then
  $\det(I-zX)=\prod_i(1-z\lambda_i)$, and $-\log$ of a product is the sum of
  the $-\log$s, whence
  \[
    -\log\det(I-zX)
    =\sum_{i=1}^N\sum_{j\ge1}\frac{\lambda_i^j}{j}z^j
    =\sum_{j\ge1}\frac{\Tr(X^j)}{j}z^j,
  \]
  the interchange being a finite sum in each coefficient.
\end{proof}

For an example of this formal series, take $X=T_W$ and $m=3$.  Item~(2)
gives
\[
  -\log\det(I-zT_W)
  =\Tr(T_W)\,z+\frac{\Tr(T_W^2)}{2}z^2+\frac{\Tr(T_W^3)}{3}z^3+\cdots,
\]
whence
\[
  3\,[z^3]\bigl(-\log\det(I-zT_W)\bigr)=\Tr(T_W^3)=t(C_3,W)
\]
by \eqref{eq:finite-trace-cycle}.

Alternatively, one may read these series as functions.  If a series
converges in a neighbourhood of the origin, its sum is a function there, and
the coefficients are recovered from that function by differentiation,
\begin{equation}\label{eq:coeff-def}
  [z^m]F=\frac{F^{(m)}(0)}{m!}.
\end{equation}
The arguments below are the same under either reading.

\subsection{The block decomposition and the support of the spectral measure}
\label{subsec:block-decomposition}

From here to the end of \Cref{sec:plt23}, we assume $W$ to be a
step graphon, which \Cref{lem:step-reduction} permits, and all determinants
and traces are taken in the resulting finite-dimensional setting.  Fix a finite
measurable partition
\[
  [0,1]=B_1\sqcup\cdots\sqcup B_N
\]
into sets of positive measure such that $W$, and hence $U=1-W$, is
constant on each rectangle $B_a\times B_b$.

Let
\[
  \mathcal E=\operatorname{span}
  \{\one_{B_1},\ldots,\one_{B_N}\}\subseteq L^2([0,1]).
\]
This is the $N$-dimensional space of functions that are constant on every
part $B_a$.  If a step kernel $K$ has value $K_{ab}$ on
$B_a\times B_b$, then, for $x\in B_a$,
\begin{equation}\label{eq:step-operator-action}
  (T_Kf)(x)=\sum_{b=1}^N K_{ab}\int_{B_b}f(y)\dd y.
\end{equation}
Consequently, $T_Kf\in\mathcal E$ for every $f\in L^2$.  Moreover, if
$f\in\mathcal E^\perp$, then every integral in
\eqref{eq:step-operator-action} is zero, so $T_Kf=0$.  Hence $T_K$ is
determined by its restriction to the finite-dimensional space
$\mathcal E$.

The cycle densities of $K$ are computed from this finite matrix as follows.
In the orthonormal basis $e_a=\one_{B_a}/\sqrt{|B_a|}$, the matrix of
$T_K|_{\mathcal E}$ has entries
\[
  \ip{e_a}{T_Ke_b}=K_{ab}\sqrt{|B_a||B_b|}.
\]
Here $|B_a|$ denotes the measure of $B_a$.
Multiplying this matrix around a cycle shows, for every $r\ge2$, that
\begin{align}
  \Tr(T_K^r)
  &=\sum_{a_1,\ldots,a_r=1}^N
  |B_{a_1}|\cdots|B_{a_r}|
  K_{a_1a_2}\cdots K_{a_ra_1} \notag\\
  &=t(C_r,K). \label{eq:finite-trace-cycle}
\end{align}

\tk{TODO (structure): move the shared identities \eqref{eq:logdet-cycle} and
\eqref{eq:block-determinant-identity}, currently stated inside
\Cref{subsec:pge23-two-sided-schur}, into this section, next to the
formal-series definitions, so that \Cref{sec:plt23} no longer depends on
\Cref{sec:pge23}.  Their labels are deliberately kept unprefixed for this
reason.}

We now split $\mathcal E$ into the constant direction and its orthogonal
complement, and decompose $T_U$ accordingly.  The constant function $\one$
lies in $\mathcal E$ and has norm one, because the underlying probability
space has total mass one, so setting
\[
  \mathcal E_0=\mathcal E\cap\one^\perp
\]
gives the orthogonal decomposition $\mathcal E=\R\one\oplus\mathcal E_0$.
Put
\[
  g=T_U\one-q\one\in\mathcal E_0,
  \qquad
  A=P_{\mathcal E_0}T_UP_{\mathcal E_0},
\]
so that $A$ is a self-adjoint map on $\mathcal E_0$, equivalently a real
symmetric matrix in an orthonormal basis of that space.  These two determine
$T_U$: for $f=a\one+h$ with $h\in\mathcal E_0$,
\[
  T_Uf=\bigl(aq+\ip gh\bigr)\one+\bigl(ag+Ah\bigr),
\]
since $\ip\one{T_Uh}=\ip{T_U\one}h=\ip gh$ and $P_{\mathcal E_0}T_Uh=Ah$.
In block form,
\begin{equation}\label{eq:block-U}
  T_U=
  \begin{pmatrix}
    q&g^*\\
    g&A
  \end{pmatrix},
\end{equation}
where $g^*$ denotes the linear functional $h\mapsto\ip gh$.  This $A$ plays
the role of the operator in \Cref{subsec:path-moments}: for a step graphon,
the compression $P_{\one^\perp}T_UP_{\one^\perp}$ used there is zero on
$\mathcal E^\perp$, and its restriction to $\mathcal E_0$ is $A$.


\begin{lemma}[Operator norm bound]\label{lem:half-norm}
  The spectrum of $A$ lies in $[-1/2,1/2]$; equivalently,
  \begin{equation}\label{eq:A-half}
    \norm A_{\mathrm{op}}\le\frac{1}{2}.
  \end{equation}
\end{lemma}

\begin{proof}
  Let $f\in\mathcal E_0$ with $\norm f=1$, and write $f=f_+-f_-$ for nonnegative
  functions $f_+$ and $f_-$ on $[0,1]$.  Since
  $\int f=0$,
  \[
    \int f_+=\int f_-=:a,
    \qquad 2a=\lVert f\rVert_1\le\lVert f\rVert_2=1.
  \]
  We also have
  \begin{align*}
    \ip f{T_Uf}
    ={}&\ip {f_+}{T_Uf_+}+\ip {f_-}{T_Uf_-}
    -2\ip {f_+}{T_Uf_-}.
  \end{align*}
  Each of the three inner products on the right is nonnegative.  Moreover,
  \begin{align*}
    \ip {f_+}{T_Uf_+} &= \int_{[0,1]^2}U(x,y)f_+(x)f_+(y)\dd x\dd y \le \int_{[0,1]^2}f_+(x)f_+(y)\dd x\dd y = \left(\int f_+\right)^2,
    \\
    \ip {f_-}{T_Uf_-} &= \int_{[0,1]^2}U(x,y)f_-(x)f_-(y)\dd x\dd y \le \int_{[0,1]^2}f_-(x)f_-(y)\dd x\dd y = \left(\int f_-\right)^2,
    \\
    \ip {f_+}{T_Uf_-} &= \int_{[0,1]^2}U(x,y)f_+(x)f_-(y)\dd x\dd y \le \int_{[0,1]^2}f_+(x)f_-(y)\dd x\dd y = \left(\int f_+\right)\left(\int f_-\right).
  \end{align*}
  because $f_+$ and $f_-$ are nonnegative and so replacing $U(x,y)$ by $1$ in each of the above three cases gives the corresponding product of
  integrals.  For the upper bound, we discard the nonpositive cross term and
  bound the two diagonal terms; hence
  \[
    \ip f{T_Uf}
    \le\left(\int f_+\right)^2+\left(\int f_-\right)^2
    =2a^2\le\frac{1}{2}.
  \]
  For the lower bound, we instead discard the two nonnegative diagonal terms
  and bound the cross term, obtaining
  \[
    \ip f{T_Uf}
    \ge-2\left(\int f_+\right)\left(\int f_-\right)
    =-2a^2\ge-\frac{1}{2}.
  \]
  The upper and lower bounds just proved concern the quadratic form of $T_U$
  on $\mathcal E_0$, and this is also the quadratic form of $A$: since
  $P_{\mathcal E_0}f=f$,
  \[
    \ip f{Af}=\ip f{P_{\mathcal E_0}T_UP_{\mathcal E_0}f}=\ip f{T_Uf},
  \]
  so $\bigl|\ip f{Af}\bigr|\le\tfrac{1}{2}$ for every unit $f\in\mathcal E_0$.
  As $A$ is self-adjoint, its operator norm is the supremum of
  $\bigl|\ip f{Af}\bigr|$ over such $f$, whence
  $\norm A_{\mathrm{op}}\le\tfrac{1}{2}$.
\end{proof}

In the remaining section, $A$ and $g$ will appear through the quantities of the form $\ip g{F(A)g}$, which can be expressed as an integral against a measure supported on the spectrum of $A$.
Specifically, since $A$ is a finite real symmetric matrix, it has an orthogonal eigenspace decomposition
\[
  A=\sum_i\lambda_iP_i,
\]
where the distinct $\lambda_i$ are its eigenvalues and $P_i$ is the
orthogonal projection onto the corresponding eigenspace.  
Define the finite positive measure
\begin{equation}\label{eq:mu-def}
  \mu=\sum_i\norm{P_i g}^2\,\delta_{\lambda_i},
\end{equation}
where $\delta_{\lambda_i}$ denotes a unit point mass at $\lambda_i$; thus
$\mu$ assigns to each eigenvalue the weight measuring how much of $g$ lies
in its eigenspace, and by \Cref{lem:half-norm} it is supported on
$[-1/2,1/2]$.  For any function $F$ defined on the eigenvalues of $A$,
orthogonality of the eigenspaces gives
\[
  \ip g{F(A)g}
  =\sum_i F(\lambda_i)\norm{P_i g}^2
  =\int F(\lambda)\dd\mu(\lambda).
\]

In particular, expanding $(I-zA)^{-1}=\sum_{j\ge0}z^jA^j$ by
\Cref{lem:series-expansions}\,(1) and taking $F(\lambda)=\lambda^j$ in each
coefficient,
\begin{equation}\label{eq:resolvents-common}
  \ip g{(I-zA)^{-1}g}
  =\sum_{j\ge0}z^j\int\lambda^j\dd\mu(\lambda)
  =\int_{-1/2}^{1/2}\frac{\dd\mu(\lambda)}{1-\lambda z}.
\end{equation}

\section{Above the threshold \texorpdfstring{$p\ge2/3$}{p>=2/3}: Dirichlet averages and gamma smoothing}
\label{sec:pge23}

Throughout this section,
\begin{equation}\label{eq:pge23-q}
  q = 1 - p, \qquad 0 \le q \le \frac{1}{3},
\end{equation}
and $m\ge3$ is odd.  We prove \eqref{eq:main-goal} for all such $m$.

\subsection{A two-sided Schur-complement identity and cancellation of the odd
trace}\label{subsec:pge23-two-sided-schur}

By \Cref{lem:series-expansions}\,(2) and \eqref{eq:finite-trace-cycle}, every
step kernel $K$ satisfies
\begin{equation}\label{eq:logdet-cycle}
  -\log\det(I-zT_K)=\sum_{j\ge1}\frac{\Tr(T_K^j)}{j}z^j,
  \qquad
  m\,[z^m]\bigl(-\log\det(I-zT_K)\bigr)=t(C_m,K)
\end{equation}
for every $m\ge2$. We expand
these two series along the block decomposition \eqref{eq:block-U}. 
As the proof of the next lemma shows, each of them splits into three
summands: one coming from the mean-zero block $\mp A$, one from the scalar
entry, and the third coming from the Schur complement:
\begin{align}
  L_-(z)&=-\log\left(1-\frac{z^2\ip g{(I+zA)^{-1}g}}{1-pz}\right),
  \label{eq:pge23-Lminus}\\
  L_+(z)&=-\log\left(1-\frac{z^2\ip g{(I-zA)^{-1}g}}{1-qz}\right).
  \label{eq:pge23-Lplus}
\end{align}
Writing
\begin{align}
  R_-(z)&=\ip g{(I+zA)^{-1}g}
  =\int\frac{\dd\mu(\lambda)}{1+\lambda z},\label{eq:pge23-Rminus}\\
  R_+(z)&=\ip g{(I-zA)^{-1}g}
  =\int\frac{\dd\mu(\lambda)}{1-\lambda z},\label{eq:pge23-Rplus}
\end{align}
as in \eqref{eq:resolvents-common} with $A$ replaced by $\mp A$, the inner
series of $L_\pm$ is $z^2R_\pm(z)/(1-cz)$, with $c=p$ or $q$.  It has order
at least two, so $L_\pm$ are well defined formal power series, as discussed
in \Cref{subsec:analytic-reading}.  For odd $m$, put
\begin{equation}\label{eq:pge23-Sm}
  S_m=m[z^m]\bigl(L_-(z)+L_+(z)\bigr).
\end{equation}

\begin{lemma}[Sum of the two cycle densities]\label{lem:pge23-two-sided-identity}
  For every odd $m\ge3$,
  \begin{equation}\label{eq:pge23-two-sided-identity}
    t(C_m,W)+t(C_m,U)=p^m+q^m+S_m.
  \end{equation}
  Consequently,
  \begin{equation}\label{eq:pge23-one-sided-identity}
    t(C_m,W)-\bigl(p^m-pq^{m-1}\bigr)
    =q^{m-1}+S_m-t(C_m,U).
  \end{equation}
\end{lemma}

\begin{proof}
  We first put the two matrices $I-zT_U$ and $I-zT_W$ in block form.  The
  operator $T_1$ of the constant kernel $1$ sends $a\one+h$, with
  $h\in\mathcal E_0$, to $a\one$, because $\int h=0$; thus $T_1$ has block
  matrix $\left(\begin{smallmatrix}1&0\\0&0\end{smallmatrix}\right)$ and
  $T_W=T_1-T_U$ has block matrix
  \[
    T_W=
    \begin{pmatrix}1&0\\0&0\end{pmatrix}
    -
    \begin{pmatrix}q&g^*\\g&A\end{pmatrix}
    =
    \begin{pmatrix}p&-g^*\\-g&-A\end{pmatrix}.
  \]
  Consequently
  \[
    I-zT_U=
    \begin{pmatrix}1-qz&-zg^*\\-zg&I-zA\end{pmatrix},
    \qquad
    I-zT_W=
    \begin{pmatrix}1-pz&zg^*\\zg&I+zA\end{pmatrix}.
  \]

  The determinants decompose along this block structure by the following
  elementary identity.  If $D$ is invertible, then
  \begin{equation}\label{eq:block-determinant-identity}
    \det\begin{pmatrix}\alpha&v^*\\v&D\end{pmatrix}
    =\det(D)\bigl(\alpha-v^*D^{-1}v\bigr),
  \end{equation}
  the scalar $\alpha-v^*D^{-1}v$ being the Schur complement of $D$.  
  To see this, multiply on the left by the block matrix
  $\left(\begin{smallmatrix}1&-v^*D^{-1}\\0&I\end{smallmatrix}\right)$, which
  has determinant one and makes the product block triangular; the determinant
  of the latter is the product of its diagonal-block determinants.  We may
  take $D=I\mp zA$ here: these blocks have constant term $I$, hence are
  invertible over $\R[[z]]$ by \Cref{lem:series-expansions}\,(1).  Applying
  \eqref{eq:block-determinant-identity} with $D=I-zA$ and $v=-zg$, and then
  with $D=I+zA$ and $v=zg$, factorizes the determinants into three parts,
  \begin{align*}
    \det(I-zT_U)&=\det(I-zA)\bigl(1-qz-z^2R_+(z)\bigr) =\det(I-zA)\,(1-qz)
    \left(1-\frac{z^2R_+(z)}{1-qz}\right),\\
    \det(I-zT_W)&=\det(I+zA)\bigl(1-pz-z^2R_-(z)\bigr)=\det(I+zA)\,(1-pz)
    \left(1-\frac{z^2R_-(z)}{1-pz}\right).
  \end{align*}
  Taking $-\log$ turns each product into a sum of three series,
  \begin{align*}
    -\log\det(I-zT_U)&=-\log\det(I-zA)-\log(1-qz)+L_+(z),\\
    -\log\det(I-zT_W)&=-\log\det(I+zA)-\log(1-pz)+L_-(z),
  \end{align*}
  and we extract the coefficient of $z^m$ from each summand.  On the left,
  \eqref{eq:logdet-cycle} gives
  \[
    m\,[z^m]\bigl(-\log\det(I-zT_U)\bigr)=t(C_m,U),
    \qquad
    m\,[z^m]\bigl(-\log\det(I-zT_W)\bigr)=t(C_m,W).
  \]
  For the first summand on the right, \Cref{lem:series-expansions}\,(2) gives,
  for odd $m$,
  \[
    m\,[z^m]\bigl(-\log\det(I-zA)\bigr)=\Tr(A^m),
    \qquad
    m\,[z^m]\bigl(-\log\det(I+zA)\bigr)=\Tr\bigl((-A)^m\bigr)=-\Tr(A^m),
  \]
  so the two cancel when the identities for $U$ and $W$ are added.  For the
  second summand,
  \[
    m\,[z^m]\bigl(-\log(1-qz)\bigr)=q^m,
    \qquad
    m\,[z^m]\bigl(-\log(1-pz)\bigr)=p^m,
  \]
  and the third summands give $m[z^m]\bigl(L_-(z)+L_+(z)\bigr)=S_m$ by
  \eqref{eq:pge23-Sm}.  Adding the two identities therefore proves
  \eqref{eq:pge23-two-sided-identity}.

  To obtain the second identity, rearrange the first:
  \begin{align*}
    t(C_m,W)-\bigl(p^m-pq^{m-1}\bigr)
    &=q^m+S_m-t(C_m,U)+pq^{m-1}\\
    &=(p+q)q^{m-1}+S_m-t(C_m,U)\\
    &=q^{m-1}+S_m-t(C_m,U),
  \end{align*}
  using $p+q=1$. 
\end{proof}

\subsection{Replacing a complement cycle by a path}

Let $P_j$ denote the path with $j$ edges and set
\begin{equation}\label{eq:pge23-xj}
  x_j=t(P_j,U)=\ip\one{T_U^j\one}.
\end{equation}
As in \Cref{sec:short-cycles}, deleting one edge from the $m$-cycle leaves a
path with $m-1$ edges, so $0\le U\le1$ gives
\begin{equation}\label{eq:pge23-cycle-path}
  t(C_m,U)\le x_{m-1}.
\end{equation}

By \eqref{eq:pge23-one-sided-identity} and \eqref{eq:pge23-cycle-path}, it is enough to
prove
\begin{equation}\label{eq:pge23-Phi}
  \Phi_m:=q^{m-1}+S_m-x_{m-1}\ge0.
\end{equation}
Two of the three terms are already expressed through $q$ and $R_\pm$, so it
remains to do the same for the path densities.  By
\Cref{lem:series-expansions}\,(1),
\[
  \sum_{j\ge0}x_jz^j
  =\sum_{j\ge0}z^j\ip\one{T_U^j\one}
  =\ip\one{(I-zT_U)^{-1}\one}.
\]
Write $(I-zT_U)^{-1}\one=a\one+h$ with $h\in\mathcal E_0$, so that
$\ip\one{(I-zT_U)^{-1}\one}=a$.  In the block form of $I-zT_U$ this is
\[
  \begin{pmatrix}1-qz&-zg^*\\-zg&I-zA\end{pmatrix}
  \begin{pmatrix}a\\h\end{pmatrix}
  =\begin{pmatrix}1\\0\end{pmatrix}.
\]
Multiply on the left by the triangular matrix of
\eqref{eq:block-determinant-identity}, here
$\left(\begin{smallmatrix}1&zg^*(I-zA)^{-1}\\0&I\end{smallmatrix}\right)$.
It fixes $\left(\begin{smallmatrix}1\\0\end{smallmatrix}\right)$ and
produces the Schur complement in the upper left corner, so the system
becomes
\[
  \begin{pmatrix}1-qz-z^2R_+(z)&0\\-zg&I-zA\end{pmatrix}
  \begin{pmatrix}a\\h\end{pmatrix}
  =\begin{pmatrix}1\\0\end{pmatrix}.
\]
Its first row shows $\bigl(1-qz-z^2R_+(z)\bigr)a=1$, whence
\begin{equation}\label{eq:pge23-path-gf}
  \sum_{j\ge0}x_jz^j=\frac{1}{1-qz-z^2R_+(z)}.
\end{equation}

\subsection{Expansion in complete homogeneous polynomials}

In this section, we will give the explicit form of $\Phi_m$ as an integral of a polynomial in $q$ and the eigenvalues of $A$.

For variables $y_1,\ldots,y_k$, the \emph{complete homogeneous symmetric
polynomial} of degree $d$ is
\[
  h_d(y_1,\ldots,y_k)
  =\sum_{\substack{\alpha_1,\ldots,\alpha_k\ge0\\
  \alpha_1+\cdots+\alpha_k=d}}
  y_1^{\alpha_1}\cdots y_k^{\alpha_k},
  \qquad h_0=1,
\]
the sum of all monomials of degree $d$ in $y_1,\ldots,y_k$.  Equivalently,
\begin{equation}\label{eq:pge23-h-gf}
  \prod_{i=1}^k\frac{1}{1-y_i z}
  =\sum_{d\ge0}h_d(y_1,\ldots,y_k)z^d.
\end{equation}
The notation $a^{\times r}$ means that $a$ is repeated $r$ times.

\begin{proposition}[Exact coefficient expansion]\label{prop:pge23-expansion}
  For odd $m\ge3$,
  \begin{equation}\label{eq:pge23-expansion}
    \Phi_m=\sum_{r=1}^{(m-1)/2}
    \int\!\cdots\!\int
    P_{m,r}(q;\lambda_1,\ldots,\lambda_r)
    \prod_{i=1}^r\dd\mu(\lambda_i),
  \end{equation}
  where $n=m-2r$ and
  \begin{align}
    P_{m,r}(q;\lambda_1,\ldots,\lambda_r)
    ={}&\frac {m}{r}\Bigl[
    h_n(p^{\times r},-\lambda_1,\ldots,-\lambda_r)
    +h_n(q^{\times r},\lambda_1,\ldots,\lambda_r)
    \Bigr]\notag\\
    &-h_{n-1}(q^{\times(r+1)},\lambda_1,\ldots,\lambda_r).
    \label{eq:pge23-Pmr}
  \end{align}
\end{proposition}

\begin{proof}
  Put
  \[
    Y_-(z)=\frac{z^2R_-(z)}{1-pz},
    \qquad
    Y_+(z)=\frac{z^2R_+(z)}{1-qz}.
  \]
  From $-\log(1-Y)=\sum_{r\ge1}Y^r/r$,
  \begin{align*}
    m[z^m]L_-(z) &=m[z^m] (- \log (1-Y_-(z)))
    \\
    {} &= m[z^m]\sum_{r\ge1}\frac{Y_-(z)^r}{r}
    \\
    {} &=\sum_{r\ge1}\frac {m}{r} [z^{m-2r}]\frac{R_-(z)^r}{(1-pz)^r}.
  \end{align*}
  The spectral representation of $R_-(z) = \int \frac{d \mu(\lambda)}{1 + \lambda z}$ gives
  \begin{align*}
    R_-(z)^r &=\int\!\cdots\!\int
    \prod_{i=1}^r\frac{1}{1+\lambda_i z}
    \prod_{i=1}^r\dd\mu(\lambda_i),
    \\
    \frac{R_-(z)^r}{(1-pz)^r} &= \int\!\cdots\!\int \frac{1}{(1-pz)^r} \prod_{i=1}^r \frac{1}{1+\lambda_i z} \prod_{i=1}^r \dd\mu(\lambda_i).
  \end{align*}
  By \eqref{eq:pge23-h-gf}, the $r$-th coefficient is 
  \[
    [z^{m-2r}] \frac{R_-(z)^r}{(1-pz)^r} = \int\!\cdots\!\int h_{m-2r}(p^{\times r}, -\lambda_1, \ldots, -\lambda_r) \prod_{i=1}^r \dd\mu(\lambda_i).
  \]
  The same argument for $L_+$ produces the second term in brackets in
  \eqref{eq:pge23-Pmr}.

  On the other hand, factoring $1-qz$ out of the denominator of
  \eqref{eq:pge23-path-gf} and expanding the geometric series gives
  \begin{align*}
    \sum_{j\ge0}x_jz^j
    &=\frac{1}{(1-qz)\left(1-\dfrac{z^2R_+(z)}{1-qz}\right)}
    \\
    {}&=\sum_{r\ge0}\frac{z^{2r}R_+(z)^r}{(1-qz)^{r+1}},
  \end{align*}
  so that
  \[
    x_{m-1}=\sum_{r\ge0}[z^{m-1-2r}]\frac{R_+(z)^r}{(1-qz)^{r+1}}.
  \]
  The term $r=0$ is $[z^{m-1}](1-qz)^{-1}=q^{m-1}$, which cancels the first
  term of \eqref{eq:pge23-Phi}.  For $r\ge1$, the spectral representation of
  $R_+(z)=\int\frac{\dd\mu(\lambda)}{1-\lambda z}$ gives, as above,
  \[
    \frac{R_+(z)^r}{(1-qz)^{r+1}}
    =\int\!\cdots\!\int\frac{1}{(1-qz)^{r+1}}
    \prod_{i=1}^r\frac{1}{1-\lambda_iz}
    \prod_{i=1}^r\dd\mu(\lambda_i),
  \]
  and \eqref{eq:pge23-h-gf}, applied now to the $r+1$ copies of $q$ together
  with $\lambda_1,\ldots,\lambda_r$, yields
  \[
    [z^{m-1-2r}]\frac{R_+(z)^r}{(1-qz)^{r+1}}
    =\int\!\cdots\!\int
    h_{n-1}(q^{\times(r+1)},\lambda_1,\ldots,\lambda_r)
    \prod_{i=1}^r\dd\mu(\lambda_i),
  \]
  because $m-1-2r=n-1$.  Subtracting these terms from the contributions of
  $L_-$ and $L_+$ gives \eqref{eq:pge23-expansion}.  All three sums are
  finite and stop at $r=(m-1)/2$, since their coefficients are zero once
  $2r>m-1$.
\end{proof}

\subsection{Diagonalization by a Dirichlet average}

In this section, we will simplify the integrand $P_{m,r}$ of \eqref{eq:pge23-expansion} by replacing the eigenvalues $\lambda_1,\ldots,\lambda_r$ by a single variable $\ell$.
Define the diagonal polynomial
\begin{equation}\label{eq:pge23-Ptilde}
  \widetilde P_{m,r}(q,\ell)
  =P_{m,r}(q;\ell,\ldots,\ell).
\end{equation}

\begin{lemma}[Dirichlet diagonalization]\label{lem:pge23-dirichlet}
  Let $(\Theta_1,\ldots,\Theta_r)$ be uniformly distributed on the simplex
  \[
    \Theta_i\ge0,\qquad \sum_{i=1}^r\Theta_i=1.
  \]
  Then
  \begin{equation}\label{eq:pge23-dirichlet-identity}
    P_{m,r}(q;\lambda_1,\ldots,\lambda_r)
    =\E\,\widetilde P_{m,r}
    \left(q,\sum_{i=1}^r\Theta_i\lambda_i\right).
  \end{equation}
  Consequently,
  \begin{equation}\label{eq:pge23-diagonal-min}
    \min_{\lambda_i\in[-1/2,1/2]}
    P_{m,r}(q;\lambda_1,\ldots,\lambda_r)
    =\min_{\ell\in[-1/2,1/2]}\widetilde P_{m,r}(q,\ell).
  \end{equation}
\end{lemma}

\begin{proof}
  Write $S=\sum_{i=1}^r\Theta_i\lambda_i$.  The proof rests on two properties
  of $h_j$.

  First, the moments of $S$ are the complete homogeneous polynomials, up to a
  constant:
  \begin{equation}\label{eq:pge23-dirichlet-moment}
    \E\,S^j
    =\frac{h_j(\lambda_1,\ldots,\lambda_r)}
    {\binom{j+r-1}{r-1}}.
  \end{equation}
  To see this, expand
  \[
    S^j=\Bigl(\sum_{i=1}^r\Theta_i\lambda_i\Bigr)^j
    =\sum_{\alpha_1+\cdots+\alpha_r=j}
    \frac{j!}{\alpha_1!\cdots\alpha_r!}
    \prod_{i=1}^r\Theta_i^{\alpha_i}\lambda_i^{\alpha_i},
  \]
  and take expectations.  From
  \[
    \E\prod_{i=1}^r\Theta_i^{\alpha_i}
    =\frac{(r-1)!\,\alpha_1!\cdots\alpha_r!}{(j+r-1)!},
    \qquad \alpha_1+\cdots+\alpha_r=j,
  \]
  we obtain
  \begin{align*}
    \E\,S^j
    &=\sum_{\alpha_1+\cdots+\alpha_r=j}
    \frac{j!}{\alpha_1!\cdots\alpha_r!}
    \cdot\frac{(r-1)!\,\alpha_1!\cdots\alpha_r!}{(j+r-1)!}
    \prod_{i=1}^r\lambda_i^{\alpha_i}\\
    &=\frac{j!\,(r-1)!}{(j+r-1)!}
    \sum_{\alpha_1+\cdots+\alpha_r=j}\prod_{i=1}^r\lambda_i^{\alpha_i}
    =\frac{h_j(\lambda_1,\ldots,\lambda_r)}{\binom{j+r-1}{r-1}}.
  \end{align*}

  Second, setting all variables equal in the definition of $h_j$ gives the
  diagonal value
  \begin{equation}\label{eq:pge23-h-diagonal}
    h_j(\ell,\ldots,\ell)=\binom{j+r-1}{r-1}\ell^j,
  \end{equation}
  the binomial coefficient being the number of monomials of degree $j$ in $r$
  variables.

  It remains to express the terms of
  $P_{m,r}$ through the
  $h_j(\lambda_1,\ldots,\lambda_r)$.  In the generating function
  \eqref{eq:pge23-h-gf} of the $h_d$, the $k$ equal variables and the
  remaining ones give separate factors, so
  \[
    \sum_{d\ge0}h_d(a^{\times k},b_1,\ldots,b_r)z^d
    =\frac{1}{(1-az)^k}\prod_{i=1}^r\frac{1}{1-b_iz}
    =\Bigl(\sum_{s\ge0}\binom{s+k-1}{k-1}a^sz^s\Bigr)
    \Bigl(\sum_{j\ge0}h_j(b_1,\ldots,b_r)z^j\Bigr),
  \]
  and comparing the coefficients of $z^d$ gives the splitting identity
  \begin{equation}\label{eq:pge23-h-split}
    h_d(a^{\times k},b_1,\ldots,b_r)
    =\sum_{j=0}^d
    \binom{d-j+k-1}{k-1}a^{d-j}h_j(b_1,\ldots,b_r).
  \end{equation}
  Apply this to the three terms of
  $P_{m,r}$, the first two with $k=r$ and $a=p$ or $q$, the third with
  $k=r+1$ and $a=q$.  Since $h_j$ is homogeneous of degree $j$, we have
  $h_j(-\lambda_1,\ldots,-\lambda_r)=(-1)^jh_j(\lambda_1,\ldots,\lambda_r)$,
  and therefore
  \begin{align}
    P_{m,r}(q;\lambda_1,\ldots,\lambda_r)
    ={}&\frac {m}{r}\sum_{j=0}^{n}\binom{n-j+r-1}{r-1}
    \bigl((-1)^jp^{n-j}+q^{n-j}\bigr)
    h_j(\lambda_1,\ldots,\lambda_r) \notag\\
    &-\sum_{j=0}^{n-1}\binom{n-1-j+r}{r}q^{n-1-j}
    h_j(\lambda_1,\ldots,\lambda_r). \label{eq:pge23-P-in-h}
  \end{align}

  Write the expansion \eqref{eq:pge23-P-in-h} as
  \[
    P_{m,r}(q;\lambda_1,\ldots,\lambda_r)
    =\sum_{j=0}^nc_j\,h_j(\lambda_1,\ldots,\lambda_r),
  \]
  where $c_j$ subsumes the two coefficients of $h_j$ above, which do not
  involve the $\lambda_i$.  Setting $\lambda_1=\cdots=\lambda_r=\ell$ and
  using the diagonal value \eqref{eq:pge23-h-diagonal},
  \begin{equation}\label{eq:pge23-Ptilde-in-c}
    \widetilde P_{m,r}(q,\ell)
    =\sum_{j=0}^nc_j\binom{j+r-1}{r-1}\ell^j.
  \end{equation}
  On the other hand, the moment identity \eqref{eq:pge23-dirichlet-moment}
  gives
  \[
    P_{m,r}(q;\lambda_1,\ldots,\lambda_r)
    =\sum_{j=0}^nc_j\binom{j+r-1}{r-1}\E\,S^j
    =\E\sum_{j=0}^nc_j\binom{j+r-1}{r-1}S^j
    =\E\,\widetilde P_{m,r}(q,S),
  \]
  the last equality being $\widetilde P_{m,r}$ with $\ell$ replaced by $S$.

  For the second claim, if the $\lambda_i$ lie in $[-1/2,1/2]$ then so does every convex
  combination $S$, so the right side of the identity
  \eqref{eq:pge23-dirichlet-identity} is
  at least $\min_\ell\widetilde P_{m,r}(q,\ell)$; conversely
  $\lambda_1=\cdots=\lambda_r=\ell$ attains each diagonal value.
\end{proof}

Thus \eqref{eq:pge23-expansion} will be nonnegative once we prove
\begin{equation}\label{eq:pge23-diagonal-target}
  \widetilde P_{m,r}(q,\ell)\ge0
\end{equation}
for
\[
  0\le q\le\frac{1}{3},
  \quad 1\le r\le\frac{m-1}{2},
  \quad -\frac{1}{2}\le\ell\le\frac{1}{2}.
\]
This is the first major simplification: an $r$-variable symmetric
polynomial is replaced by one scalar variable polynomial.

\subsection{A beta representation}

In this section, we will further simplify the diagonal polynomial $\widetilde P_{m,r}$ by representing it as an expectation over a beta distribution.

\begin{lemma}[Beta formula]\label{lem:pge23-beta}
  Let $1\le r\le\frac{m-1}{2}$ and put
  \begin{equation}\label{eq:pge23-n}
    n=m-2r,
  \end{equation}
  so that $n\ge1$ is odd.  Let $X\sim\BetaLaw(r,r)$ and
  \begin{equation}\label{eq:pge23-rho}
    \rho_{n,m}(u)
    =\frac {m}{n}\bigl(u^n+(1-u)^n\bigr)-u^{n-1}.
  \end{equation}
  Then
  \begin{equation}\label{eq:pge23-beta-second}
    \widetilde P_{m,r}(q,\ell)
    =\binom{n+2r-1}{2r-1}\frac {n}{r}
    \E\left[
    X^n\rho_{n,m}\left(q+\ell\frac{1-X}{X}\right)
    \right].
  \end{equation}
  If $\ell>0$, then
  \begin{equation}\label{eq:pge23-beta-integral}
    \widetilde P_{m,r}(q,\ell)
    =\frac{\Gamma(m)}{r\Gamma(n)\Gamma(r)^2}\,
    \ell^{n+r}
    \int_0^\infty
    \frac{s^{r-1}}{(\ell+s)^m}\rho_{n,m}(q+s)\dd s.
  \end{equation}
\end{lemma}

\begin{proof}
  For arbitrary $a,b$ we first claim that
  \begin{equation}\label{eq:pge23-beta-h}
    h_n(a^{\times r},b^{\times r})
    =\binom{n+2r-1}{2r-1}
    \E\bigl(aX+b(1-X)\bigr)^n.
  \end{equation}
  Both sides are polynomials in $a$ and $b$, so it suffices to compare the
  coefficient of $a^kb^{n-k}$.  On the left, the splitting identity
  \eqref{eq:pge23-h-split} and the diagonal value
  \eqref{eq:pge23-h-diagonal} give
  \[
    h_n(a^{\times r},b^{\times r})
    =\sum_{k=0}^n h_k(a^{\times r})\,h_{n-k}(b^{\times r})
    =\sum_{k=0}^n\binom{k+r-1}{r-1}\binom{n-k+r-1}{r-1}a^kb^{n-k}.
  \]
  On the right, the binomial theorem and the beta moments
  \[
    \E\bigl[X^k(1-X)^{n-k}\bigr]
    =\frac{\Gamma(2r)}{\Gamma(r)^2}
    \cdot\frac{\Gamma(r+k)\,\Gamma(r+n-k)}{\Gamma(2r+n)}
    =\frac{(2r-1)!}{((r-1)!)^2}
    \cdot\frac{(k+r-1)!\,(n-k+r-1)!}{(n+2r-1)!}
  \]
  give
  \[
    \binom{n+2r-1}{2r-1}\E\bigl(aX+b(1-X)\bigr)^n
    =\sum_{k=0}^n\binom{n+2r-1}{2r-1}\binom nk
    \E\bigl[X^k(1-X)^{n-k}\bigr]a^kb^{n-k}.
  \]
  The two coefficients agree, since writing every factor in factorials turns
  both into
  $(r+k-1)!\,(r+n-k-1)!/\bigl(k!\,(n-k)!\,((r-1)!)^2\bigr)$; on the right the
  factors $(2r-1)!$, $n!$, and $(2r+n-1)!$ cancel against those of
  $\binom{n+2r-1}{2r-1}\binom nk$.

  Next we differentiate both sides of \eqref{eq:pge23-beta-h} with respect to $a$.
  On the left, differentiating the generating function
  \eqref{eq:pge23-h-gf} in the form
  $\sum_{d\ge0}h_d(a^{\times r},b^{\times r})z^d=(1-az)^{-r}(1-bz)^{-r}$
  gives
  \[
    \sum_{d\ge0}\frac{\partial h_d(a^{\times r},b^{\times r})}{\partial a}z^d
    =rz\,(1-az)^{-(r+1)}(1-bz)^{-r}
    =r\sum_{d\ge0}h_d(a^{\times(r+1)},b^{\times r})z^{d+1},
  \]
  so that
  $\partial_a h_n(a^{\times r},b^{\times r})
  =r\,h_{n-1}(a^{\times(r+1)},b^{\times r})$.
  On the right, $\partial_a\bigl(aX+b(1-X)\bigr)^n
  =nX\bigl(aX+b(1-X)\bigr)^{n-1}$.  Hence
  \begin{equation}\label{eq:pge23-beta-h-deriv}
    h_{n-1}(a^{\times(r+1)},b^{\times r})
    =\binom{n+2r-1}{2r-1}\frac {n}{r}
    \E\left[X\bigl(aX+b(1-X)\bigr)^{n-1}\right].
  \end{equation}

  By the definition \eqref{eq:pge23-Ptilde} of the diagonal polynomial,
  \[
    \widetilde P_{m,r}(q,\ell)
    =\frac {m}{r}\bigl[h_n(p^{\times r},(-\ell)^{\times r})
    +h_n(q^{\times r},\ell^{\times r})\bigr]
    -h_{n-1}(q^{\times(r+1)},\ell^{\times r}),
  \]
  and we apply \eqref{eq:pge23-beta-h} to the first two terms, with
  $(a,b)=(p,-\ell)$ and $(a,b)=(q,\ell)$, and
  \eqref{eq:pge23-beta-h-deriv} to the third, with $(a,b)=(q,\ell)$.
  Using $\frac {m}{r}=\frac {n}{r}\cdot\frac {m}{n}$, this gives
  \begin{align*}
    &\widetilde P_{m,r}(q,\ell)
    \\
    {} = &\binom{n+2r-1}{2r-1}\frac {n}{r}
    \left[
    \frac {m}{n}\E\Bigl(\bigl(pX-\ell(1-X)\bigr)^n
                       +\bigl(qX+\ell(1-X)\bigr)^n\Bigr)
    -\E\Bigl(X\bigl(qX+\ell(1-X)\bigr)^{n-1}\Bigr)
    \right].
  \end{align*}
  Abbreviating the two linear forms by
  \[
    V_X=qX+\ell(1-X),
    \qquad
    W_X=pX-\ell(1-X),
    \qquad\text{so that } V_X+W_X=X,
  \]
  this simplifies to
  \begin{equation}\label{eq:pge23-beta-first}
    \widetilde P_{m,r}(q,\ell)
    =\binom{n+2r-1}{2r-1}\frac {n}{r}
    \left[
    \frac {m}{n}\E\bigl(V_X^n+W_X^n\bigr)
    -\E\bigl(X V_X^{n-1}\bigr)
    \right].
  \end{equation}

  The three terms $V_X^n$, $W_X^n$, and $XV_X^{n-1}$ all have degree $n$ in
  $(V_X,W_X,X)$, and $V_X+W_X=X$.  They therefore depend on $X$ only
  through the proportion of $X$ taken by $V_X$.  Writing that proportion as
  \[
    u_X=\frac{V_X}{X}=q+\ell\,\frac{1-X}{X},
    \qquad
    V_X=Xu_X,
    \qquad
    W_X=X(1-u_X),
  \]
  the last equality again because $V_X+W_X=X$, we obtain
  \[
    V_X^n+W_X^n=X^n\bigl(u_X^n+(1-u_X)^n\bigr),
    \qquad
    XV_X^{n-1}=X\cdot X^{n-1}u_X^{n-1}=X^nu_X^{n-1}.
  \]
  By linearity, the two expectations in \eqref{eq:pge23-beta-first} combine
  into the expectation of
  \[
    \frac {m}{n}\bigl(V_X^n+W_X^n\bigr)-XV_X^{n-1}
    =X^n\left[\frac {m}{n}\bigl(u_X^n+(1-u_X)^n\bigr)-u_X^{n-1}\right]
    =X^n\rho_{n,m}(u_X),
  \]
  by the definition \eqref{eq:pge23-rho} of $\rho_{n,m}$, and this proves
  \eqref{eq:pge23-beta-second}.

  For $\ell>0$, expand the expectation in \eqref{eq:pge23-beta-second} as an
  integral against the beta density,
  \[
    \E\left[X^n\rho_{n,m}\left(q+\ell\frac{1-X}{X}\right)\right]
    =\frac{\Gamma(2r)}{\Gamma(r)^2}
    \int_0^1 x^{n+r-1}(1-x)^{r-1}
    \rho_{n,m}\left(q+\ell\frac{1-x}{x}\right)\dd x,
  \]
  and apply the change of variable
  \[
    s=\ell\frac{1-x}{x},
    \qquad x=\frac{\ell}{\ell+s},
    \qquad 1-x=\frac {s}{\ell+s},
    \qquad \dd x=-\frac{\ell}{(\ell+s)^2}\dd s,
  \]
  under which $x=1$ corresponds to $s=0$, and $x\to0^+$ to
  $s\to\infty$.  Thus
  \begin{align*}
    &\int_0^1 x^{n+r-1}(1-x)^{r-1}
     \rho_{n,m}\left(q+\ell\frac{1-x}{x}\right)\dd x\\
    &\qquad=\int_\infty^0
     \left(\frac{\ell}{\ell+s}\right)^{n+r-1}
     \left(\frac{s}{\ell+s}\right)^{r-1}
     \rho_{n,m}(q+s)
     \left(-\frac{\ell}{(\ell+s)^2}\right)\dd s\\
    &\qquad=\int_0^\infty
     \frac{\ell^{n+r}s^{r-1}}{(\ell+s)^{n+2r}}\rho_{n,m}(q+s)\dd s
     =\ell^{n+r}\int_0^\infty
     \frac{s^{r-1}}{(\ell+s)^{m}}\rho_{n,m}(q+s)\dd s,
  \end{align*}
  where in the last line the minus sign has been used to interchange the
  endpoints of integration, and the exponent $n+2r$ has been replaced by
  $m$ according to \eqref{eq:pge23-n}.  Altogether
  \[
    \widetilde P_{m,r}(q,\ell)
    =\binom{m-1}{2r-1}\frac {n}{r}\frac{\Gamma(2r)}{\Gamma(r)^2}\,
     \ell^{n+r}\int_0^\infty
     \frac{s^{r-1}}{(\ell+s)^{m}}\rho_{n,m}(q+s)\dd s,
  \]
  and since
  \[
    \binom{m-1}{2r-1}\frac {n}{r}
    \frac{\Gamma(2r)}{\Gamma(r)^2}
    =\frac{(m-1)!}{(2r-1)!\,n!}\cdot\frac {n}{r}
     \cdot\frac{(2r-1)!}{((r-1)!)^2}
    =\frac{\Gamma(m)}{r\Gamma(n)\Gamma(r)^2},
  \]
  this proves \eqref{eq:pge23-beta-integral}.
\end{proof}

\subsection{The case \texorpdfstring{$\ell\le0$}{l<=0}: pointwise nonnegativity}

\begin{lemma}\label{lem:pge23-ell-negative}
  If $0\le q\le1/2$ and $\ell\le0$, then
  $\widetilde P_{m,r}(q,\ell)\ge0$.
\end{lemma}

\begin{proof}
  From the definition \eqref{eq:pge23-rho} of $\rho_{n,m}$,
  \begin{align*}
    \rho_{n,m}(u)
    &=\frac {m}{n}\bigl(u^n+(1-u)^n\bigr)-u^{n-1} \nonumber \\
    &=\frac {m}{n}(1-u)\bigl((1-u)^{n-1}-u^{n-1}\bigr)
    +\left(\frac {m}{n}-1\right)u^{n-1},
  \end{align*}
  which can be shown by expanding
  $(1-u)\bigl((1-u)^{n-1}-u^{n-1}\bigr)=(1-u)^n-u^{n-1}+u^n$.  
  Recall that $\frac{m}{n}-1=\frac{2r}{n}>0$ by \eqref{eq:pge23-n}, and $n-1$ is even. 
  The second term is therefore always nonnegative, and the first is
  nonnegative whenever $1-u\ge0$ and $|1-u|\ge|u|$, which hold for $u \le 1/2$.

  Since $X\sim\BetaLaw(r,r)$ takes values in $(0,1)$, we have $(1-X)/X\ge0$, so $\ell\le0$
  gives
  \[
    q+\ell\,\frac{1-X}{X}\le q\le\frac{1}{2}.
  \]
  Hence $X^n\rho_{n,m}\bigl(q+\ell(1-X)/X\bigr)\ge0$, and
  \eqref{eq:pge23-beta-second} gives $\widetilde P_{m,r}(q,\ell)\ge0$.
\end{proof}

So only the case $\ell>0$ remains.
There the argument $q+s$ of $\rho_{n,m}$ in \eqref{eq:pge23-beta-integral} ranges over $[q,\infty)$, 
and unlike the $\ell\le0$ case, $\rho_{n,m}$ can be negative on that range: for instance $\rho_{5,7}(0.7)<0$.
The integrand can therefore no longer be handled pointwise, and nonnegativity has to be derived from the integral as a whole.

\subsection{Laplace--gamma representation}

Since the constant and the factor $\ell^{n+r}$ in
\eqref{eq:pge23-beta-integral} are positive, what has to be shown
nonnegative for $\ell>0$ is
\[
  \int_0^\infty
  \frac{s^{r-1}}{(\ell+s)^m}\rho_{n,m}(q+s)\dd s.
\]
Here $\ell$ and $s$ are coupled through $(\ell+s)^{-m}$, which is what
makes the integral difficult to handle.  They can be separated by an
inverse Laplace transform: the inverse Laplace transform of
$x\mapsto x^{-m}$ is $t\mapsto t^{m-1}/\Gamma(m)$, so that
\begin{equation}\label{eq:pge23-laplace-kernel}
  \frac{1}{(\ell+s)^m}
  =\frac{1}{\Gamma(m)}\int_0^\infty t^{m-1}e^{-t(\ell+s)}\dd t,
\end{equation}
which can be shown by evaluating the integral under the substitution
$w=t(\ell+s)$.

\begin{lemma}[Laplace--gamma representation]\label{lem:pge23-laplace-gamma}
  Let $Y\sim\GammaLaw(r,1)$.  If $\ell>0$, then
  \begin{equation}\label{eq:pge23-laplace-gamma}
    \widetilde P_{m,r}(q,\ell)
    =\frac{\ell^{n+r}}{r\Gamma(n)\Gamma(r)}
    \int_0^\infty
    t^{m-r-1}e^{-\ell t}
    \E\rho_{n,m}\left(q+\frac {Y}{t}\right)\dd t.
  \end{equation}
\end{lemma}

\begin{proof}
  Insert \eqref{eq:pge23-laplace-kernel} into \eqref{eq:pge23-beta-integral} and
  exchange the two integrals:
  \begin{align*}
    \widetilde P_{m,r}(q,\ell)
    &=\frac{\Gamma(m)}{r\Gamma(n)\Gamma(r)^2}\,\ell^{n+r}
      \int_0^\infty s^{r-1}\rho_{n,m}(q+s)
      \left(\frac{1}{\Gamma(m)}\int_0^\infty t^{m-1}e^{-t(\ell+s)}\dd t\right)
      \dd s\\
    &=\frac{\ell^{n+r}}{r\Gamma(n)\Gamma(r)^2}
      \int_0^\infty t^{m-1}e^{-\ell t}
      \left(\int_0^\infty s^{r-1}e^{-ts}\rho_{n,m}(q+s)\dd s\right)\dd t.
  \end{align*}
  This exchange is justified by Fubini's theorem, given the absolute
  integrability of the double integral.  To verify absolute integrability,
  integrate the absolute value in $t$ first: by \eqref{eq:pge23-laplace-kernel},
  \[
    \int_0^\infty\!\!\int_0^\infty
    s^{r-1}t^{m-1}e^{-t(\ell+s)}\bigl|\rho_{n,m}(q+s)\bigr|\dd t\dd s
    =\Gamma(m)\int_0^\infty \frac{s^{r-1}}{(\ell+s)^m}
    \bigl|\rho_{n,m}(q+s)\bigr|\dd s.
  \]
  It remains to show that the integral on the right is finite.  Recall that
  $\rho_{n,m}$ is a polynomial of degree $n-1$, since the terms $u^n$ and
  $-u^n$ of $u^n+(1-u)^n$ cancel for odd $n$.  As $s\to\infty$, the
  integrand is therefore $O\bigl(s^{r+n-2-m}\bigr)=O\bigl(s^{-r-2}\bigr)$
  by \eqref{eq:pge23-n}, which is integrable at infinity.  As $s\to0$, the
  integrand is $O\bigl(s^{r-1}\bigr)$, which is integrable because
  $r\ge1$.

  It remains to compute the inner integral.  Substituting $y=ts$, so that
  $s=y/t$ and $\dd s=\dd y/t$, gives
  \[
    \int_0^\infty s^{r-1}e^{-ts}\rho_{n,m}(q+s)\dd s
    =t^{-r}\int_0^\infty y^{r-1}e^{-y}
     \rho_{n,m}\left(q+\frac {y}{t}\right)\dd y
    =\Gamma(r)\,t^{-r}
     \E\rho_{n,m}\left(q+\frac {Y}{t}\right),
  \]
  where the last equality holds because $Y$ has density
  $y^{r-1}e^{-y}/\Gamma(r)$ on $(0,\infty)$.  Substituting this for the
  inner integral obtained after the exchange gives
  \eqref{eq:pge23-laplace-gamma}.
\end{proof}

In \eqref{eq:pge23-laplace-gamma} the parameter $\ell$ occurs only in the
factor $\ell^{n+r}$ and in the weight $t^{m-r-1}e^{-\ell t}$, both of which
are positive; the expectation $\E\rho_{n,m}(q+Y/t)$ does not involve $\ell$
at all.  It is therefore enough to prove that this expectation is
nonnegative for every $t>0$, and the conclusion then follows for all
$\ell>0$ at once.  This is what the next two subsections do.

\subsection{The shifted-gamma moment inequality}

To prove that $\E\rho_{n,m}(q+xY)$ is nonnegative, we center $\rho_{n,m}$ at $u=1/2$ and expand.
This produces the terms with positive constant factors
\[
  2(r+j)v^{2j}-jv^{2j-1},
  \qquad j\ge1,
\]
evaluated at the random $v=q+xY-1/2$.  
In each of them, the even power has nonnegative expectation, while the odd power has no definite sign; 
what is needed is therefore a bound of the odd moment by the even one. 
The following lemma supplies such a bound, uniformly in the scaling of $Y$.

\begin{lemma}[Uniform odd-to-even gamma moment bound]\label{lem:pge23-gamma-moment}
  Let $r,j\ge1$ be integers and $Y\sim\GammaLaw(r,1)$.  For every
  $z\ge0$,
  \begin{equation}\label{eq:pge23-gamma-moment}
    3j\,\E(zY-1)^{2j-1}
    \le(r+j)\,\E(zY-1)^{2j}.
  \end{equation}
\end{lemma}

We postpone its proof to the next subsection and first show exactly how it proves the nonnegativity of $\E\rho_{n,m}(q+xY)$.

\begin{proposition}[Shifted-gamma positivity]\label{prop:pge23-gamma-positive}
  Let $n\ge1$ be odd, $r\ge1$, and $m=n+2r$.  If $0\le q\le1/3$, then
  for every $x\ge0$,
  \begin{equation}\label{eq:pge23-shifted-positive}
    \E\rho_{n,m}(q+xY)\ge0,
    \qquad Y\sim\GammaLaw(r,1).
  \end{equation}
\end{proposition}

\begin{proof}
  Set
  \[
    A_n(u)=u^n+(1-u)^n.
  \]
  We rewrite $\rho_{n,m}$ through $A_n$ and $A_n'$. 
  From $A_n'(u)=n\bigl(u^{n-1}-(1-u)^{n-1}\bigr)$,
  \[
    (1-u)A_n'(u)
    =n(1-u)u^{n-1}-n(1-u)^n
    =nu^{n-1}-nu^n-n(1-u)^n
    =nu^{n-1}-nA_n(u),
  \]
  that is, $u^{n-1}=A_n(u)+\frac1n(1-u)A_n'(u)$. 
  Substituting this into the definition \eqref{eq:pge23-rho} of $\rho_{n,m}$ and using $m-n=2r$,
  \begin{equation}\label{eq:pge23-rho-derivative}
    n\rho_{n,m}(u)=2rA_n(u)-(1-u)A_n'(u).
  \end{equation}

  Now expand around $u=1/2$. 
  With $u=\frac12+v$, the binomial theorem gives
  \begin{equation}\label{eq:pge23-even-expansion}
    A_n\left(\frac{1}{2}+v\right)
    =\left(\frac12+v\right)^n+\left(\frac12-v\right)^n
    =2\sum_{2j\le n}\binom n{2j}\left(\frac12\right)^{n-2j}v^{2j}
    =\sum_{j=0}^{(n-1)/2}a_jv^{2j},
  \end{equation}
  where $a_j=2^{2j+1-n}\binom n{2j}>0$. 
  Differentiating in $v$, we also have $A_n'\bigl(\frac12+v\bigr)=\sum_{j\ge1}2ja_jv^{2j-1}$. 
  Substituting both into \eqref{eq:pge23-rho-derivative}, with $1-u=\frac12-v$, gives
  \begin{align}
    n\rho_{n,m}\left(\frac{1}{2}+v\right)
    &=2r\sum_{j\ge0}a_jv^{2j}
     -\left(\frac12-v\right)\sum_{j\ge1}2ja_jv^{2j-1} \nonumber \\
    &=2ra_0+
     \sum_{j=1}^{(n-1)/2}
     a_j\bigl(2rv^{2j}-jv^{2j-1}+2jv^{2j}\bigr) \nonumber \\
    &=2ra_0+
     \sum_{j=1}^{(n-1)/2}
     a_j\bigl(2(r+j)v^{2j}-jv^{2j-1}\bigr).\label{eq:pge23-rho-centered}
  \end{align}
  
  At $u=q+xY$, the $v$ in \eqref{eq:pge23-rho-centered} equals
  \[
    v=q+xY-\frac{1}{2}=xY-\delta=\delta(zY-1),
    \quad \text{where} \quad
    \delta=\frac{1}{2}-q\ge\frac{1}{2}-\frac{1}{3}=\frac{1}{6},
    \qquad
    z=\frac{x}{\delta},
  \]
  which is the form occurring in \Cref{lem:pge23-gamma-moment}.  Taking
  expectations in \eqref{eq:pge23-rho-centered}, it suffices to bound each
  summand.  For $j\ge1$,
  \begin{align*}
    &\E\left[
    2(r+j)\bigl(\delta(zY-1)\bigr)^{2j}
    -j\bigl(\delta(zY-1)\bigr)^{2j-1}
    \right]\\
    &\qquad=2(r+j)\delta^{2j}\E(zY-1)^{2j}
      -j\delta^{2j-1}\E(zY-1)^{2j-1}\\
    &\qquad\ge2(r+j)\delta^{2j}\E(zY-1)^{2j}
      -\frac{r+j}{3}\delta^{2j-1}\E(zY-1)^{2j}\\
    &\qquad=\delta^{2j}(r+j)
    \left(2-\frac{1}{3\delta}\right)
    \E(zY-1)^{2j}
    \ge0,
  \end{align*}
  where the inequality is \Cref{lem:pge23-gamma-moment} and the last step uses
  $2\delta\ge1/3$ together with $\E(zY-1)^{2j}\ge0$.  Since every $a_j$ is
  positive, we conclude
  \[
    n\,\E\rho_{n,m}(q+xY)\ge2ra_0>0,
  \]
  which proves \eqref{eq:pge23-shifted-positive}.
\end{proof}

\begin{remark}[Tightness of the density requirement $p \ge 2/3$]\label{rem:pge23-threshold}
  This is the only place in the proof where the restriction $p\ge 2/3$ is used.
  As this inequality is tight, this proof only works for $p \ge 2/3$, and we require a different argument for $p<2/3$.
\end{remark}

\subsection{Proof of the gamma moment inequality}\label{subsec:pge23-gamma-moment-proof}

It remains to prove \Cref{lem:pge23-gamma-moment}: for integers $r,j\ge1$, for
$Y\sim\GammaLaw(r,1)$, and for every $z\ge0$,
\begin{equation}
  3j\,\E(zY-1)^{2j-1}
  \le(r+j)\,\E(zY-1)^{2j}. \label{eq:pge23-gamma-moment2}
\end{equation}

\begin{proof}[Proof of \Cref{lem:pge23-gamma-moment}]
  The case $z=0$ is immediate, since the left hand side equals $-3j < 0$ and the right hand side equals $r+j > 0$.
  
  Now assume $z>0$, put $b=1/z$, and define
  \[
    F(b)=\E(Y-b)^{2j},
    \qquad\text{so that}\qquad
    F'(b)=-2j\,\E(Y-b)^{2j-1}.
  \]
  Since $zY-1=z(Y-b)$, dividing \eqref{eq:pge23-gamma-moment2} by $z^{2j}>0$ gives
  the equivalent inequality
  \begin{equation}\label{eq:pge23-gamma-H}
    H(b):=(r+j)F(b)+\frac{3b}{2}F'(b)\ge0.
  \end{equation}

  For $I_k(b)=\E(Y-b)^k$, apply the Stein identity of the gamma
  distribution $Y\sim\GammaLaw(r,1)$,
  \begin{equation}\label{eq:pge23-gamma-Stein}
    \E\bigl[Yf'(Y)\bigr]=\E\bigl[(Y-r)f(Y)\bigr],
  \end{equation}
  to $f(y)=(y-b)^k$.  Its two sides are
  \begin{align*}
    \E[Y f'(Y)] &= k\E\bigl[(Y-b)(Y-b)^{k-1}+b(Y-b)^{k-1}\bigr] = kI_k(b)+kbI_{k-1}(b),\\
    \E[(Y-r)f(Y)] &= \E\bigl[(Y-b)(Y-b)^k + (b-r)(Y-b)^k\bigr] = I_{k+1}(b)-(r-b)I_k(b),
  \end{align*}
  so that
  \begin{equation}\label{eq:pge23-gamma-recurrence}
    I_{k+1}(b)=(r+k-b)I_k(b)+kbI_{k-1}(b).
  \end{equation}
  Taking $k=2j-1$, and using $I_{k+1} = F$, $I_{k}=-\frac{F'}{2j}$ and
  $I_{k-1}=\frac{F''}{2j(2j-1)}$, gives the differential equation
  \begin{equation}\label{eq:pge23-gamma-ODE}
    bF''(b)+(b-r-2j+1)F'(b)-2jF(b)=0.
  \end{equation}

  We next determine the sign of $H'$ at a zero of $H$.  Differentiating
  \eqref{eq:pge23-gamma-H} and then eliminating $F''$ by \eqref{eq:pge23-gamma-ODE}, we obtain
  \[
    H'(b)=\left(r+j+\frac{3}{2}\right)F'(b)+\frac{3b}{2}F''(b)
    =3jF(b)+\left(\frac{5r}{2}+4j-\frac{3b}{2}\right)F'(b).
  \]
  At a zero of $H$ we have $F'(b)=-\frac{2(r+j)}{3b}F(b)$ by
  \eqref{eq:pge23-gamma-H}, and substituting this leaves
  \begin{equation}\label{eq:pge23-gamma-crossing}
    \frac{H'(b)}{F(b)}
    =\frac{3(r+4j)b-(r+j)(5r+8j)}{3b}.
  \end{equation}
  Set
  \begin{equation}\label{eq:pge23-gamma-bstar}
    b_*:=\frac{(r+j)(5r+8j)}{3(r+4j)}.
  \end{equation}
  Since $F(b)=\E(Y-b)^{2j}>0$, formula \eqref{eq:pge23-gamma-crossing} shows that
  at a zero $b$ of $H$ we have $H'(b)>0$ if $b>b_*$, and $H'(b)<0$ if
  $b<b_*$.  Now
  \[
    H(0)=(r+j)\E Y^{2j}>0,
    \qquad
    H(b)=(r+4j)b^{2j}+O(b^{2j-1})
    \qquad(b\to\infty),
  \]
  so $H$ is positive at both ends of $[0,\infty)$.  If $H$ were negative
  somewhere, it would have a first zero, where $H'\le0$, and a later zero,
  where $H'\ge0$; the first therefore lies below $b_*$ and the second above
  it, so the interval where $H$ is negative would contain $b_*$.  It is
  therefore enough to prove
  \begin{equation}\label{eq:pge23-gamma-Hbstar}
    H(b_*)\ge0.
  \end{equation}

  We first write $H(b_*)$ as an integral.  Since
  $F'(b)=-2j\E(Y-b)^{2j-1}$, \eqref{eq:pge23-gamma-H} gives
  \[
    H(b_*)=(r+j)\E(Y-b_*)^{2j}-3jb_*\E(Y-b_*)^{2j-1}
    =(r+j)\,\E\bigl[(Y-b_*)^{2j-1}(Y-b_*-c)\bigr],
  \]
  where $c=\frac{3jb_*}{r+j}$. 
  Expanding the expectation as an integral and substituting $x=y-b_*$,
  \begin{equation}\label{eq:pge23-gamma-H-integral}
    H(b_*) =\frac{r+j}{\Gamma(r) e^{b_*}}\int_{-b_*}^{\infty}
    x^{2j-1}(x-c)(x+b_*)^{r-1}e^{-x}\dd x.
  \end{equation}

  Set
  \begin{equation}\label{eq:pge23-gamma-G}
    G(x)=x^{2j}(x+b_*)^re^{-x},
    \qquad -b_*\le x<\infty,
  \end{equation}
  so that
  \[
    G'(x)=x^{2j-1}(x+b_*)^{r-1}e^{-x}
    \bigl(-x^2+(r+2j-b_*)x+2jb_*\bigr).
  \]
  The quadratic factor here factors further as $(c-x)(x+d)$, where
  \[
    d=\frac{2jb_*}{c}=\frac{2(r+j)}{3}.
  \]
  Thus
  \begin{equation}\label{eq:pge23-gamma-Gprime}
    G'(x)=x^{2j-1}(x+b_*)^{r-1}e^{-x}(c-x)(x+d),
  \end{equation}
  and the integrand of \eqref{eq:pge23-gamma-H-integral} is $-G'(x)/(x+d)$, whence
  \[
    \frac{\Gamma(r)e^{b_*}}{r+j}H(b_*)
    =-\int_{-b_*}^{\infty}\frac{G'(x)}{x+d}\dd x.
  \]
  While $-d$ lies inside the range of integration, the quotient
  $-G'(x)/(x+d)$ is the integrand of \eqref{eq:pge23-gamma-H-integral} and is
  therefore harmless there.

  We integrate by parts on the two sides of $-d$.  Away from $x=-d$,
  \[
    \frac{\dd}{\dd x}\left(-\frac{G(x)-G(-d)}{x+d}\right)
    =-\frac{G'(x)}{x+d}+\frac{G(x)-G(-d)}{(x+d)^2}
  \]
  holds, so for $0<\varepsilon<b_*-d$ we have
  \begin{align*}
    -\int_{-b_*}^{-d-\varepsilon}\frac{G'(x)}{x+d}\dd x
    &=\left[-\frac{G(x)-G(-d)}{x+d}\right]_{-b_*}^{-d-\varepsilon}
    -\int_{-b_*}^{-d-\varepsilon}\frac{G(x)-G(-d)}{(x+d)^2}\dd x,\\
    -\int_{-d+\varepsilon}^{\infty}\frac{G'(x)}{x+d}\dd x
    &=\left[-\frac{G(x)-G(-d)}{x+d}\right]_{-d+\varepsilon}^{\infty}
    -\int_{-d+\varepsilon}^{\infty}\frac{G(x)-G(-d)}{(x+d)^2}\dd x.
  \end{align*}
  Here $G(-b_*)=0$ and $G(x)$ converges to zero as $x\to\infty$, so the
  boundary terms at $-b_*$ and at $\infty$ contribute $G(-d)/(b_*-d)$; the boundary
  terms at $-d-\varepsilon$ and at $-d+\varepsilon$ vanish as
  $\varepsilon\downarrow0$, because $G(-d\pm\varepsilon)-G(-d)
  =O(\varepsilon^2)$ by $G'(-d)=0$; and $(G(x)-G(-d))/(x+d)^2$ extends
  continuously to $x=-d$ for the same reason, combining the two remaining
  integrals.  Hence
  \begin{align}
    -\int_{-b_*}^{\infty}\frac{G'(x)}{x+d}\dd x
    &=\lim_{\varepsilon\downarrow0}
    \left(-\int_{-b_*}^{-d-\varepsilon}\frac{G'(x)}{x+d}\dd x
    -\int_{-d+\varepsilon}^{\infty}\frac{G'(x)}{x+d}\dd x\right)\notag\\
    &=\lim_{\varepsilon\downarrow0}
    \left(\left[-\frac{G(x)-G(-d)}{x+d}\right]_{-b_*}^{-d-\varepsilon}
    +\left[-\frac{G(x)-G(-d)}{x+d}\right]_{-d+\varepsilon}^{\infty}\right.\notag\\
    &\qquad\left.-\int_{-b_*}^{-d-\varepsilon}\frac{G(x)-G(-d)}{(x+d)^2}\dd x
    -\int_{-d+\varepsilon}^{\infty}\frac{G(x)-G(-d)}{(x+d)^2}\dd x\right)\notag\\
    &=\frac{G(-d)}{b_*-d}
    -\int_{-b_*}^{\infty}\frac{G(x)-G(-d)}{(x+d)^2}\dd x.
    \label{eq:pge23-gamma-ibp-max}
  \end{align}
  Here $b_*-d=r(r+j)/(r+4j)>0$ and $G(-d)=d^{2j}(b_*-d)^re^{d}>0$, so the
  first term is positive.  For the second term, $(x+d)^{-2}$ is positive as
  well, so its sign is decided by that of $G(x)-G(-d)$.
  Consequently, if $G(x)-G(-d)\le0$ for every $x\ge-b_*$, that is, if $G(-d)$
  is the global maximum of $G$, then the second term is nonnegative and
  \eqref{eq:pge23-gamma-Hbstar} follows.

  By \eqref{eq:pge23-gamma-Gprime}, the two nonzero local maxima occur at $-d$ and
  $c$.  Their ratio is
  \begin{equation}\label{eq:pge23-gamma-max-ratio}
    \frac{G(c)}{G(-d)}
    =\left(\frac {c}{d}\right)^{2j}
    \left(\frac{b_*+c}{b_*-d}\right)^r e^{-(c+d)}.
  \end{equation}
  Let $t=j/r>0$.  To see that $G$ attains its maximum at $-d$ we prove that
  the logarithm of \eqref{eq:pge23-gamma-max-ratio}, divided by $r$, namely
  \begin{align}
    L(t)={}&2t\log\left(
    \frac{3t(5+8t)}{2(1+t)(1+4t)}\right)
    +\log\left(
    \frac{(1+4t)(5+8t)}{3(1+t)}\right)\notag\\
    &-\frac{32t^2+25t+2}{3(1+4t)},
    \label{eq:pge23-gamma-L}
  \end{align}
  is negative for every $t>0$.  Direct differentiation gives, with
  $A(t)=\frac{3t(5+8t)}{2(1+t)(1+4t)}$,
  \begin{align}
    L'(t)
    &=2\log A(t)
    -\frac{(32t^2+16t-7)(32t^2+25t+2)}
    {3(t+1)(4t+1)^2(8t+5)},\label{eq:pge23-gamma-Lprime}\\
    L''(t)
    &=-\frac{(32t^2+25t+2)
    (128t^4+224t^3+120t^2-28t-25)}
    {t(t+1)^2(4t+1)^3(8t+5)^2}.
    \label{eq:pge23-gamma-Lsecond}
  \end{align}
  Every factor of \eqref{eq:pge23-gamma-Lsecond} other than the quartic
  $P(t)=128t^4+224t^3+120t^2-28t-25$ is positive for $t>0$, so $L''$ and $P$
  have opposite signs.  The coefficients of $P$ change sign once, so by
  Descartes' rule of signs $P$ has at most one positive root.
  Since $P(0) = -25 < 0$ and $P(t)\to\infty$, it has exactly one positive root
  $t_0$, and $L'$ increases on $(0,t_0)$ and decreases on $(t_0,\infty)$.

  We evaluate $L'$ at three places to see its sign changes.
  As $t\downarrow0$ the argument $A(t)$ tends to $0$ while the rational term of \eqref{eq:pge23-gamma-Lprime} tends to
  $-14/15$, so $L'(t)\to-\infty$.  At $t=1$ we have $A(1)=39/20$ and the
  rational term equals $2419/1950$, so the bound
  \[
    \log x\ge\frac{2(x-1)}{x+1}\qquad(x\ge1)
  \]
  gives $2\log A(1)\ge\frac{76}{59}$ and
  \[
    L'(1)\ge\frac{76}{59}-\frac{2419}{1950} \approxeq 0.047622... >0.
  \]
  At $t=3/2$ we have $A(3/2)=153/70$ and the rational
  term equals $19847/12495$; concavity of the logarithm at $x=2$ gives
  \[
    \log\frac{153}{70}\le\log2+\frac{1}{2}\left(\frac{153}{70}-2\right)
    =\log2+\frac{13}{140},
  \]
  so with $2 \log 2 < 7/5$,
  \[
    L'\left(\frac{3}{2}\right)
    <\frac{7}{5} + \frac{13}{70}-\frac{19847}{12495}\approxeq -0.002681... <0.
  \]

  Since $L'$ increases and then decreases, and is negative near $0$, positive at $t=1$ and negative at $t=3/2$, it has two roots: 
  at some $t_1\in(0,1)$ with $L''(t_1) > 0$, and at some $t_2\in(1,3/2)$ with $L''(t_2) < 0$. 
  Thus $L$ decreases on $(0,t_1)$, increases on $(t_1,t_2)$ and decreases on
  $(t_2,\infty)$, so its supremum over $(0,\infty)$ is the larger of $L(t_2)$
  and the limit of $L$ at $0$.  The latter is negative,
  \[
    \lim_{t\downarrow0}L(t)=\log\frac{5}{3}-\frac{2}{3}<0,
  \]
  and $t_2$ lies in $[1,3/2]$, so it remains to control $L$ on $[1,3/2]$.  Let
  \[
    B(t)=\frac{(1+4t)(5+8t)}{3(1+t)}
  \]
  so that $L(t)=2t\log A(t)+\log B(t)-(32t^2+25t+2)/(3(4t+1))$.
  By concavity of $\log$,
  \[
    \log A\le\log2+\frac{A-2}{2},
    \qquad
    \log B\le\log13+\frac{B-13}{13}.
  \]
  Using $\log2<7/10$ and $\log13<13/5$, exact simplification gives
  \begin{equation}\label{eq:pge23-gamma-L-rational}
    L(t)<
    \frac{3(96t^3-191t^2-16t+46)}
    {130(t+1)(4t+1)}.
  \end{equation}
  On $[1,3/2]$ we have $96t-191<0$ and $t^2\ge1$, so the cubic numerator is at
  most $(96t-191)-16t+46=80t-145\le-25$, and $L(t)<0$ there as well.

  To summarize the proof: the signs of $L'$ and $L''$ in
  \eqref{eq:pge23-gamma-Lprime} and \eqref{eq:pge23-gamma-Lsecond} leave two candidates for
  the supremum of $L$, namely its limit at $0$, which is negative, and a single
  local maximum, which lies in $[1,3/2]$.  Confined to that interval, $L$ is
  bounded by the rational function \eqref{eq:pge23-gamma-L-rational}, which is
  negative there; hence $L<0$ on all of $(0,\infty)$.  The ratio
  \eqref{eq:pge23-gamma-max-ratio} is therefore smaller than $1$, so $G$ attains its
  maximum at $-d$ and the integrand $G(x)-G(-d)$ in \eqref{eq:pge23-gamma-ibp-max} is
  nonpositive.  The right side of \eqref{eq:pge23-gamma-ibp-max} is then positive,
  giving $H(b_*)>0$, and the crossing argument yields $H(b)\ge0$ for every
  $b>0$, which is \eqref{eq:pge23-gamma-H} and hence implies \eqref{eq:pge23-gamma-moment2}.
\end{proof}

\subsection{Completion of the proof for \texorpdfstring{$p\ge2/3$}{p>=2/3}}

\begin{theorem}[The theorem for $p\ge2/3$]\label{thm:pge23-main}
  If $p\ge2/3$, then \eqref{eq:main-goal} holds for every odd $m\ge3$.
\end{theorem}

\begin{proof}
  Equivalently, $q\le1/3$.
  \Cref{lem:pge23-two-sided-identity}, the bound of $t(C_m,U)$ by the path density
  $x_{m-1}$, and the definition of $\Phi_m$ give
  \begin{align*}
    t(C_m,W)-\bigl(p^m-pq^{m-1}\bigr)
    &=q^{m-1}+S_m-t(C_m,U)\\
    &\ge q^{m-1}+S_m-x_{m-1}=\Phi_m,
  \end{align*}
  so the theorem follows from the nonnegativity of $\Phi_m$.  By the expansion
  of \Cref{prop:pge23-expansion}, it is enough to show that
  $P_{m,r}\ge0$ on $[-1/2,1/2]^r$; and by \Cref{lem:pge23-dirichlet}, which
  presents $P_{m,r}$ as an expectation of the diagonal values
  $\widetilde P_{m,r}(q,\ell)$ at points $\ell$ of $[-1/2,1/2]$, it is enough
  to prove $\widetilde P_{m,r}(q,\ell)\ge0$ for $1\le r\le(m-1)/2$ and
  $-1/2\le\ell\le1/2$.

  When $\ell\le0$, \Cref{lem:pge23-ell-negative} proves
  $\widetilde P_{m,r}(q,\ell)\ge0$.  When $\ell>0$, the beta formula
  \Cref{lem:pge23-beta} and the Laplace--gamma representation
  \Cref{lem:pge23-laplace-gamma} write $\widetilde P_{m,r}(q,\ell)$ as a
  positive weight times an average of the expectations
  $\E\rho_{n,m}(q+xY)$, each of which is nonnegative by
  \Cref{prop:pge23-gamma-positive}.  Hence $\widetilde P_{m,r}(q,\ell)\ge0$
  throughout that range, so $\Phi_m\ge0$.
\end{proof}

\section{Below the threshold \texorpdfstring{$1/2<p<2/3$}{1/2<p<2/3}: the single-outlier reduction for \texorpdfstring{$m\ge9$}{m>=9}}
\label{sec:plt23}

\tk{TODO (structure): make this section readable independently of
\Cref{sec:pge23}, so that a reader interested in one region need not
read the other.  Remaining dependencies: the generating function $L_-$ and the
block determinant identity \eqref{eq:block-determinant-identity}.}

From now on, we assume
\begin{equation}\label{eq:plt23-range}
  \frac{1}{2}<p<\frac{2}{3}, \qquad \frac{1}{3}<q<\frac{1}{2},
\end{equation}
and $m\ge9$ is odd. 
We continue with the block decomposition
\eqref{eq:block-U} and the operator norm bound of \Cref{lem:half-norm}.  
We will also use the following bound on the spectrum of $A$, in which
$\Tr(A^2)=\sum_i\lambda_i^2$.

\begin{lemma}[Hilbert--Schmidt bound]\label{lem:plt23-HS-bound}
  \begin{equation}\label{eq:plt23-HS-bound}
    \Tr(A^2)+2\norm g^2\le pq.
  \end{equation}
\end{lemma}

\begin{proof}
  Since $0\le U\le1$,
  \[
    \Tr(T_U^2)=\int_{[0,1]^2}U(x,y)^2\dd x\dd y
    \le\int_{[0,1]^2}U(x,y)\dd x\dd y=q.
  \]
  Expanding $T_U^2$ in the block decomposition \eqref{eq:block-U},
  \[
    \Tr(T_U^2)=\Tr\begin{pmatrix}q^2 + \norm g^2 & q g^* + g^* A \\ qg + Ag & g g^* + A^2\end{pmatrix}=q^2+2\norm g^2+\Tr(A^2).
  \]
  Subtracting $q^2$ from the both sides proves the claim.
\end{proof}

\subsection{A one-sided Schur-complement identity}\label{subsec:plt23-one-sided-schur}

In \Cref{subsec:pge23-two-sided-schur} we introduced the generating functions $L_-$ and
$L_+$, attached to $W$ and to $U$.  Only $L_-$ is used here.

By \eqref{eq:logdet-cycle},
\[
  t(C_m,W)=m\,[z^m]\bigl(-\log\det(I-zT_W)\bigr),
\]
and the block decomposition \eqref{eq:block-U} gives
\[
  I-zT_W=\begin{pmatrix}1-pz&zg^*\\zg&I+zA\end{pmatrix}.
\]
Applying the block determinant identity
\eqref{eq:block-determinant-identity} with $D=I+zA$ and $v=zg$ factorizes
this determinant into three parts,
\begin{equation}\label{eq:plt23-one-sided-factorization}
  \det(I-zT_W)=\det(I+zA)\,(1-pz)\bigl(1-Y(z)\bigr),
  \qquad
  Y(z)=\frac{z^2\ip g{(I+zA)^{-1}g}}{1-pz},
\end{equation}
the last factor being the Schur complement, divided by $1-pz$.

\begin{lemma}[One-sided identity]\label{lem:plt23-one-sided-identity}
  For every odd $m\ge3$,
  \begin{equation}\label{eq:plt23-one-sided-identity}
    t(C_m,W)=p^m-\Tr(A^m)+m[z^m]L_-(z),
  \end{equation}
  where
  \begin{equation}\label{eq:plt23-Lminus-def}
    L_-(z)=-\log\bigl(1-Y(z)\bigr).
  \end{equation}
  Moreover, every coefficient of $L_-(z)$ is nonnegative.
\end{lemma}

\begin{proof}
  Taking $-\log$ in \eqref{eq:plt23-one-sided-factorization} turns the product into a
  sum of three terms,
  \[
    -\log\det(I-zT_W)=-\log\det(I+zA)-\log(1-pz)+L_-(z),
  \]
  and for odd $m$, \Cref{lem:series-expansions}\,(2) gives
  \[
    m\,[z^m]\bigl(-\log\det(I+zA)\bigr)=\Tr\bigl((-A)^m\bigr)=-\Tr(A^m),
    \qquad
    m\,[z^m]\bigl(-\log(1-pz)\bigr)=p^m.
  \]
  Together with the third summand this is \eqref{eq:plt23-one-sided-identity}.

  It remains to check nonnegativity of coefficients.  We first show that the
  coefficients of $Y$ are nonnegative, which implies the same for
  $L_-=-\log(1-Y)=\sum_{k\ge1}Y^k/k$.

  Expanding the two factors of $Y$,
  \[
    \ip g{(I+zA)^{-1}g}=\sum_{j\ge0}z^j\int(-\lambda)^j\dd\mu(\lambda),
    \qquad
    \frac{1}{1-pz}=\sum_{k\ge0}p^kz^k,
  \]
  where the first follows from \eqref{eq:resolvents-common} with $A$ replaced
  by $-A$ and $\mu=\sum_i\norm{P_ig}^2\delta_{\lambda_i}$ is the measure
  \eqref{eq:mu-def}, gives the coefficients of $Y$ as
  \[
    [z^r]Y(z)
    =\int\sum_{j+k=r-2}p^k(-\lambda)^j\dd\mu(\lambda)
    \qquad(r\ge2).
  \]
  Summing the geometric series inside the integral,
  \begin{equation}\label{eq:plt23-Y-coeff}
    [z^r]Y(z)
    =\int\frac{p^{r-1}-(-\lambda)^{r-1}}{p+\lambda}\dd\mu(\lambda).
  \end{equation}
  By \Cref{lem:half-norm}, the support of $\mu$ lies in $[-1/2,1/2]$, so
  $|\lambda|<p$ on the whole range of integration and both the numerator and
  the denominator on the right of \eqref{eq:plt23-Y-coeff} are positive.  The
  coefficients of $Y$ are therefore nonnegative.
\end{proof}


\subsection{At most one eigenvalue above \texorpdfstring{$q$}{q}}

\begin{lemma}[Eigenvalues above $q$]\label{lem:plt23-eigenvalues-above-q}
  There is at most one eigenvalue of $A$ that is larger than $q$.  If none of
  them is larger than $q$, then \eqref{eq:main-goal} holds.
\end{lemma}

\begin{proof}
  If two eigenvalues of $A$ were larger than $q$, then
  \[
    \Tr(A^2)=\sum_i\lambda_i^2>2q^2>pq,
  \]
  the last inequality because $2q>2/3>p$, which contradicts
  \Cref{lem:plt23-HS-bound}.  This proves the first claim.

  Assume now that no eigenvalue of $A$ is larger than $q$.  Then every
  eigenvalue $\lambda_i$ of $A$ satisfies
  \[
    \lambda_i^m\le q^{m-2}\lambda_i^2,
  \]
  which is immediate for $0\le\lambda_i\le q$, while for $\lambda_i<0$ the left side is negative and the right side is positive.
  Therefore,
  \[
    \Tr(A^m) = \sum_{i}\lambda_i^m\le q^{m-2} \sum_{i} \lambda_i^2 = q^{m-2} \Tr(A^2) \le q^{m-2} \bigl(\Tr(A^2) + 2\norm g^2\bigr) \le pq^{m-1},
  \]
  where we used \Cref{lem:plt23-HS-bound} in the last inequality.
  Since the coefficients of $L_-$ are nonnegative, \eqref{eq:plt23-one-sided-identity} gives
  \[
    t(C_m,W) = p^m - \Tr(A^m) + m[z^m]L_-(z) \ge p^m - pq^{m-1},
  \]
  proving \eqref{eq:main-goal} in this case.
\end{proof}

We henceforth assume that $A$ has exactly one eigenvalue
\begin{equation}\label{eq:plt23-Lambda-above-q}
  \Lambda>q.
\end{equation}
Define
\begin{equation}\label{eq:plt23-M-def}
  M=\sqrt{pq-\Lambda^2},
\end{equation}
which by \eqref{eq:plt23-HS-bound} controls the remaining eigenvalues:
\[
  \sum_{\lambda_i\ne\Lambda}\lambda_i^2\le pq-\Lambda^2=M^2,
\]
so in particular $|\lambda_i|\le M$ for every $\lambda_i\ne\Lambda$.

These quantities are ordered as
\begin{equation}\label{eq:plt23-M-Lambda-order}
  0\le M<q<\Lambda<\frac{1}{2}<p,
\end{equation}
where $M<q$ because $M^2<pq-q^2=q(1-2q)<q^2$ for $q>1/3$, and
$\Lambda<1/2$ comes from \Cref{lem:half-norm}.

\subsection{A lower bound via splitting \texorpdfstring{$g$}{g} along the
eigenfunction of \texorpdfstring{$\Lambda$}{Lambda}}\label{subsec:plt23-splitting-lower-bound}

Let $\phi$ be a unit eigenfunction of $A$ for $\Lambda$, which is determined up to sign. 
We will choose the sign of $\phi$ later.  
Split $g$ into its component along $\phi$ and the rest,
\begin{equation}\label{eq:plt23-g-decomp}
  g=\gamma\phi+g_s,
  \qquad \gamma=\ip g\phi,
  \qquad g_s\perp\phi,
\end{equation}
so that $\norm g^2=\gamma^2+\norm{g_s}^2$.

\begin{proposition}[Lower bound via splitting $g$ along the eigenfunction of $\Lambda$]\label{prop:plt23-splitting-lower-bound}
  For odd $m$, let
  \begin{align}
    k_m(\lambda)&=\frac{p^{m-1}-\lambda^{m-1}}{p+\lambda},
    \label{eq:plt23-km-def}\\
    \Lambda_m&=2M^{m-2}+mk_m(\Lambda),\label{eq:plt23-Lambdam-def}\\
    M_m&=2M^{m-2}+mk_m(M),\label{eq:plt23-Mm-def}\\
    R_m&=\Lambda^m+M^m-pq^{m-1}.\label{eq:plt23-Rm-def}
  \end{align}
  Then
  \begin{align}
    t(C_m,W)-\bigl(p^m-pq^{m-1}\bigr)
    &= pq^{m-1} - \Tr(A^m)+m[z^m]L_-(z) \nonumber
    \\
    {} &\ge -R_m+\Lambda_m\gamma^2+M_m\norm{g_s}^2.\label{eq:plt23-splitting-lower-bound}
  \end{align}
  Moreover,
  \begin{equation}\label{eq:plt23-Lambdam-Mm-positive}
    \Lambda_m>0,
    \qquad
    M_m>0.
  \end{equation}
\end{proposition}

\begin{proof}
  We first bound $m[z^m]L_-(z)$ from below.  The coefficients of $Y$ are
  nonnegative, hence so are those of $L_-(z)-Y(z)=\sum_{k\ge2}Y(z)^k/k$.
  Therefore
  \begin{align}
    m[z^m]L_-(z)&\ge m[z^m]Y(z)
    =m\int\frac{p^{m-1}-(-\lambda)^{m-1}}{p+\lambda}\dd\mu(\lambda)\notag\\
    &=m\int k_m(\lambda)\dd\mu(\lambda),
    \label{eq:plt23-Lminus-linear-lower}
  \end{align}
  where the first equality follows from \eqref{eq:plt23-Y-coeff} and the second equality holds because $m-1$ is even.

  The decomposition \eqref{eq:plt23-g-decomp} can be expressed in terms of the spectral measure $\mu$ of $A$ as
  \[
    \mu(\{\Lambda\})=\gamma^2,
    \qquad
    \int_{[-M,M]}\dd\mu(\lambda)=\norm{g_s}^2,
  \]
  so \eqref{eq:plt23-Lminus-linear-lower} can be decomposed as
  \[
    m\int k_m(\lambda)\dd\mu(\lambda)
    =mk_m(\Lambda)\gamma^2
    +m\int_{[-M,M]}k_m(\lambda)\dd\mu(\lambda).
  \]
  The remaining integral can be bounded from below by $k_m(M)\norm{g_s}^2$, because $k_m$ attains its minimum over $[-M,M]$ at $\lambda=M$. 
  To see this, for $0\le s<p$, we first have
  \[
    k_m(-s)=\frac{p^{m-1}-s^{m-1}}{p-s}
    >\frac{p^{m-1}-s^{m-1}}{p+s}=k_m(s),
  \]
  so it suffices to minimize over $0\le\lambda\le M$. 
  In this range, $k_m$ is strictly decreasing: with $n=m-1\ge2$,
  \[
    k_m'(\lambda)
    =-\frac{p^n+np\lambda^{n-1}+(n-1)\lambda^n}{(p+\lambda)^2}<0,
  \]
  since every term of the numerator is nonnegative for $\lambda\ge0$ and
  $p^n>0$.
  Hence
  \begin{equation}\label{eq:plt23-km-min-at-M}
    k_m(\lambda)\ge k_m(M)
    \qquad\text{whenever }|\lambda|\le M,
  \end{equation}
  and altogether
  \begin{equation}\label{eq:plt23-Lminus-lower}
    m[z^m]L_-(z)
    \ge mk_m(\Lambda)\gamma^2+mk_m(M)\norm{g_s}^2.
  \end{equation}

  We next bound $\Tr(A^m)$.  Splitting off the eigenvalue $\Lambda$,
  \[
    \Tr(A^m)=\Lambda^m+\sum_{\lambda_i\ne\Lambda}\lambda_i^m,
  \]
  where $|\lambda_i|\le M$ gives
  \[
    \sum_{\lambda_i\ne\Lambda}\lambda_i^m
    \le M^{m-2}\sum_{\lambda_i\ne\Lambda}\lambda_i^2.
  \]
  From \eqref{eq:plt23-HS-bound},
  \[
    \Lambda^2+\sum_{\lambda_i\ne\Lambda}\lambda_i^2+2\norm g^2
    =\sum_i\lambda_i^2+2\norm g^2\le pq,
  \]
  so $M^2=pq-\Lambda^2$ gives
  $\sum_{\lambda_i\ne\Lambda}\lambda_i^2\le M^2-2\norm g^2$.  Summing up,
  \begin{align}
    -\Tr(A^m)
    &\ge-\Lambda^m-M^{m-2}\sum_{\lambda_i\ne\Lambda}\lambda_i^2\notag\\
    &\ge-\Lambda^m-M^{m-2}\bigl(M^2-2\norm g^2\bigr)\notag\\
    &=-\Lambda^m-M^m+2M^{m-2}\bigl(\gamma^2+\norm{g_s}^2\bigr),
    \label{eq:plt23-trace-upper}
  \end{align}
  where the last equality follows from \eqref{eq:plt23-g-decomp}.
  
  It remains to combine the two bounds.  
  Subtracting $p^m-pq^{m-1}$ in \eqref{eq:plt23-one-sided-identity} gives the first line of \eqref{eq:plt23-splitting-lower-bound},
  \[
    t(C_m,W)-\bigl(p^m-pq^{m-1}\bigr)
    =pq^{m-1}-\Tr(A^m)+m[z^m]L_-(z),
  \]
  and substituting \eqref{eq:plt23-trace-upper} and \eqref{eq:plt23-Lminus-lower},
  \begin{align*}
    &pq^{m-1}-\Tr(A^m)+m[z^m]L_-(z)
    \\
    \ge{}&pq^{m-1}-\Lambda^m-M^m
    +2M^{m-2}\bigl(\gamma^2+\norm{g_s}^2\bigr) +mk_m(\Lambda)\gamma^2+mk_m(M)\norm{g_s}^2\\
    ={}&-\underbrace{(\Lambda^m+M^m-pq^{m-1})}_{\eqqcolon R_m}
    +\underbrace{\bigl(2M^{m-2}+mk_m(\Lambda)\bigr)}_{\eqqcolon \Lambda_m}\gamma^2
    +\underbrace{\bigl(2M^{m-2}+mk_m(M)\bigr)}_{\eqqcolon M_m}\norm{g_s}^2,
  \end{align*}
  which proves \eqref{eq:plt23-splitting-lower-bound}.

  Finally, since $k_m(p)=0$ and $k_m$ is strictly decreasing on $[0,p]$, $k_m$ is positive on $[0,p)$. 
  From \eqref{eq:plt23-M-Lambda-order}, we have $M, \Lambda \in [0, p)$, so both $k_m(\Lambda)$ and $k_m(M)$ are positive.
  Since the other terms in the definitions of $\Lambda_m$ and $M_m$ are also positive, \eqref{eq:plt23-Lambdam-Mm-positive} holds.
\end{proof}

\subsection{Minimizing
\texorpdfstring{$\Lambda_m\gamma^2+M_m\norm{g_s}^2$}{Lambda\_m gamma\^{}2 + M\_m
|g\_s|\^{}2}}
\label{subsec:plt23-minimization}

For fixed $q$ and $\Lambda$, let $S$ be the set of pairs
$(\gamma,\norm{g_s})$ that are realizable by some graphon satisfying the assumptions we made so far. 
By deriving inequalities on $\gamma$ and $\norm{g_s}$, we construct a larger
set $\overline S\supseteq S$, over which the minimum can be computed, and
\begin{equation}\label{eq:plt23-overapprox}
  \Lambda_m\gamma^2+M_m\norm{g_s}^2
  \ge\min_{(\gamma,\norm{g_s})\in\overline S}
  \bigl(\Lambda_m\gamma^2+M_m\norm{g_s}^2\bigr).
\end{equation}
Specifically, \Cref{lem:plt23-gamma-upper-bound} bounds $\gamma$ from above and
\Cref{lem:plt23-gs-lower-bound} bounds $\norm{g_s}$ from below.  Both are linear in
$\gamma$ and $\norm{g_s}$, and $\Lambda_m\gamma^2+M_m\norm{g_s}^2$ is a
convex quadratic in them by \eqref{eq:plt23-Lambdam-Mm-positive}, so the minimization in
\eqref{eq:plt23-overapprox} is convex.

Both lemmas bound $\gamma$ and $\norm{g_s}$ through the decomposition of
$|\phi|$ along $\one$ and $\phi$,
\begin{equation}\label{eq:plt23-absphi-decomp}
  |\phi|=a_\phi\one+b\phi+\phi_\perp,
  \qquad
  a_\phi=\int|\phi|,
  \qquad
  b=\ip{|\phi|}\phi,
\end{equation}
where $\phi_\perp$ is orthogonal to $\one$ and $\phi$.  We fix the sign of
$\phi$, left open in \Cref{subsec:plt23-splitting-lower-bound}, so that $b\ge0$.
Note that
\begin{equation}\label{eq:plt23-absphi-normalization}
  a_\phi^2+b^2+\norm{\phi_\perp}^2=\norm{|\phi|}^2 = \norm{\phi}^2 = 1.
\end{equation}
Besides this identity, we will need a lower bound on $a_\phi$.  Let $\phi_+$
and $\phi_-$ be the positive and negative parts of $\phi$, so that
$\phi=\phi_+-\phi_-$ and $|\phi|=\phi_++\phi_-$.  Since $\phi\perp\one$, we
have $\int\phi=0$, and therefore
\begin{equation}\label{eq:plt23-phi-parts}
  \int\phi_+=\int\phi_-=\frac{a_\phi}{2}.
\end{equation}
Expanding $\ip\phi{T_U\phi}$ in these two parts and dropping the cross term,
which is nonnegative because $U\ge0$,
\begin{align*}
  \Lambda=\ip\phi{T_U\phi}
  &=\iint U(x,y)\bigl(\phi_+(x)\phi_+(y)+\phi_-(x)\phi_-(y)\bigr)\dd x\dd y\\
  &\qquad-2\iint U(x,y)\phi_+(x)\phi_-(y)\dd x\dd y\\
  &\le\iint U(x,y)
  \bigl(\phi_+(x)\phi_+(y)+\phi_-(x)\phi_-(y)\bigr)\dd x\dd y\\
  &\le\left(\int\phi_+\right)^2+\left(\int\phi_-\right)^2
  =\frac{a_\phi^2}{2},
\end{align*}
the last inequality because $U\le1$ and the last equality by
\eqref{eq:plt23-phi-parts}.  Thus
\begin{equation}\label{eq:plt23-aphi-lower}
  a_\phi^2\ge2\Lambda.
\end{equation}

\begin{lemma}[Upper bound for $\gamma$]\label{lem:plt23-gamma-upper-bound}
  \begin{equation}\label{eq:plt23-gamma-upper-bound}
    \gamma\le u_\gamma:=
    \frac{a_\phi^2-2\Lambda-2\Lambda b}{2a_\phi}.
  \end{equation}
\end{lemma}

\begin{proof}
  For every $x$, the following holds
  \[
    (T_U \phi)(x) = \int U(x, y) \phi(y) \dd y \le \int U(x, y) \phi_+(y) \dd y \le \int \phi_+(y) \dd y = \frac{a_\phi}{2},
  \]
  where we used $U \ge 0$ and $\phi \le \phi_+$ in the first inequality, and $U \le 1$ in the second.
  Taking the inner product of both sides with $\phi_+\ge0$ gives
  \begin{equation}\label{eq:plt23-phiplus-pairing}
    \ip{\phi_+}{T_U\phi}
    \le\frac{a_\phi}{2}\int\phi_+
    =\frac{a_\phi^2}{4}.
  \end{equation}
  To expand the left side, the block form \eqref{eq:block-U} gives
  \[
    T_U\phi=\begin{pmatrix}
      q & g^* \\
      g & A
    \end{pmatrix} \begin{pmatrix}
      0 \\ \phi
    \end{pmatrix}=\ip g\phi\,\one+A\phi=\gamma\one+\Lambda\phi,
  \]
  using $\gamma=\ip g\phi$ from \eqref{eq:plt23-g-decomp} and $A\phi=\Lambda\phi$.
  Hence, by \eqref{eq:plt23-phi-parts},
  \[
    \ip{\phi_+}{T_U\phi}
    =\gamma\ip{\phi_+}\one+\Lambda\ip{\phi_+}\phi
    =\frac{\gamma a_\phi}{2}+\Lambda\ip{\phi_+}\phi.
  \]
  Finally, from $\phi_+(x) \phi_-(x) = 0$, we have $\ip{\phi_+}\phi=\ip{\phi_+}{\phi_+} - \ip{\phi_+}{\phi_-} = \norm{\phi_+}^2$, and
  \[
    \norm{\phi_+}^2+\norm{\phi_-}^2=\norm\phi^2=1,
    \qquad
    \norm{\phi_+}^2-\norm{\phi_-}^2=\ip{|\phi|}\phi=b,
  \]
  so $\ip{\phi_+}\phi=(1+b)/2$.
  Substituting into \eqref{eq:plt23-phiplus-pairing} gives \eqref{eq:plt23-gamma-upper-bound}.
\end{proof}

\begin{lemma}[Lower bound for $\norm{g_s}$]\label{lem:plt23-gs-lower-bound}
  The following inequality holds:
  \begin{equation}\label{eq:plt23-gs-inner-lower}
    b\gamma+\ip{g_s}{\phi_\perp}
    \ge\frac{(\Lambda-q)a_\phi^2+(\Lambda-M)\norm{\phi_\perp}^2}{2a_\phi}
    \eqqcolon\cH.
  \end{equation}
  Consequently,
  \begin{equation}\label{eq:plt23-gs-norm-lower}
    \norm{g_s}\norm{\phi_\perp}\ge \ip{g_s}{\phi_\perp}\ge\cH-b\gamma.
  \end{equation}
\end{lemma}

\begin{proof}
  From $|\phi|=\phi_++\phi_-$ and $\phi=\phi_+-\phi_-$, we have
  \begin{align*}
    \ip{|\phi|}{T_U|\phi|}
    &=\ip{\phi_+}{T_U\phi_+}+\ip{\phi_-}{T_U\phi_-}
    +\ip{\phi_+}{T_U\phi_-}+\ip{\phi_-}{T_U\phi_+},\\
    \ip\phi{T_U\phi}
    &=\ip{\phi_+}{T_U\phi_+}+\ip{\phi_-}{T_U\phi_-}
    -\ip{\phi_+}{T_U\phi_-}-\ip{\phi_-}{T_U\phi_+},
  \end{align*}
  so the two inner products differ only in the sign of their cross terms. 
  Since $T_U$ is self-adjoint, $\ip{\phi_-}{T_U\phi_+}=\ip{\phi_+}{T_U\phi_-}$, 
  and therefore
  \begin{align*}
    \ip{|\phi|}{T_U|\phi|}-\ip\phi{T_U\phi}
    &=2\ip{\phi_+}{T_U\phi_-}+2\ip{\phi_-}{T_U\phi_+}
    \\
    &=4\ip{\phi_+}{T_U\phi_-} 
    \\
    &= 4 \iint \phi_+(x) U(x, y) \phi_-(y) \dd x \dd y
    \\
    &\ge0,
  \end{align*}
  where the inequality holds because $U\ge0$ and $\phi_\pm\ge0$.  
  We also have
  \[
    \ip{\phi}{T_U \phi} = \ip{\phi}{\gamma \one + \Lambda \phi} = \Lambda
  \]
  since $T_U\phi=\gamma\one+\Lambda\phi$ and $\phi\perp\one$. 
  Hence
  \begin{equation}\label{eq:plt23-absphi-lower}
    \ip{|\phi|}{T_U|\phi|}\ge\ip\phi{T_U\phi}=\Lambda.
  \end{equation}

  We expand the left side of \eqref{eq:plt23-absphi-lower}.  In
  $|\phi|=a_\phi\one+b\phi+\phi_\perp$ the last two terms are orthogonal to
  $\one$, so the block form \eqref{eq:block-U} gives
  \[
    T_U|\phi|=\begin{pmatrix}
      q & g^* \\
      g & A
    \end{pmatrix} \begin{pmatrix}
      a_\phi \\ b\phi+\phi_\perp
    \end{pmatrix}
    =\bigl(a_\phi g+A(b\phi+\phi_\perp)\bigr)+\bigl(qa_\phi+\ip g{b\phi+\phi_\perp}\bigr)\one
    .
  \]
  Taking the inner product of this with $|\phi|$, and using
  $\ip{|\phi|}\one=a_\phi$, $\one\perp g,\phi,\phi_\perp$, and that $A$ maps
  into $\one^\perp$,
  \begin{align*}
    &\ip{|\phi|}{T_U|\phi|}
    \\
    {} = &\ip{a_\phi \one}{a_\phi g + A(b\phi + \phi_\perp)}+ \ip{b\phi+\phi_\perp}{a_\phi g+A(b\phi+\phi_\perp)}+\bigl(qa_\phi+\ip g{b\phi+\phi_\perp}\bigr) \ip{|\phi|}\one
    \\
    {} =&\ip{b\phi+\phi_\perp}{a_\phi g+A(b\phi+\phi_\perp)}+\bigl(qa_\phi+\ip g{b\phi+\phi_\perp}\bigr) \ip{|\phi|}\one
    \\
    {}=&\ip{b\phi+\phi_\perp}{a_\phi g+A(b\phi+\phi_\perp)}+a_\phi\bigl(qa_\phi+\ip g{b\phi+\phi_\perp}\bigr)
    \\
    {} =&qa_\phi^2+2a_\phi\ip g{b\phi+\phi_\perp}
    +\ip{b\phi+\phi_\perp}{A(b\phi+\phi_\perp)}.
  \end{align*}
  The two remaining terms expand by $g=\gamma\phi+g_s$, $A\phi=\Lambda\phi$
  and $\phi_\perp\perp\phi$:
  \begin{align*}
    \ip g{b\phi+\phi_\perp}
    &=b\ip g\phi+\ip{\gamma\phi+g_s}{\phi_\perp}
    =b\gamma+\ip{g_s}{\phi_\perp},\\
    \ip{b\phi+\phi_\perp}{A(b\phi+\phi_\perp)}
    &=b^2\ip\phi{A\phi}+2b\Lambda\ip{\phi_\perp}\phi
    +\ip{\phi_\perp}{A\phi_\perp}
    =\Lambda b^2+\ip{\phi_\perp}{A\phi_\perp},
  \end{align*}
  which gives 
  \[
    \ip{|\phi|}{T_U |\phi|} = q a_\phi^2 + 2 a_\phi\bigl(b \gamma + \ip{g_s}{\phi_\perp}\bigr) + \Lambda b^2 + \ip{\phi_\perp}{A \phi_\perp}.
  \]
  Every eigenvalue of $A$ other than $\Lambda$ is at most $M$ in absolute value, 
  so $\ip{\phi_\perp}{A\phi_\perp}\le M\norm{\phi_\perp}^2$, and therefore
  \[
    \Lambda\le qa_\phi^2+2a_\phi\bigl(b\gamma+\ip{g_s}{\phi_\perp}\bigr)
    +\Lambda b^2+M\norm{\phi_\perp}^2.
  \]
  By \eqref{eq:plt23-absphi-normalization},
  $\Lambda-\Lambda b^2=\Lambda\bigl(a_\phi^2+\norm{\phi_\perp}^2\bigr)$, 
  so rearranging gives
  \[
    (\Lambda-q)a_\phi^2+(\Lambda-M)\norm{\phi_\perp}^2
    \le2a_\phi\bigl(b\gamma+\ip{g_s}{\phi_\perp}\bigr),
  \]
  which is \eqref{eq:plt23-gs-inner-lower}. 
  The consequence follows from \eqref{eq:plt23-gs-inner-lower} and Cauchy--Schwarz,
  $\norm{g_s}\norm{\phi_\perp}\ge\ip{g_s}{\phi_\perp}\ge\cH-b\gamma$.
\end{proof}

By \eqref{eq:plt23-gamma-upper-bound} and \eqref{eq:plt23-gs-norm-lower}, the pair
$(\gamma,\norm{g_s})$ lies in
\begin{equation}\label{eq:plt23-Sbar}
  \overline S=\bigl\{(t,s):
  t\le u_\gamma,\quad
  bt+s\norm{\phi_\perp}\ge\cH,\quad
  s\ge0\bigr\},
\end{equation}
a polyhedron determined by $a_\phi$, $b$ and $\norm{\phi_\perp}$.  Hence
\begin{equation}\label{eq:plt23-min-over-Sbar}
  \Lambda_m\gamma^2+M_m\norm{g_s}^2
  \ge\min_{(t,s)\in\overline S}
  \bigl(\Lambda_mt^2+M_ms^2\bigr),
\end{equation}
and the right side is the minimum of a convex quadratic over a polyhedron.

That minimum still depends on $\phi$, through $a_\phi$, $b$ and
$\norm{\phi_\perp}$.  The following theorem evaluates it and bounds it from
below by a quantity depending only on $q$, $\Lambda$, $M$ and $m$.

\tk{Revised until here.}

\begin{theorem}[Reduction to $\psi(\xi,\rho)$]\label{thm:plt23-huber-elim}
  Let
  \begin{align}
    \Omega_m&=\frac{M_m(\Lambda-M)\sqrt{2\Lambda}\,(1-2\Lambda)^2}{4\Lambda^2},
    \label{eq:plt23-C-xi-rho}\\
    \xi&=\frac{4\Lambda^2(\Lambda-q)}{(1-2\Lambda)^2},
    \label{eq:plt23-xi-def}\\
    \rho&=\frac{\Lambda_m}{M_m}\frac{\sqrt\Lambda}{2\sqrt2\,(\Lambda-M)},
    \label{eq:plt23-rho-def}\\
    \psi(\xi,\rho)&=
    \min_{0\le v\le1}
    \left\{\rho v^2+\pos{\xi-v+v^2}\right\}.
    \label{eq:plt23-psi-def}
  \end{align}
  Then
  \begin{equation}\label{eq:plt23-huber-payment}
    \Lambda_m\gamma^2+M_m\norm{g_s}^2
    \ge \Omega_m\psi(\xi,\rho).
  \end{equation}
\end{theorem}

\begin{proof}
  We first evaluate the minimum in \eqref{eq:plt23-min-over-Sbar} in the variable $s$.
  The objective increases in $s$, so at a minimizer $s$ is as small as
  $\overline S$ allows, that is, $s\norm{\phi_\perp}=\pos{\cH-bt}$.  For
  $\phi_\perp\ne0$ this gives $s^2=\pos{\cH-bt}^2/\norm{\phi_\perp}^2$, and
  \begin{equation}\label{eq:plt23-one-d-t}
    \min_{(t,s)\in\overline S}\bigl(\Lambda_mt^2+M_ms^2\bigr)
    =\min_{t\le u_\gamma}
    \left\{
    \Lambda_mt^2+\frac{M_m}{\norm{\phi_\perp}^2}\pos{\cH-bt}^2
    \right\};
  \end{equation}
  for $\phi_\perp=0$ the constraint reads $bt\ge\cH$ instead, and that case
  is treated at the end.  We write $\gamma$ for $t$ again below.

  Assume $\phi_\perp\ne0$.  Suppose first that $u_\gamma\le0$.  Every
  admissible $\gamma$ is nonpositive, so $\cH-b\gamma\ge\cH$, and the right
  side of \eqref{eq:plt23-one-d-t} is at least
  \[
    M_m\frac{\cH^2}{\norm{\phi_\perp}^2}.
  \]
  By \eqref{eq:plt23-gs-inner-lower},
  \[
    \frac{\cH^2}{\norm{\phi_\perp}^2}
    =\frac{\bigl((\Lambda-q)a_\phi^2+(\Lambda-M)\norm{\phi_\perp}^2\bigr)^2}
    {4a_\phi^2\norm{\phi_\perp}^2}\ge(\Lambda-q)(\Lambda-M),
  \]
  because this inequality is equivalent to
  $\bigl((\Lambda-q)a_\phi^2-(\Lambda-M)\norm{\phi_\perp}^2\bigr)^2\ge0$.  On
  the other hand, taking $v=0$ in \eqref{eq:plt23-psi-def} gives
  \[
    \Omega_m\psi(\xi,\rho)
    \le \Omega_m\xi=M_m(\Lambda-M)(\Lambda-q)\sqrt{2\Lambda}
    \le M_m(\Lambda-M)(\Lambda-q).
  \]
  Thus \eqref{eq:plt23-huber-payment} holds in this case.

  Now suppose $u_\gamma>0$.  A minimizer in \eqref{eq:plt23-one-d-t} may be taken
  with $\gamma\ge0$: replacing a negative $\gamma$ by $0$ decreases both
  terms.  The bound \eqref{eq:plt23-gamma-upper-bound} gives
  \[
    2\gamma a_\phi+2\Lambda b
    \le a_\phi^2-2\Lambda=1-2\Lambda-b^2-\norm{\phi_\perp}^2
    \le1-2\Lambda-b^2.
  \]
  Since $a_\phi^2\ge2\Lambda$,
  \begin{equation}\label{eq:plt23-b-gamma-budget}
    b^2+2\Lambda b+2\gamma\sqrt{2\Lambda}\le1-2\Lambda.
  \end{equation}
  It follows that
  \begin{equation}\label{eq:plt23-beta-gamma}
    0\le\gamma\le\frac{1-2\Lambda}{2\sqrt{2\Lambda}},
    \qquad
    b\le\beta(\gamma):=
    \sqrt{(1-\Lambda)^2-2\gamma\sqrt{2\Lambda}}-\Lambda.
  \end{equation}
  Rationalizing the square root,
  \begin{equation}\label{eq:plt23-beta-rational}
    \beta(\gamma)
    \le\frac{1-2\Lambda-2\gamma\sqrt{2\Lambda}}{2\Lambda}.
  \end{equation}
  Also, since $2\Lambda\le a_\phi^2\le1$,
  \begin{equation}\label{eq:plt23-H-lower}
    \cH\ge\frac{(\Lambda-q)\sqrt{2\Lambda}}{2}
    +\frac{(\Lambda-M)\norm{\phi_\perp}^2}{2}.
  \end{equation}
  For every real $w$,
  \begin{equation}\label{eq:plt23-K-AMGM}
    \inf_{t>0}\frac{\pos{w+(\Lambda-M)t/2}^2}{t}=2(\Lambda-M)\pos w.
  \end{equation}
  Indeed, for $w\ge0$ this is the arithmetic--geometric mean inequality,
  with equality at $t=2w/(\Lambda-M)$; for $w<0$, both sides vanish after
  choosing $0<t\le-2w/(\Lambda-M)$.  Using \eqref{eq:plt23-one-d-t},
  \eqref{eq:plt23-H-lower}, and \eqref{eq:plt23-K-AMGM}, and then maximizing $b$ by
  \eqref{eq:plt23-beta-gamma}, we obtain
  \begin{equation}\label{eq:plt23-one-d-gamma}
    \Lambda_m\gamma^2+M_m\norm{g_s}^2
    \ge
    \min_{0\le\gamma\le(1-2\Lambda)/(2\sqrt{2\Lambda})}
    \left\{
    \Lambda_m\gamma^2+2M_m(\Lambda-M)
    \pos{\frac{(\Lambda-q)\sqrt{2\Lambda}}{2}-\gamma\beta(\gamma)}
    \right\}.
  \end{equation}
  Set
  \[
    v=\frac{2\gamma\sqrt{2\Lambda}}{1-2\Lambda}\in[0,1].
  \]
  By \eqref{eq:plt23-beta-rational},
  \[
    \gamma\beta(\gamma)
    \le\frac{(1-2\Lambda)^2v(1-v)}{4\Lambda\sqrt{2\Lambda}}.
  \]
  Therefore
  \[
    \frac{(\Lambda-q)\sqrt{2\Lambda}}{2}-\gamma\beta(\gamma)
    \ge\frac{\sqrt{2\Lambda}\,(1-2\Lambda)^2}{8\Lambda^2}
    (\xi-v+v^2).
  \]
  Furthermore,
  \[
    \Lambda_m\gamma^2=\Omega_m\rho v^2,
    \qquad
    2M_m(\Lambda-M)\frac{\sqrt{2\Lambda}\,(1-2\Lambda)^2}{8\Lambda^2}
    =\Omega_m.
  \]
  Substitution in \eqref{eq:plt23-one-d-gamma} proves
  \eqref{eq:plt23-huber-payment}.

  Finally suppose $\phi_\perp=0$.  The constraint in \eqref{eq:plt23-Sbar} then
  reads $b\gamma\ge\cH$, that is,
  \[
    b\gamma\ge\frac{(\Lambda-q)a_\phi}{2}
    \ge\frac{(\Lambda-q)\sqrt{2\Lambda}}{2}>0,
  \]
  so $b,\gamma>0$ and \eqref{eq:plt23-b-gamma-budget} still applies.  With the
  same variable $v$,
  \[
    \frac{(\Lambda-q)\sqrt{2\Lambda}}{2}
    \le b\gamma
    \le\frac{(1-2\Lambda)^2v(1-v)}{4\Lambda\sqrt{2\Lambda}},
  \]
  which is equivalent to $\xi\le v-v^2$.  Thus the positive-part term in
  \eqref{eq:plt23-psi-def} vanishes at this $v$, and
  \[
    \Lambda_m\gamma^2=\Omega_m\rho v^2\ge \Omega_m\psi(\xi,\rho).
  \]
  This completes the proof.
\end{proof}

It remains to say which pairs $(q,\Lambda)$ can occur.  The eigenfunction
$\phi$ restricts them further, through a lower bound on $\norm g$.

\begin{lemma}[Variance lower bound]\label{lem:plt23-forced-variance}
  One has
  \begin{equation}\label{eq:plt23-forced-variance}
    \norm g^2\ge
    \frac{\Lambda(\Lambda-q)^2}{2(1-2\Lambda)}.
  \end{equation}
  Consequently,
  \begin{equation}\label{eq:plt23-alpha-ceiling-poly}
    \Lambda^2+q\Lambda-q\le0,
  \end{equation}
  and hence
  \begin{equation}\label{eq:plt23-alpha-ceiling}
    q<\Lambda\le r(q):=
    \frac{\sqrt{q^2+4q}-q}{2}.
  \end{equation}
  Moreover $M>0$.
\end{lemma}

\begin{proof}
  Write $|\phi|=a_\phi\one+h$ with $h\perp\one$, so that
  $\norm h^2=1-a_\phi^2$.  The block form \eqref{eq:block-U} gives
  \[
    \ip{|\phi|}{T_U|\phi|}
    =qa_\phi^2+2a_\phi\ip gh+\ip h{Ah},
  \]
  and bounding $\ip gh\le\norm g\norm h$ by Cauchy--Schwarz and
  $\ip h{Ah}\le\Lambda\norm h^2$ by the maximality of $\Lambda$,
  \eqref{eq:plt23-absphi-lower} gives
  \[
    \Lambda\le qa_\phi^2+2a_\phi\norm g\sqrt{1-a_\phi^2}
    +\Lambda(1-a_\phi^2).
  \]
  Subtracting $\Lambda(1-a_\phi^2)$ and dividing by $a_\phi>0$,
  \[
    (\Lambda-q)a_\phi
    \le2\norm g\sqrt{1-a_\phi^2}.
  \]
  Squaring and writing $t=a_\phi^2$,
  \[
    \norm g^2\ge\frac{(\Lambda-q)^2}{4}\cdot\frac{t}{1-t},
  \]
  and $t\mapsto t/(1-t)$ is increasing, so \eqref{eq:plt23-aphi-lower} allows $t$ to
  be replaced by $2\Lambda$, which is \eqref{eq:plt23-forced-variance}.

  For the second claim, \eqref{eq:plt23-HS-bound} gives
  $2\norm g^2\le pq-\Lambda^2$, so \eqref{eq:plt23-forced-variance} and
  $1-2\Lambda>0$ give
  \[
    \Lambda(\Lambda-q)^2\le\bigl(pq-\Lambda^2\bigr)(1-2\Lambda).
  \]
  With $p=1-q$, the difference of the two sides is
  $-(1-\Lambda-q)(\Lambda^2+q\Lambda-q)$, and the factor $1-\Lambda-q$ is
  positive by \eqref{eq:plt23-M-Lambda-order}, so \eqref{eq:plt23-alpha-ceiling-poly} follows.
  The roots of $\Lambda^2+q\Lambda-q$ are $(-q\pm\sqrt{q^2+4q})/2$, whence
  \eqref{eq:plt23-alpha-ceiling}.

  Finally, because $\Lambda>q$, the right side of \eqref{eq:plt23-forced-variance}
  is strictly positive.  Thus $g\ne0$, and the bound \eqref{eq:plt23-HS-bound}
  gives $\Lambda^2<pq$.  Consequently $M>0$.
\end{proof}

Together, \Cref{thm:plt23-huber-elim} and \Cref{lem:plt23-forced-variance} show that
\eqref{eq:main-goal} follows once
\begin{equation}\label{eq:plt23-scalar-target}
  R_m\le \Omega_m\psi(\xi,\rho)
\end{equation}
is proved for every pair $(q,\Lambda)$ satisfying \eqref{eq:plt23-range} and
\eqref{eq:plt23-alpha-ceiling}.

\begin{proposition}[Dual form of $\psi$]\label{prop:plt23-huber-dual}
  For $\xi\ge0$ and $\rho>0$,
  \begin{equation}\label{eq:plt23-psi-dual}
    \psi(\xi,\rho)=
    \max_{0\le\lambda\le1}
    \left\{
    \lambda\xi-\frac{\lambda^2}{4(\rho+\lambda)}
    \right\}.
  \end{equation}
\end{proposition}

\begin{proof}
  Use
  \[
    \pos x=\max_{0\le\lambda\le1}\lambda x.
  \]
  The resulting function
  \[
    (\rho+\lambda)v^2-\lambda v+\lambda\xi
  \]
  is convex in $v$, affine in $\lambda$, and the two domains are compact
  convex intervals.  The elementary finite-dimensional minimax theorem (or
  Sion's theorem) permits the interchange of $\min_v$ and $\max_\lambda$.
  For fixed $\lambda$, the minimizer is
  \[
    v=\frac{\lambda}{2(\rho+\lambda)}\in[0,1/2],
  \]
  and substitution gives \eqref{eq:plt23-psi-dual}.
\end{proof}

The function $\psi$ has the quadratic-versus-linear structure of Huber's
robust loss~\cite{Huber1964}, and \eqref{eq:plt23-psi-dual} is an interchange of
the convex--concave type treated by Sion~\cite{Sion1958}; the special case
needed here is the one proved above.

Two choices of $\lambda$ will be used throughout.

\begin{corollary}[Two lower bounds for $\psi$]\label{cor:plt23-two-witnesses}
  If $2\rho\xi\le1$, then
  \begin{equation}\label{eq:plt23-witness-quadratic}
    \Omega_m\psi(\xi,\rho)
    \ge
    \frac{2\Lambda^3\Lambda_m(\Lambda-q)^2}{(1-2\Lambda)^2}
    \frac{1+4\xi}{1+2\xi}.
  \end{equation}
  If $2\rho\xi>1$, then
  \begin{equation}\label{eq:plt23-witness-linear}
    \Omega_m\psi(\xi,\rho)
    >M_m(\Lambda-M)\sqrt{2\Lambda}\,(\Lambda-q)
    \frac{1+4\xi}{2+4\xi}.
  \end{equation}
\end{corollary}

\begin{proof}
  In the first case, take $\lambda=2\rho\xi$ in
  \eqref{eq:plt23-psi-dual}.  This gives
  \[
    \psi(\xi,\rho)\ge
    \rho\xi^2\frac{1+4\xi}{1+2\xi}.
  \]
  A direct calculation from \eqref{eq:plt23-C-xi-rho} gives
  \[
    \Omega_m\rho\xi^2=\frac{2\Lambda^3\Lambda_m(\Lambda-q)^2}{(1-2\Lambda)^2}.
  \]

  In the second case, take $\lambda=1$.  Since
  $\rho>1/(2\xi)$,
  \[
    \frac{1}{4(1+\rho)}<\frac{\xi}{2+4\xi}.
  \]
  Therefore
  \[
    \psi(\xi,\rho)
    \ge\xi-\frac{1}{4(1+\rho)}
    >\xi\frac{1+4\xi}{2+4\xi}.
  \]
  Finally,
  \[
    \Omega_m\xi=M_m(\Lambda-M)\sqrt{2\Lambda}\,(\Lambda-q).
  \]
\end{proof}

\subsection{The dimensionless chart}\label{subsec:plt23-chart}

\tk{TODO (notation): $N$ defined below collides with the partition size $N$ of
\Cref{subsec:block-decomposition} (the $N\times N$ matrices of
\Cref{lem:series-expansions}); several later subsection titles use this $N$,
so rename one of the two.}
Set
\begin{equation}\label{eq:plt23-N-def}
  N=m-2.
\end{equation}
Thus $N\ge7$ is odd.  Introduce
\begin{equation}\label{eq:plt23-chart-basic}
  a=\frac {q}{\Lambda},
  \qquad
  \ell=\frac {M}{\Lambda},
  \qquad
  u=1-a,
  \qquad
  D=1+a^2+\ell^2,
  \qquad
  \sigma=\frac {p}{\Lambda}.
\end{equation}
The relation $M^2=pq-\Lambda^2$ is equivalent to
\begin{equation}\label{eq:plt23-chart-inverse-1}
  \Lambda=\frac {a}{D},
  \qquad q=\frac{a^2}{D},
  \qquad \sigma=\frac{1+\ell^2}{a}.
\end{equation}
Moreover,
\begin{equation}\label{eq:plt23-chart-d-e}
  \Lambda-q=\frac{au}{D},
  \qquad 1-2\Lambda=\frac{u^2+\ell^2}{D}.
\end{equation}

\begin{lemma}[Chart domain]\label{lem:plt23-chart-domain}
  The admissible parameters satisfy
  \begin{align*}
      &0<\ell<a<1, \qquad D<3a^2, \qquad \ell^2\ge a(1-a)=au, \nonumber
      \\
      &a>\frac{1+\sqrt{13}}{6}>\frac{3}{4}, \qquad 1<\sigma<2. \label{eq:plt23-chart-domain}
  \end{align*}
\end{lemma}

\begin{proof}
  The order $0<M<q<\Lambda$ gives $0<\ell<a<1$.  Since
  $q=a^2/D>1/3$, we have $D<3a^2$.  The ceiling
  $\Lambda^2+q\Lambda-q\le0$, after substitution from
  \eqref{eq:plt23-chart-inverse-1}, is exactly
  \[
    \ell^2\ge a-a^2=au.
  \]
  Hence $D\ge1+a$.  Combining this with $D<3a^2$ gives
  $3a^2-a-1>0$, which is the stated lower bound on $a$.  Finally,
  $\sigma=p/\Lambda>1$ because $p>1/2>\Lambda$, while
  $p<2/3$ and $\Lambda>q>1/3$ give $\sigma<2$.
\end{proof}

Now put
\begin{equation}\label{eq:plt23-zeta-v-def}
  \zeta=\frac{\ell^2}{u},
  \qquad v=\sigma-1,
  \qquad H_\zeta=\zeta+1+v.
\end{equation}
Solving these relations gives
\begin{equation}\label{eq:plt23-zeta-v-inverse}
  u=\frac {v}{H_\zeta},
  \qquad
  a=\frac{\zeta+1}{H_\zeta},
  \qquad
  \ell^2=\frac{\zeta v}{H_\zeta}.
\end{equation}
The frontier ceiling $\ell^2\ge au$ becomes
\begin{equation}\label{eq:plt23-zeta-domain}
  \zeta(\zeta+v)\ge1.
\end{equation}
Equivalently,
\begin{equation}\label{eq:plt23-zeta-simple-lower}
  \zeta\ge\frac{\sqrt{v^2+4}-v}{2}\ge1-\frac {v}{2}.
\end{equation}
In particular,
\begin{equation}\label{eq:plt23-zeta-lower}
  \zeta\ge a>\frac{3}{4},
  \qquad 0<v<1.
\end{equation}

For later use, define
\begin{equation}\label{eq:plt23-Q-chart}
  Q=v^2+v\zeta+2v+2\zeta+2.
\end{equation}
A direct substitution gives
\begin{equation}\label{eq:plt23-D-Q}
  D=\frac{(\zeta+1)Q}{H_\zeta^2}.
\end{equation}
Since $D<3a^2$,
\begin{equation}\label{eq:plt23-Q-upper}
  Q<3(\zeta+1).
\end{equation}
We shall repeatedly use the identity
\begin{equation}\label{eq:plt23-H-identity}
  \zeta H_\zeta^2-(\zeta+1)(\zeta+v)^2
  =\zeta^2+\zeta-v^2>0.
\end{equation}
The positivity follows from \eqref{eq:plt23-zeta-domain}:
\[
  \zeta^2+\zeta-v^2
  \ge1+\zeta(1-v)-v^2
  =(1-v)(1+v+\zeta)>0.
\]

\subsection{The normalized defect and finite quotients}

Normalize $R_m$ by $\Lambda^m$:
\begin{equation}\label{eq:plt23-F-def}
  F_N:=\frac{R_m}{\Lambda^m}
  =1-a^N-\ell^2(a^N-\ell^N).
\end{equation}
Define
\begin{equation}\label{eq:plt23-T-def}
  T_N=N-\zeta(a^N-\ell^N).
\end{equation}
Since $1-a^N\le N(1-a)=Nu$,
\begin{equation}\label{eq:plt23-F-T}
  F_N\le uT_N.
\end{equation}
If $T_N\le0$, then $R_m\le0$, and the scalar target
\eqref{eq:plt23-scalar-target} is immediate.  We may therefore assume
\begin{equation}\label{eq:plt23-T-positive}
  T_N>0
\end{equation}
whenever convenient.

The normalized versions of $\Lambda_m$ and $M_m$ are
\begin{align}
  \frac{\Lambda_m}{\Lambda^N}
  &=2\ell^N+(N+2)K_\Lambda,
  &
  K_\Lambda&=\frac{\sigma^{N+1}-1}{\sigma+1},
  \label{eq:plt23-KA-def}\\
  \frac{M_m}{\Lambda^N}
  &=2\ell^N+(N+2)K_M,
  &
  K_M&=\frac{\sigma^{N+1}-\ell^{N+1}}{\sigma+\ell}.
  \label{eq:plt23-KL-def}
\end{align}
Because $N+1$ is even,
\begin{align}
  K_\Lambda&=(\sigma-1)
  \sum_{j=0}^{(N-1)/2}\sigma^{2j},
  \label{eq:plt23-KA-sum}\\
  K_M&=(\sigma-\ell)
  \sum_{j=0}^{(N-1)/2}
  \sigma^{N-1-2j}\ell^{2j}.
  \label{eq:plt23-KL-sum}
\end{align}

\subsection{Three elementary bounds for \texorpdfstring{$T_N$}{TN}}

\begin{lemma}[Bounds for $T_N$]\label{lem:plt23-T-bounds}
  For every admissible chart point and every odd $N\ge7$,
  \begin{align}
    T_N&\le N(1+v)-\zeta(1-v^{N/2}),
    \label{eq:plt23-T-bound-one}\\
    T_N&<N(1-\ell^{N+1}).
    \label{eq:plt23-T-bound-two}
  \end{align}
  If $T_N>0$, then also
  \begin{equation}\label{eq:plt23-T-tangent-bound}
    \zeta T_N<\frac{3N(N+1)}{2}(1+v).
  \end{equation}
\end{lemma}

\begin{proof}
  By Bernoulli's inequality,
  \[
    a^N\ge1-Nu.
  \]
  Also $\ell^2=\zeta v/H_\zeta<v$, so
  $\ell^N<v^{N/2}$.  Since $\zeta u< v$,
  \begin{align*}
    T_N
    &=N-\zeta(a^N-\ell^N)\\
    &\le N-\zeta+N\zeta u+\zeta v^{N/2}\\
    &\le N(1+v)-\zeta(1-v^{N/2}),
  \end{align*}
  proving \eqref{eq:plt23-T-bound-one}.

  Next, $D<3a^2$ gives
  \[
    \ell^2<2a^2-1<(2a-1)^2.
  \]
  Thus $a-\ell>1-a=u$.  By convexity of $x\mapsto x^N$,
  \[
    a^N-\ell^N>N(a-\ell)\ell^{N-1}>Nu\ell^{N-1}.
  \]
  Multiplication by $\zeta=\ell^2/u$ gives
  \eqref{eq:plt23-T-bound-two}.

  Assume now $T_N>0$.  From \eqref{eq:plt23-T-bound-one},
  \begin{equation}\label{eq:plt23-zeta-upper-Z}
    \zeta<Z:=\frac{N(1+v)}{1-v^{N/2}}.
  \end{equation}
  Put $r=(N+1)/2$ and
  \[
    g(t)=t\left[1-
    \left(\frac{tv}{t+1+v}\right)^r\right].
  \]
  If $\theta=(1+v)/(t+1+v)$, then
  \[
    g'(t)=1-[v(1-\theta)]^r(1+r\theta)\ge0,
  \]
  because $v^r\le1$ and Bernoulli's inequality gives
  $(1-\theta)^{-r}\ge1+r\theta$.  Therefore,
  using \eqref{eq:plt23-T-bound-two} and \eqref{eq:plt23-zeta-upper-Z},
  \begin{equation}\label{eq:plt23-gZ-step}
    \zeta T_N<N g(\zeta)<Ng(Z).
  \end{equation}
  Let
  \[
    \eta=v^{N/2},
    \qquad
    \omega=\frac{Nv}{N+1-\eta}.
  \]
  A calculation gives
  \[
    Ng(Z)=
    \frac{N^2(1+v)}{1-\eta}
    \left(1-\omega^{(N+1)/2}\right).
  \]
  As a function of $\eta\in(0,1)$,
  \[
    f(\eta)=\omega^{(N+1)/2}
    =\left(\frac{N}{N+1-\eta}\right)^{(N+1)/2}
    \eta^{(N+1)/N}
  \]
  is convex.  Indeed, with
  $r=(N+1)/2$, $s=(N+1)/N$, and $d=N+1-\eta$,
  \[
    \frac{f''(\eta)}{f(\eta)}
    =\frac{r(r+1)}{d^2}
    +\frac{2rs}{d\eta}
    +\frac{s(s-1)}{\eta^2}>0.
  \]
  Moreover,
  \[
    f(1)=1,
    \qquad
    f'(1)=\frac{3(N+1)}{2N}.
  \]
  The tangent inequality for a convex function therefore yields
  \[
    1-f(\eta)
    \le\frac{3(N+1)}{2N}(1-\eta).
  \]
  Substitution in \eqref{eq:plt23-gZ-step} proves
  \eqref{eq:plt23-T-tangent-bound}.
\end{proof}

\subsection{The regime \texorpdfstring{$2\rho\xi\le1$}{2 rho xi <= 1}}\label{subsec:plt23-quadratic}

Assume
\begin{equation}\label{eq:plt23-quadratic-branch}
  2\rho\xi\le1.
\end{equation}
We shall prove \eqref{eq:plt23-scalar-target} for every odd $N\ge7$, hence for
every odd $m\ge9$.

By \eqref{eq:plt23-witness-quadratic}, discarding the factor
$(1+4\xi)/(1+2\xi)>1$, and using \eqref{eq:plt23-KA-def},
\begin{equation}\label{eq:plt23-quadratic-payment-start}
  \frac{\Omega_m\psi(\xi,\rho)}{\Lambda^m}
  \ge2(N+2)\Lambda\left(\frac {d}{e}\right)^2K_\Lambda.
\end{equation}
Since the summands in \eqref{eq:plt23-KA-sum} have geometric mean
$\sigma^{(N-1)/2}$, AM--GM gives
\begin{equation}\label{eq:plt23-KA-AMGM}
  K_\Lambda\ge\frac{N+1}{2}
  v\sigma^{(N-1)/2}.
\end{equation}
Using \eqref{eq:plt23-chart-inverse-1}, \eqref{eq:plt23-chart-d-e}, and
$v=u(1+\zeta)/a$, the right side of
\eqref{eq:plt23-quadratic-payment-start} is at least
\begin{equation}\label{eq:plt23-quad-payment-M}
  u(N+2)(N+1)
  \frac{a^2(1+\zeta)}{D(u+\zeta)^2}
  \sigma^{(N-1)/2}.
\end{equation}

\begin{lemma}[Coefficient estimate]\label{lem:plt23-quad-coeff}
  One has
  \begin{equation}\label{eq:plt23-M-lower}
    \frac{a^2(1+\zeta)}{D(u+\zeta)^2}
    >\frac{1}{3\zeta}.
  \end{equation}
\end{lemma}

\begin{proof}
  Using \eqref{eq:plt23-zeta-v-inverse} and \eqref{eq:plt23-D-Q}, the left side is
  \[
    M=\frac{H_\zeta^2}{(\zeta+v)^2Q}.
  \]
  By \eqref{eq:plt23-H-identity},
  \[
    \frac{H_\zeta^2}{(\zeta+v)^2}
    >\frac{\zeta+1}{\zeta},
  \]
  while \eqref{eq:plt23-Q-upper} gives $Q<3(\zeta+1)$.  Multiplication proves
  \eqref{eq:plt23-M-lower}.
\end{proof}

Combining \eqref{eq:plt23-quad-payment-M} and \eqref{eq:plt23-M-lower},
\begin{equation}\label{eq:plt23-quad-payment-final}
  \frac{\Omega_m\psi(\xi,\rho)}{\Lambda^m}
  >u\frac{(N+2)(N+1)}{3\zeta}
  (1+v)^{(N-1)/2}.
\end{equation}
In view of \eqref{eq:plt23-F-T}, it remains to pay for $uT_N$.

\begin{proposition}[The regime $2\rho\xi\le1$]\label{prop:plt23-quadratic-branch}
  For every odd $N\ge7$, assumption \eqref{eq:plt23-quadratic-branch} implies
  \[
    R_m\le \Omega_m\psi(\xi,\rho).
  \]
\end{proposition}

\begin{proof}
  If $T_N\le0$, the claim follows from \eqref{eq:plt23-F-T}.  Assume
  $T_N>0$.

  First, \eqref{eq:plt23-T-bound-one} implies
  \[
    \zeta T_N
    \le N(1+v)\zeta-(1-v^{N/2})\zeta^2
    \le\frac{N^2(1+v)^2}{4(1-v^{N/2})}.
  \]
  Thus \eqref{eq:plt23-quad-payment-final} is sufficient whenever
  \begin{equation}\label{eq:plt23-E-N}
    E_N(v):=
    \frac{4(N+2)(N+1)}{3N^2}
    (1-v^{N/2})(1+v)^{(N-5)/2}\ge1.
  \end{equation}
  Second, the tangent estimate \eqref{eq:plt23-T-tangent-bound} shows that
  \eqref{eq:plt23-quad-payment-final} is sufficient whenever
  \begin{equation}\label{eq:plt23-H-N}
    H_N(v):=
    \frac{2(N+2)}{9N}(1+v)^{(N-3)/2}\ge1.
  \end{equation}
  We now cover the whole interval $0<v<1$.

  If $N\ge9$ and $v\le3/5$, then
  \[
    E_N(v)\ge\frac{4}{3}\left(1-\left(\frac{3}{5}\right)^4\right)
    =\frac{2176}{1875}>1.
  \]
  If $N\ge9$ and $v\ge3/5$, then
  \[
    \frac{H_{N+2}(v)}{H_N(v)}
    =(1+v)\frac{N(N+4)}{(N+2)^2}>1,
  \]
  so $H_N(v)$ increases under $N\mapsto N+2$, and
  \[
    H_9(3/5)=\frac{11264}{10125}>1.
  \]

  It remains to take $N=7$.  For $v\le9/10$,
  \[
    E_7(v)\ge\frac{96}{49}(1-v^3)(1+v)\ge1.
  \]
  Indeed, $1+v-v^3-v^4$ is concave, so its minimum on
  $[0,9/10]$ occurs at an endpoint; at $v=9/10$ it equals
  $5149/10000>49/96$.  For $v\ge9/10$,
  \[
    H_7(v)\ge\frac{2}{7}\left(\frac{19}{10}\right)^2
    =\frac{361}{350}>1.
  \]
  Thus either \eqref{eq:plt23-E-N} or \eqref{eq:plt23-H-N} holds in every case, proving
  the proposition.
\end{proof}

\subsection{The regime \texorpdfstring{$2\rho\xi>1$}{2 rho xi > 1}}\label{subsec:plt23-linear}

Assume
\begin{equation}\label{eq:plt23-linear-branch}
  2\rho\xi>1.
\end{equation}
Abbreviate
\begin{equation}\label{eq:plt23-varphi-cxi}
  \varphi=\frac{1+4\xi}{2+4\xi},
  \qquad
  c_\xi=\sqrt{2\Lambda}\,\varphi.
\end{equation}
By \eqref{eq:plt23-witness-linear}, \eqref{eq:plt23-KL-def}, and
$\Lambda-M=\Lambda(1-\ell)$, $\Lambda-q=\Lambda u$,
\begin{equation}\label{eq:plt23-linear-payment-start}
  \frac{\Omega_m\psi(\xi,\rho)}{\Lambda^m}
  >u(N+2)c_\xi(1-\ell)K_M.
\end{equation}
From \eqref{eq:plt23-KL-sum},
\begin{equation}\label{eq:plt23-KL-first-term}
  K_M\ge(\sigma-\ell)\sigma^{N-1}.
\end{equation}
In view of \eqref{eq:plt23-F-T}, it is therefore enough to prove
\begin{equation}\label{eq:plt23-linear-core}
  T_N\le
  (N+2)c_\xi(1-\ell)(\sigma-\ell)\sigma^{N-1}.
\end{equation}

The factor $c_\xi$ may be small, but only where a short power of
$\sigma$ compensates for it.

\begin{lemma}[Compensation]\label{lem:plt23-compensation}
  On the whole admissible chart,
  \begin{equation}\label{eq:plt23-compensation}
    c_\xi\sigma^{\zeta/4}\ge1-\ell^2.
  \end{equation}
\end{lemma}

\begin{proof}
  We first show
  \begin{equation}\label{eq:plt23-sqrt-compensation}
    \sqrt{2\Lambda}\ge1-\ell^2.
  \end{equation}
  Using $\Lambda=a/D$, the square of \eqref{eq:plt23-sqrt-compensation} is
  \[
    2a\ge(1+a^2+y)(1-y)^2,
    \qquad y=\ell^2.
  \]
  The difference between the two sides has derivative
  \[
    (1-y)(1+2a^2+3y)>0.
  \]
  The smallest admissible $y$ is $a(1-a)$, where the difference is
  \[
    (1-a)(a^4+a^2+2a-1)>0
  \]
  because $a>3/4$.  This proves \eqref{eq:plt23-sqrt-compensation}.

  Next we show $\varphi\sigma^{\zeta/4}\ge1$.  From
  \eqref{eq:plt23-C-xi-rho}, \eqref{eq:plt23-zeta-v-inverse}, and
  \eqref{eq:plt23-D-Q},
  \begin{equation}\label{eq:plt23-xi-zeta-v}
    \xi=\frac{4H_\zeta^2}
    {v(\zeta+v)^2Q}.
  \end{equation}
  Equations \eqref{eq:plt23-Q-upper} and \eqref{eq:plt23-H-identity} imply
  \begin{equation}\label{eq:plt23-xi-comp-bound}
    4\zeta v\xi>\frac{16}{3}.
  \end{equation}
  We claim that
  \begin{equation}\label{eq:plt23-key-log-comp}
    \zeta v(1+4\xi)>2(v+2).
  \end{equation}
  Indeed, the left side is larger than $\zeta v+16/3$.  If
  $\zeta\ge2$, this is immediate.  If $\zeta<2$, then
  \[
    v(2-\zeta)<\frac{5}{4}<\frac{4}{3}
  \]
  because $v<1$ and $\zeta>3/4$, which again proves
  \eqref{eq:plt23-key-log-comp}.

  The elementary inequality
  \[
    \log(1+v)\ge\frac{2v}{v+2}
  \]
  now gives
  \[
    \frac{\zeta}{4}\log\sigma>
    \frac{1}{1+4\xi}.
  \]
  Exponentiating and using $e^t>1+t$,
  \[
    \sigma^{\zeta/4}>
    1+\frac{1}{1+4\xi}
    =\frac{2+4\xi}{1+4\xi}=\frac{1}{\varphi}.
  \]
  Together with \eqref{eq:plt23-sqrt-compensation}, this proves
  \eqref{eq:plt23-compensation}.
\end{proof}

\subsection{The cases \texorpdfstring{$\zeta\ge N$}{zeta >= N} and \texorpdfstring{$v\ge5/8$}{v >= 5/8}}

We first dispose of the range $\zeta\ge N$ when $v$ is not large.

\begin{lemma}[The regime $\zeta\ge N$]\label{lem:plt23-linear-high-zeta}
  If $\zeta\ge N$, $0<v\le5/8$, and $N\ge7$ is odd, then
  \eqref{eq:plt23-linear-core} holds.
\end{lemma}

\begin{proof}
  By \eqref{eq:plt23-T-bound-one} and $\zeta\ge N$,
  \begin{equation}\label{eq:plt23-T-high-zeta}
    T_N\le N(v+v^{N/2}).
  \end{equation}
  Put $y=\sqrt v$.  Since $\ell^2<v$,
  \[
    (1-\ell)(\sigma-\ell)
    \ge(1-y)(1-y+y^2).
  \]
  Also $q<\Lambda$ and $p=(1+v)\Lambda$ imply
  $\Lambda>1/(2+v)$, while $\varphi>1/2$.  Hence
  \begin{equation}\label{eq:plt23-cxi-crude}
    c_\xi>\frac{1}{\sqrt{2(2+v)}}.
  \end{equation}
  It is enough to show
  \begin{equation}\label{eq:plt23-h-sufficient}
    h(y):=
    \frac{(1-y)(1-y+y^2)(1+y^2)^5}{y^2}
    >\sqrt{2(2+y^2)}\,(1+y^{N-2}).
  \end{equation}
  Indeed, the omitted factors
  $(N+2)/N$ and $(1+y^2)^{N-6}$ are at least $1$.

  We claim that $h$ is decreasing on $(0,4/5]$.  Differentiation gives
  \[
    h'(y)=-\frac{(1+y^2)^4}{y^3}S(y),
  \]
  where
  \[
    S(y)=11y^5-20y^4+19y^3-8y^2-2y+2.
  \]
  For $0\le y\le1/2$,
  \[
    S(y)\ge2-2y-8y^2+9y^3\ge\frac{1}{8}.
  \]
  Here the first inequality follows from
  $y^3(11y^2-20y+10)\ge0$, and the cubic on the right is decreasing on this
  interval.  For $1/2\le y\le4/5$, write $y=2/3+t$.  If $t\ge0$, then $0\le t\le2/15$, and
  \begin{align*}
    S(y)&=\frac{58}{243}-\frac{14}{81}t
    +\frac{250}{27}t^2+\frac{131}{9}t^3
    +\frac{50}{3}t^4+11t^5\\
    &\ge \frac{58}{243}-\frac{14}{81}\frac{2}{15}
    =\frac{262}{1215}>0.
  \end{align*}
  If $t=-s\le0$, where $0\le s\le1/6$, then
  \begin{align*}
    S(y)={}&\frac{58}{243}+\frac{14}{81}s
    +\frac{250}{27}s^2-\frac{131}{9}s^3
    +\frac{50}{3}s^4-11s^5\\
    \ge{}&\frac{58}{243}+\frac{14}{81}s
    +\frac{1465}{216}s^2>0.
  \end{align*}
  Thus $h'\le0$.

  Since $v\le5/8$, we have $y<4/5$, and therefore
  \[
    h(y)\ge h(4/5)
    =\frac{2432980221}{781250000}>\frac{46}{15}.
  \]
  On the other hand, $N\ge7$ gives
  \[
    \sqrt{2(2+y^2)}(1+y^{N-2})
    \le\sqrt{\frac{21}{4}}
    \left(1+\left(\frac{4}{5}\right)^5\right)
    <\frac{46}{15};
  \]
  after squaring, the final difference is
  $52766011/351562500>0$.  This proves
  \eqref{eq:plt23-h-sufficient} and hence \eqref{eq:plt23-linear-core}.
\end{proof}

The next estimate handles all $\zeta$ once $v$ is large.

\begin{lemma}[Large $v$]\label{lem:plt23-linear-large-v}
  If $v\ge5/8$ and $N\ge7$ is odd, then
  \eqref{eq:plt23-linear-core} holds.
\end{lemma}

\begin{proof}
  By \eqref{eq:plt23-T-bound-two},
  \begin{equation}\label{eq:plt23-T-geometric-cancel}
    T_N<N(N+1)(1-\ell).
  \end{equation}
  Using \eqref{eq:plt23-cxi-crude}, $\ell\le\sqrt v$, and
  $\sigma=1+v$, it is enough, after cancelling $1-\ell$, to prove
  \begin{equation}\label{eq:plt23-large-v-target}
    \frac{N+2}{\sqrt{2(2+v)}}
    (1+v-\sqrt v)(1+v)^{N-1}
    \ge N(N+1).
  \end{equation}
  The left side is increasing in $v$ for $v\ge5/8$: the factor
  $1+v-\sqrt v$ is increasing for $v>1/4$, and
  \[
    \frac{\dd}{\dd v}\log\frac{(1+v)^{N-1}}{\sqrt{2+v}}
    =\frac{N-1}{1+v}-\frac{1}{2(2+v)}>0.
  \]
  After division by $N(N+1)$, its ratio under $N\mapsto N+2$ is at least
  \[
    \left(\frac{13}{8}\right)^2
    \frac{N(N+1)(N+4)}{(N+2)^2(N+3)}>1,
  \]
  because, after clearing denominators, the excess is
  \[
    105N^3+397N^2-348N-768>0
    \qquad(N\ge7).
  \]
  It remains to check $N=7$, $v=5/8$.  The estimates
  $\sqrt{10}<19/6$ and $\sqrt{21}<23/5$ give
  \[
    \frac{9}{\sqrt{21/4}}
    \left(\frac{13}{8}-\sqrt{\frac{5}{8}}\right)
    \left(\frac{13}{8}\right)^6
    >\frac{75}{23}\left(\frac{13}{8}\right)^6
    >56,
  \]
  the final excess being $24369203/6029312$.  Thus
  \eqref{eq:plt23-large-v-target} holds.
\end{proof}

\subsection{The regime \texorpdfstring{$\zeta\le N$}{zeta<=N}: the growth lemma}

It remains, except for $N=7$, to handle $\zeta\le N$.  The compensation
lemma reduces this to one variable.

\begin{lemma}[One-variable growth]\label{lem:plt23-J-growth}
  Let $N\ge9$ be odd.  If $0\le x<1$ and
  \begin{equation}\label{eq:plt23-J-domain}
    1-x\ge\frac{2x^2}{N},
  \end{equation}
  then
  \begin{equation}\label{eq:plt23-J-ineq}
    J_N(x):=
    \frac{N+2}{N}(1-x)^2(1+x^3)
    \left(\frac{N(1+x^2)}{N-x^2}\right)^{3N/4-1}
    \ge1.
  \end{equation}
\end{lemma}

\begin{proof}
  Put $A=3N/4-1$.

  \smallskip
  \noindent\emph{Step 1: $0\le x\le2/5$.}
  Taking logarithms in \eqref{eq:plt23-J-ineq}, and using
  \begin{align*}
    \log(1+2/N)&\ge\frac{2}{N}-\frac{2}{N^2},\\
    \log(1+x^2)&\ge x^2-\frac{x^4}{2},\\
    -\log(1-x^2/N)&\ge\frac{x^2}{N},
  \end{align*}
  together with
  \begin{equation}\label{eq:plt23-log-shape-bound}
    2\log(1-x)+\log(1+x^3)
    \ge-2x-x^2-\frac{x^4}{1-x}-\frac{x^6}{2},
  \end{equation}
  we obtain
  \begin{align}
    \log J_N(x)\ge\mathcal L_N(x):={}&
    \frac{2}{N}-\frac{2}{N^2}-2x \notag\\
    &+\left[A\left(1+\frac{1}{N}\right)-1\right]x^2
    -\frac {A}{2}x^4-\frac{x^4}{1-x}-\frac{x^6}{2}. \label{eq:plt23-L-N-def}
  \end{align}
  For completeness, \eqref{eq:plt23-log-shape-bound} follows by writing
  \[
    2\log(1-x)
    =-2x-x^2-\frac{2}{3}x^3-2\sum_{j\ge4}\frac{x^j}{j}
    \ge-2x-x^2-\frac{2}{3}x^3-\frac{x^4}{1-x}
  \]
  and using $\log(1+x^3)\ge x^3-x^6/2$.

  On $0\le x\le1/4$,
  \[
    x^4\le\frac{x^2}{16},
    \qquad
    \frac{x^4}{1-x}\le\frac{x^2}{12},
    \qquad
    \frac{x^6}{2}\le\frac{x^2}{512}.
  \]
  Thus
  \[
    \mathcal L_N(x)\ge c_N-2x+B_Nx^2,
  \]
  where
  \[
    c_N=\frac{2}{N}-\frac{2}{N^2},
    \qquad
    B_N=\frac{1116N^2-2003N-1536}{1536N}.
  \]
  For $N\ge9$, $B_N>0$, and
  \begin{equation}\label{eq:plt23-discriminant-check}
    B_Nc_N-1
    =\frac{348N^3-3119N^2+467N+1536}{768N^3}>0.
  \end{equation}
  The numerator equals $6792$ at $N=9$; its derivative is $28889$
  there and is increasing afterwards.  Hence the quadratic has positive global
  minimum, and $\mathcal L_N(x)>0$ on $[0,1/4]$.

  On $1/4\le x\le2/5$, differentiation in the real parameter $N$ gives
  \[
    \frac{\partial\mathcal L_N}{\partial N}
    \ge\frac{93}{2048}-\frac{2}{81}>0.
  \]
  It is therefore enough to take $N=9$.  Direct differentiation gives
  \begin{equation}\label{eq:plt23-L9-data}
    \mathcal L_9(1/4)=\frac{11795}{663552},
    \qquad
    \mathcal L_9'(1/4)=\frac{1295}{3072},
    \qquad
    \mathcal L_9''(2/5)=\frac{49}{3375}.
  \end{equation}
  Moreover,
  \[
    \mathcal L_9'''(x)=-\frac{3xR(x)}{(1-x)^4},
  \]
  where
  \[
    R(x)=20x^6-80x^5+143x^4-174x^3+166x^2-104x+31.
  \]
  For $0\le x\le2/5$,
  \[
    R(x)\ge31-104x+\frac{2282}{25}x^2
    \ge\frac{2503}{625}>0.
  \]
  Indeed, the difference from the first quadratic is
  \[
    \frac{x^2}{25}
    (500x^4-2000x^3+3575x^2-4350x+1868).
  \]
  For $x\le2/5$, the polynomial in parentheses is at least
  \[
    500x^4+(3575-800)x^2+(1868-1740)>0,
  \]
  because $-2000x^3\ge-800x^2$ and $-4350x\ge-1740$.
  The quadratic $31-104x+(2282/25)x^2$ is decreasing on this interval and
  has value $2503/625$ at $x=2/5$.  Thus
  $\mathcal L_9'''<0$.  Since $\mathcal L_9''(2/5)>0$, we have
  $\mathcal L_9''>0$ throughout the interval; the first two inequalities in
  \eqref{eq:plt23-L9-data} then imply $\mathcal L_9>0$.

  \smallskip
  \noindent\emph{Step 2: $2/5\le x\le\beta_N$.}
  Here $\beta_N$ is the positive root of
  \[
    1-\beta_N=\frac{2\beta_N^2}{N};
  \]
  condition \eqref{eq:plt23-J-domain} is exactly $x\le\beta_N$.  Since
  $N/(N-x^2)>1$, it is enough to prove
  \begin{equation}\label{eq:plt23-K-N-def}
    K_N(x):=
    \frac{N+2}{N}(1-x)^2(1+x^3)(1+x^2)^{3N/4-1}\ge1.
  \end{equation}
  The logarithmic derivative of $K_N$ has the opposite sign to
  \begin{align}
    P_N(x)={}&3Nx^5-3Nx^4+3Nx^2-3Nx
    +6x^5-2x^4\\
    &+10x^3-6x^2+4x+4.
    \label{eq:plt23-P-N}
  \end{align}
  At $x=2/5$,
  \[
    P_N(2/5)=-\frac{2(1197N-8266)}{3125}<0,
  \]
  while
  \[
    P_N''(2/5)=\frac{6(17N+66)}{25}>0
  \]
  and
  \[
    P_N'''(x)=12\left[N(15x^2-6x)+30x^2-4x+5\right]>0
    \qquad(x\ge2/5).
  \]
  Thus $P_N$ is convex on the interval and, starting negative, can cross
  zero at most once.  Therefore $K_N$ has no interior minimum: it is enough
  to check $x=2/5$ and $x=\beta_N$.

  At $x=2/5$,
  \[
    K_9(2/5)=\frac{1463}{3125}
    \left(\frac{29}{25}\right)^{23/4}.
  \]
  Taylor's formula with positive third derivative gives
  \[
    \left(1+\frac{4}{25}\right)^{23/4}
    \ge1+\frac{23}{4}\frac{4}{25}
    +\frac{1}{2}\frac{23}{4}\frac{19}{4}
    \left(\frac{4}{25}\right)^2
    =\frac{2837}{1250}.
  \]
  Hence
  \[
    K_9(2/5)\ge\frac{4150531}{3906250}>1.
  \]
  Furthermore,
  \[
    \frac{K_{N+2}(2/5)}{K_N(2/5)}
    =\frac{N(N+4)}{(N+2)^2}
    \left(\frac{29}{25}\right)^{3/2}
    \ge\frac{117}{121}\frac{31}{25}>1,
  \]
  where Bernoulli's inequality was used in the last step.  Thus the first
  endpoint is settled for every odd $N\ge9$.

  For the second endpoint, use
  $(1-\beta_N)^2=4\beta_N^4/N^2$.  The function
  $h_N(x)=1-x-2x^2/N$ is strictly decreasing, and $h_N(\beta_N)=0$.
  The exact evaluations
  \[
    h_9(21/25)=\frac{2}{625},\qquad
    h_{11}(6/7)=\frac{5}{539},\qquad
    h_{13}(7/8)=\frac{3}{416}
  \]
  show that $\beta_9>21/25$, $\beta_{11}>6/7$, and
  $\beta_N\ge7/8$ for every $N\ge13$.  When $N=9$, hence
  \begin{align*}
    K_9(\beta_9)
    &>\frac{11}{9}\frac{4}{81}
    \frac{497}{1000}\frac{159}{100}
    \left(\frac{17}{10}\right)^5\frac{37}{25}\\
    &=\frac{15221984467459}{15187500000000}>1.
  \end{align*}
  Here we used
  \[
    \left(\frac{21}{25}\right)^4>\frac{497}{1000},
    \quad
    1+\left(\frac{21}{25}\right)^3>\frac{159}{100},
    \quad
    1+\left(\frac{21}{25}\right)^2>\frac{17}{10},
  \]
  and
  \[
    \left(\frac{17}{10}\right)^{3/4}>\frac{37}{25}.
  \]
  The last inequality follows after raising both sides to the fourth power.

  When $N=11$, $\beta_{11}>6/7$, and the coarser estimates
  $\beta_{11}^4>1/2$, $1+\beta_{11}^3>3/2$, and
  $1+\beta_{11}^2>17/10$ give
  \[
    K_{11}(\beta_{11})
    >\frac{13}{11}\frac{4}{121}\frac{1}{2}\frac{3}{2}
    \left(\frac{17}{10}\right)^7
    =\frac{16003208247}{13310000000}>1.
  \]

  Finally, for $N\ge13$, $\beta_N\ge7/8$.  Dropping the factors
  $(N+2)/N$ and $1+\beta_N^3$,
  \[
    K_N(\beta_N)
    \ge\frac{4(7/8)^4}{N^2}
    \left(\frac{113}{64}\right)^{3N/4-1}.
  \]
  At $N=13$, this is larger than
  \[
    \frac{4(7/8)^4}{13^2}\left(\frac{7}{4}\right)^8
    =\frac{13841287201}{11341398016}>1.
  \]
  The ratio of these lower bounds under $N\mapsto N+2$ is at least
  \[
    \left(\frac{13}{15}\right)^2
    \left(\frac{7}{4}\right)^{3/2}>1.
  \]
  This completes all endpoint checks and the proof of
  \eqref{eq:plt23-J-ineq}.
\end{proof}

\begin{lemma}[The regime $\zeta\le N$]\label{lem:plt23-linear-low-zeta}
  If $\zeta\le N$ and $N\ge9$ is odd, then
  \eqref{eq:plt23-linear-core} holds.
\end{lemma}

\begin{proof}
  By \Cref{lem:plt23-compensation},
  \begin{align}
    &(N+2)c_\xi(1-\ell)(\sigma-\ell)\sigma^{N-1}\notag\\
    &\quad\ge
    (N+2)(1-\ell)^2(1+\ell^3)
    \sigma^{N-1-\zeta/4}.
    \label{eq:plt23-linear-comp-expanded}
  \end{align}
  Indeed, $(1+\ell)(\sigma-\ell)\ge1+\ell^3$ because
  $\sigma\ge1+\ell^2$.

  Since $\zeta=\ell^2/u\le N$,
  \[
    u\ge\frac{\ell^2}{N},
    \qquad
    a\le\frac{N-\ell^2}{N}.
  \]
  Thus
  \begin{equation}\label{eq:plt23-sigma-lower-J}
    \sigma=\frac{1+\ell^2}{a}
    \ge\frac{N(1+\ell^2)}{N-\ell^2}.
  \end{equation}
  Also $N-1-\zeta/4\ge3N/4-1$.  Finally,
  $\ell<2a-1$ from the proof of \eqref{eq:plt23-T-bound-two}, so
  \[
    1-\ell>2u=\frac{2\ell^2}{\zeta}
    \ge\frac{2\ell^2}{N}.
  \]
  Applying \Cref{lem:plt23-J-growth} with $x=\ell$ to
  \eqref{eq:plt23-linear-comp-expanded} shows that its right side is at least
  $N$.  Since $T_N<N$, \eqref{eq:plt23-linear-core} follows.
\end{proof}

\subsection{The exceptional case \texorpdfstring{$N=7$, i.e. $m=9$}{N=7, i.e. m=9}}

The preceding lemmas already cover $N=7$ whenever either
$\zeta\ge7$ or $v\ge5/8$.  It remains to prove
\eqref{eq:plt23-linear-core} under
\begin{equation}\label{eq:plt23-N7-corner}
  N=7,
  \qquad \zeta\le7,
  \qquad v\le\frac{5}{8}.
\end{equation}

\begin{lemma}[The small-$v$ part of $N=7$]\label{lem:plt23-N7-small-v}
  Under \eqref{eq:plt23-N7-corner}, if $0<v\le1/4$, then
  \eqref{eq:plt23-linear-core} holds.
\end{lemma}

\begin{proof}
  Put $s=\sqrt v$.  Since
  \[
    \frac{\ell^2}{v}=\frac{\zeta}{\zeta+1+v}
    \le\frac{7}{8+v}<\frac{7}{8}<\left(\frac{15}{16}\right)^2,
  \]
  we have
  \begin{equation}\label{eq:plt23-ell-small-v}
    \ell<\frac{15}{16}s.
  \end{equation}

  We next prove
  \begin{equation}\label{eq:plt23-xi-small-v}
    \xi\ge\frac{1}{4v}.
  \end{equation}
  Let $t=\zeta+v$.  Formula \eqref{eq:plt23-xi-zeta-v} becomes
  \[
    4v\xi=
    \frac{16(t+1)^2}{t^2(t(v+2)+2)}.
  \]
  The reciprocal factor
  \[
    G(t,v)=\frac{t^2(t(v+2)+2)}{(t+1)^2}
  \]
  is increasing in both variables, since
  \[
    \frac{\partial G}{\partial t}
    =\frac{t(t^2v+2t^2+3tv+6t+4)}{(t+1)^3}>0,
    \qquad
    \frac{\partial G}{\partial v}=\frac{t^3}{(t+1)^2}>0.
  \]
  Since
  $t\le7+1/4=29/4$, its maximum is
  \[
    \frac{246413}{17424}<16.
  \]
  This proves \eqref{eq:plt23-xi-small-v}.

  The function $(1+4\xi)/(2+4\xi)$ is increasing in $\xi$, so
  \eqref{eq:plt23-xi-small-v} gives
  \[
    \varphi\ge\frac{1+v}{1+2v}.
  \]
  Also $\Lambda>1/(2+v)$, hence
  \[
    c_\xi\sigma^6
    \ge\frac{(1+v)^7}{(1+2v)\sqrt{1+v/2}}.
  \]
  We shall use the elementary lower bound
  \begin{equation}\label{eq:plt23-cxi-small-v-poly}
    \frac{(1+v)^7}{(1+2v)\sqrt{1+v/2}}
    \ge\left(1-\frac{7v}{8}\right)(1+5v+11v^2).
  \end{equation}
  Both sides are positive.  If $M(v)$ and $R(v)$ denote respectively
  the left and right sides of \eqref{eq:plt23-cxi-small-v-poly}, then exact
  expansion gives
  \begin{align}
    &128(1+2v)^2(1+v/2)\bigl(M(v)^2-R(v)^2\bigr) \notag\\
    &\quad=v\bigl(128v^{13}+1792v^{12}+11648v^{11}+46592v^{10}
    +128128v^9 \notag\\
    &\qquad\quad+232540v^8+345884v^7+492971v^6+464196v^5 \notag\\
    &\qquad\quad+258097v^4+86924v^3+18035v^2+2254v+160\bigr), \label{eq:plt23-small-v-square-diff}
  \end{align}
  which is nonnegative.

  For the defect, use \eqref{eq:plt23-zeta-simple-lower}.  From
  \eqref{eq:plt23-T-bound-one},
  \begin{align*}
    T_7
    &\le7+7v-(1-v/2)(1-v^{7/2})\\
    &\le6+\frac{15}{2}v+\frac{1}{2}v^3,
  \end{align*}
  because $\sqrt v\le1/2$.  On the other hand,
  \eqref{eq:plt23-ell-small-v} and \eqref{eq:plt23-cxi-small-v-poly} show that the right
  side of \eqref{eq:plt23-linear-core} is at least
  \[
    9\left(1-\frac{7}{8}s^2\right)
    \left(1-\frac{15}{16}s\right)
    \left(1+s^2-\frac{15}{16}s\right)
    (1+5s^2+11s^4).
  \]
  After subtracting $6+15s^2/2+s^6/2$, we obtain the polynomial
  \begin{align}
    P_9(s)={}&\frac{10395}{128}s^9-\frac{333333}{2048}s^8
    +\frac{13635}{128}s^7+\frac{51005}{2048}s^6 \notag\\
    &-\frac{18765}{128}s^5+\frac{264969}{2048}s^4
    -\frac{4995}{64}s^3+\frac{11913}{256}s^2 \notag\\
    &-\frac{135}{8}s+3. \label{eq:plt23-P9-small}
  \end{align}
  By \Cref{lem:bernstein-P9}, $P_9(s)>0$ on $[0,1/2]$.  This proves
  \eqref{eq:plt23-linear-core}.
\end{proof}

\begin{lemma}[The middle-$v$ part of $N=7$]\label{lem:plt23-N7-middle-v}
  Under \eqref{eq:plt23-N7-corner}, if $1/4\le v\le5/8$, then
  \eqref{eq:plt23-linear-core} holds.
\end{lemma}

\begin{proof}
  We first show $\xi>1/3$.  With $t=\zeta+v$,
  \[
    \xi=\frac{4}{v}\frac{(t+1)^2}{t^2(t(v+2)+2)}.
  \]
  The quantity $vG(t,v)$, with $G$ as in the proof of
  \Cref{lem:plt23-N7-small-v}, is increasing in both variables.  Indeed,
  \[
    \frac{\partial(vG)}{\partial t}=v\frac{\partial G}{\partial t}>0,
    \qquad
    \frac{\partial(vG)}{\partial v}
    =\frac{2t^2(tv+t+1)}{(t+1)^2}>0.
  \]
  At
  $(t,v)=(61/8,5/8)$, it equals
  \[
    \frac{26214445}{2437632}<12.
  \]
  Thus $\xi>1/3$, and hence
  \begin{equation}\label{eq:plt23-varphi-7/10}
    \varphi\ge\frac{7}{10}.
  \end{equation}

  Put
  \begin{equation}\label{eq:plt23-y-middle}
    y^2=\frac{7v}{8+v}.
  \end{equation}
  Since $\zeta\le7$, we have $\ell\le y$.  The first half of the
  compensation proof, \eqref{eq:plt23-sqrt-compensation}, gives
  \[
    \sqrt{2\Lambda}\ge1-\ell^2\ge1-y^2.
  \]
  Also, $v^{7/2}\le v/2$ on $[1/4,5/8]$, because
  $v^{5/2}\le(5/8)^{5/2}<1/2$.  Together with
  \eqref{eq:plt23-zeta-simple-lower}, \eqref{eq:plt23-T-bound-one} yields
  \begin{equation}\label{eq:plt23-T7-middle}
    T_7\le6+8v.
  \end{equation}
  By \eqref{eq:plt23-varphi-7/10}, it is enough to prove
  \begin{equation}\label{eq:plt23-middle-target}
    \frac{63}{10}(1-y^2)(1-y)(1+v-y)(1+v)^6
    \ge6+8v.
  \end{equation}

  We verify that the ratio of the two sides is increasing.  Solving
  \eqref{eq:plt23-y-middle} gives
  \[
    v=v(y)=\frac{8y^2}{7-y^2}.
  \]
  For $1/4\le v\le5/8$,
  \[
    \frac{9}{20}<y<\frac{5}{7}.
  \]
  Define
  \[
    \mathcal R(y)=
    \frac{(1-y^2)(1-y)(1+v(y)-y)(1+v(y))^6}
    {6+8v(y)}.
  \]
  A direct logarithmic differentiation gives
  \begin{equation}\label{eq:plt23-R-log-derivative}
    \frac{\mathcal R'(y)}{\mathcal R(y)}
    =\frac{2Q_{10}(y)}
    {(y-1)(y+1)(y^2-7)(y^2+1)(29y^2+21)
    (y^3+7y^2-7y+7)},
  \end{equation}
  where
  \begin{align}
    Q_{10}(y)={}&58y^{10}+319y^9-1793y^8-11280y^7+9788y^6\\
    &-11354y^5+350y^4+5880y^3-6174y^2+5635y-1029.
    \label{eq:plt23-Q10}
  \end{align}
  The denominator in \eqref{eq:plt23-R-log-derivative} is positive on
  $[9/20,5/7]$, and \Cref{lem:bernstein-Q10} gives
  $Q_{10}(y)>0$ there.  Thus $\mathcal R$ is increasing.

  At $v=1/4$, one has $y_0=\sqrt{7/33}<6/13$.  Since the factors involving
  $y$ decrease with $y$,
  \begin{align*}
    &\frac{63}{10}(1-y_0^2)(1-y_0)(5/4-y_0)(5/4)^6\\
    &\quad>
    \frac{63}{10}\frac{26}{33}\frac{7}{13}\frac{41}{52}
    \left(\frac{5}{4}\right)^6
    =\frac{18834375}{2342912}>8.
  \end{align*}
  This proves \eqref{eq:plt23-middle-target} at the left endpoint.  The monotonicity
  of $\mathcal R$ proves it throughout the interval.
\end{proof}

\begin{proposition}[The regime $2\rho\xi>1$]\label{prop:plt23-linear-branch}
  For every odd $N\ge7$, assumption \eqref{eq:plt23-linear-branch} implies
  \[
    R_m\le \Omega_m\psi(\xi,\rho).
  \]
\end{proposition}

\begin{proof}
  For $N\ge9$, if $\zeta\le N$, use
  \Cref{lem:plt23-linear-low-zeta}.  If $\zeta\ge N$, use
  \Cref{lem:plt23-linear-high-zeta} when $v\le5/8$, and
  \Cref{lem:plt23-linear-large-v} when $v\ge5/8$.

  For $N=7$, the same two broad lemmas cover the cases
  $\zeta\ge7$ or $v\ge5/8$.  The remaining corner
  \eqref{eq:plt23-N7-corner} is covered by
  \Cref{lem:plt23-N7-small-v,lem:plt23-N7-middle-v}.
  In every case \eqref{eq:plt23-linear-core} holds.  Combining it with
  \eqref{eq:plt23-linear-payment-start}, \eqref{eq:plt23-KL-first-term}, and
  \eqref{eq:plt23-F-T} proves the proposition.
\end{proof}


\section{Assembly of the complete proof}\label{sec:assembly}

\begin{proof}[Proof of \Cref{thm:main}]
  Put $q=1-p$.

  If $p\le1/2$, the right side of \eqref{eq:main-goal} is nonpositive, so the
  claim follows from $t(C_m,W)\ge0$.  Assume henceforth $p>1/2$.

  The cases $m=3,5,7$ follow at every density from
  \Cref{prop:C3,prop:C5,prop:C7}.  Let $m\ge9$ be odd.

  If $p\ge2/3$, or equivalently $q\le1/3$, the result is
  \Cref{thm:pge23-main}.  It remains to consider
  $1/2<p<2/3$, equivalently $1/3<q<1/2$.
  By \Cref{lem:step-reduction}, it is enough to treat a step graphon.  If the
  compression $A$ has no eigenvalue above $q$,
  \Cref{lem:plt23-eigenvalues-above-q} proves the theorem.  Otherwise there is a unique
  frontier eigenvalue $\Lambda$, and
  \Cref{prop:plt23-splitting-lower-bound,thm:plt23-huber-elim} reduce the problem to the scalar
  target \eqref{eq:plt23-scalar-target}.  If $2\rho\xi\le1$,
  \Cref{prop:plt23-quadratic-branch} proves that target.  If $2\rho\xi>1$,
  \Cref{prop:plt23-linear-branch} proves it.  The two alternatives are exhaustive.
  The step-graphon approximation returns the result to arbitrary graphons.
\end{proof}

\begin{remark}[What is finite and what is uniform]
  The dense-region argument is uniform in every odd $m\ge3$.  The
  intermediate-region argument is uniform in every odd $m\ge9$, except for a
  small corner at $m=9$, where two fixed univariate polynomials are certified
  positive in \Cref{app:bernstein}.  There is no length-by-length computation
  for $m=11,13$, and no adaptive interval subdivision anywhere in the proof.
\end{remark}

\begin{remark}[Equality and stability]
  The manuscript is organized to prove the inequality rather than classify all
  equality cases.  A stability analysis would require tracking near-equality
  in the cycle-to-path deletion, the gamma moment inequality, the
  Hilbert--Schmidt bound, the variance lower bound, and the active Huber
  witness.  The reductions here identify those loci but do not pursue the
  classification.
\end{remark}

\section{Transferable mechanisms and the orthogonal-polynomial viewpoint}
\label{sec:transferable}

This section is not needed for logical completeness.  Its purpose is to
extract methods that may be useful in other spectral extremal problems.

\subsection{One-step stripping: the Schur complement as a Jacobi recursion}

Let $K$ be a finite-dimensional self-adjoint operator, let $e$ be a unit
vector, and decompose
\[
  K=\begin{pmatrix}a&h^*\\h&B\end{pmatrix}
  \quad\text{on}\quad \R e\oplus e^\perp.
\]
The same block inversion used in \eqref{eq:pge23-path-gf} gives
\begin{equation}\label{eq:general-stripping}
  M_{K,e}(z):=\ip e{(I-zK)^{-1}e}
  =\frac{1}{1-az-z^2\ip h{(I-zB)^{-1}h}}.
\end{equation}
If $h\ne0$, let $\nu$ be the probability spectral measure of $B$ at
$h/\norm h$.  Then
\begin{equation}\label{eq:general-stripping-normalized}
  M_{K,e}(z)^{-1}
  =1-az-\norm h^2z^2M_\nu(z).
\end{equation}
This is exactly the first step of the Jacobi continued fraction.  Iterating
on the cyclic Krylov space generated by $e$ produces a tridiagonal Jacobi
matrix and its full J-fraction.  In the centered, variance-one normalization,
\eqref{eq:general-stripping-normalized} becomes
\[
  1-M_\mu(z)^{-1}=z^2M_\nu(z),
\]
the classical once-stripped moment relation emphasized by
Anshelevich~\cite{Anshelevich2009}; coefficient stripping for scalar and
matrix Jacobi operators is surveyed by Damanik--Pushnitski--Simon
\cite{DamanikPushnitskiSimon2008}.

In our graphon proof, $e=\one$, $a=q$, $h=g$, and $B=A$.  Thus the
complement degree fluctuation $g$ is the first off-diagonal Jacobi
coefficient, while $A$ is the once-stripped tail.  The proof does not need
to tridiagonalize the tail; it retains its spectral measure $\mu$ directly.
This is useful in infinite-dimensional problems where the first Jacobi step
is rigid but the tail is not.

\subsection{First returns, irreducible excursions, and logarithmic cycles}

Equation \eqref{eq:pge23-path-gf} has the first-return form
\[
  M(z)=\frac{1}{1-H(z)},
  \qquad H(z)=qz+z^2R_+(z).
\]
Combinatorially, a closed walk from the distinguished state is a sequence of
irreducible returns: either a level step of weight $q$, or a departure into
the orthogonal tail followed by a return, encoded by $z^2R_+$.  This is
the same structural decomposition that underlies Flajolet's continued
fraction for weighted Motzkin paths~\cite{Flajolet1980} and Lehner's use of
Jacobi matrices, lattice paths, and irreducible paths in cumulant
expansions~\cite{Lehner2003}.

The logarithms $-\log(1-Y_\pm)=\sum_{r\ge1}Y_\pm^r/r$ perform a related
but cyclic assembly: $r$ excursions are concatenated with the cyclic
weight $1/r$.  The two-sided identity then applies a parity projection:
for odd coefficients, the stripped-tail terms from $A$ and $-A$ cancel.
A general design principle suggested by this proof is:
\begin{quote}
  Separate a distinguished state, encode returns by a Schur complement, and
  combine a model with a signed or complemented companion so that the
  uncontrolled tail moments cancel in the desired parity class.
\end{quote}
This ``strip, assemble, cancel'' pattern is broader than graphons.

\subsection{Dirichlet diagonalization as probabilistic polarization}

The calculation in \Cref{lem:pge23-dirichlet} is an instance of the following
general lemma.

\begin{proposition}[Dirichlet polarization]\label{prop:general-dirichlet}
  Fix $r\ge1$.  Suppose
  \[
    F(\lambda_1,\ldots,\lambda_r)
    =\sum_{j=0}^d c_jh_j(\lambda_1,\ldots,\lambda_r),
  \]
  and define
  \[
    \widetilde F(\ell)
    =\sum_{j=0}^d c_j\binom{j+r-1}{r-1}\ell^j.
  \]
  If $\Theta\sim\operatorname{Dirichlet}(1,\ldots,1)$, then
  \[
    F(\lambda_1,\ldots,\lambda_r)
    =\E\widetilde F(\Theta\cdot\lambda).
  \]
  Consequently, for every interval $I\subset\R$,
  \[
    \min_{\lambda\in I^r}F(\lambda)
    =\min_{\ell\in I}\widetilde F(\ell).
  \]
\end{proposition}

\begin{proof}
  This is exactly \eqref{eq:pge23-dirichlet-moment}, followed by linearity.
  The minimum statement follows because $\Theta\cdot\lambda\in I$, and the
  reverse inequality is obtained on the diagonal.
\end{proof}

This lemma may be viewed as a positive, probabilistic version of
polarization: the complete homogeneous basis is precisely the basis for
which restriction to the diagonal can be inverted by averaging over a
simplex.  The law of $\Theta\cdot\lambda$ is a Dirichlet average and, in
one dimension, is closely related to B-splines and the Hermite--Genocchi
formula for divided differences; see de~Boor~\cite{deBoor2005}.  The key
extremal consequence is unusually strong: positivity on a whole spectral
cube is equivalent to positivity on the diagonal whenever the polynomial
has this basis structure.

\subsection{Beta--gamma algebra and Laplace positivity transfer}

The passage from \eqref{eq:pge23-beta-integral} to
\eqref{eq:pge23-laplace-gamma} is an instance of a general positive
transform.  For suitable $f$, integers $m>r\ge1$, and $\ell>0$,
\begin{equation}\label{eq:general-laplace-transfer}
  \int_0^\infty\frac{s^{r-1}}{(\ell+s)^m}f(s)\dd s
  =\frac{\Gamma(r)}{\Gamma(m)}
  \int_0^\infty t^{m-r-1}e^{-\ell t}
  \E f(Y/t)\dd t,
  \qquad Y\sim\GammaLaw(r,1).
\end{equation}
The proof is just the Laplace representation of $(\ell+s)^{-m}$, followed
by $y=ts$.  Its value is conceptual: a parameter-dependent rational
kernel becomes a positive mixture of expectations under a fixed
probability family.  Therefore any pointwise-in-$t$ expectation inequality
for $f(Y/t)$ transfers automatically to every $\ell$.

In the present proof, the repeated variables in the complete homogeneous
polynomials first produce a beta distribution; the projective coordinate
$s=\ell(1-x)/x$ turns the beta density into a beta-prime kernel; and the
Laplace transform reveals a gamma variable.  This beta--gamma route is a
useful search heuristic whenever repeated spectral parameters and rational
powers occur together.

\subsection{The gamma Stein identity and a hidden Laguerre equation}

Equation \eqref{eq:pge23-gamma-Stein} is the standard gamma Stein identity; the
operator
\[
  \mathcal A_rf(y)=yf'(y)+(r-y)f(y)
\]
characterizes the $\GammaLaw(r,1)$ law and is widely used in Stein's
method; see, for example, Gaunt~\cite{Gaunt2016}.  Here it is used not for
distributional approximation but as a ladder relation for shifted moments.

There is an exact orthogonal-polynomial shadow.  For an integer $k\ge0$,
put
\[
  F_k(b)=\E(Y-b)^k.
\]
Expansion of the moment, or the Rodrigues formula, gives
\begin{equation}\label{eq:Laguerre-identity}
  F_k(b)=(-1)^k k!\,L_k^{(-r-k)}(-b),
\end{equation}
where $L_k^{(\alpha)}$ is the generalized Laguerre polynomial.  For
$k=2j$, the Laguerre differential equation is exactly
\eqref{eq:pge23-gamma-ODE}.  The parameter $-r-2j$ lies outside the classical
orthogonality range $\alpha>-1$, so we are using the algebraic differential
and ladder structure of Laguerre polynomials, not a positive Laguerre
orthogonality measure.  Standard background is in Szeg\H{o}
\cite{Szego1975}.

The function $G$ in \eqref{eq:pge23-gamma-G} is a weighted amplitude whose
critical points factor explicitly.  Proving that its left maximum dominates
its right maximum and then using \eqref{eq:pge23-gamma-ibp-max} is analogous in
spirit to Sonin--Sturm arguments for comparing successive extrema of
solutions to second-order equations.  The important transferable move is:
construct a first-order multiplier so that the desired linear combination
of a polynomial solution and its derivative becomes an integral of a
factored derivative; positivity is then reduced to a global maximum
comparison.

This explains the connection with the orthogonal-polynomial techniques
suggested by the two cited papers.  Anshelevich's once-stripped moment
identity is literally our first Schur step, while Lehner's Jacobi-matrix and
Motzkin-path picture explains the return decomposition of the moment series.
The later gamma lemma lands in a different classical family---Laguerre---but
uses the same general philosophy: a recurrence or differential equation is
more effective than coefficientwise estimation of high moments.

\subsection{The one-dangerous-mode principle and Huber duality}

The intermediate region exhibits a complementary mechanism.  The square
bound says that two eigenvalues above $q$ would cost more than $pq$ when
$q>1/3$, so there is at most one dangerous mode.  Entrywise constraints on
the corresponding eigenfunction force a choice: either $g$ couples
substantially to that mode, or a quantitatively controlled part of $g$
must lie in the safe subspace.  Retaining both alternatives leads to
\[
  \rho v^2+(\xi-v+v^2)_+,
\]
and minimization over the shape variable $v$ produces the Huber envelope
$\psi(\xi,\rho)$.  Its dual variable $\lambda\in[0,1]$ is an active-set
multiplier.  The witnesses $\lambda=2\rho\xi$ and $\lambda=1$ are the
quadratic and saturated linear regimes.

This suggests a general recipe for constrained spectral inequalities:
\begin{enumerate}[leftmargin=2.2em]
  \item use a trace or energy bound to isolate a finite set of dangerous
  modes;
  \item derive geometric constraints that force compensation either in direct
  coupling or in the safe complement;
  \item keep the channels simultaneously rather than estimating them
  separately;
  \item eliminate the shape by convex duality, and choose dual witnesses
  adapted to small and large activation.
\end{enumerate}
The method should extend to several dangerous modes, where the scalar Huber
envelope is replaced by a low-dimensional convex program.

\subsection{Why \texorpdfstring{$q=1/3$}{q=1/3} is the natural interface}

The same threshold appears for independent structural reasons.

On the dense side, centering at $1/2$ gives
$\delta=1/2-q$.  The gamma moment lemma pays for the odd term in
\eqref{eq:pge23-rho-centered} precisely when
\[
  2\delta\ge\frac{1}{3},
  \qquad\text{i.e.}\qquad q\le\frac{1}{3}.
\]
On the intermediate side, two eigenvalues larger than $q$ would consume
more than $2q^2$ of the bound \eqref{eq:plt23-HS-bound}, and uniqueness is guaranteed
precisely when
\[
  2q^2>pq=q(1-q),
  \qquad\text{i.e.}\qquad q>\frac{1}{3}.
\]
Thus the two mechanisms meet without overlap or gap.  This matching strongly
suggests that $q=1/3$ is an intrinsic phase boundary of the proof problem,
not an artifact of constants chosen later in the argument.
\appendix

\section{The two Bernstein certificates}\label{app:bernstein}

For completeness, this appendix records the only two fixed polynomial
positivity checks used in the proof.  The Bernstein-basis principle used
here is reviewed in Farouki~\cite{Farouki2012}.  We include the coefficients themselves,
not merely a reference to a computer calculation.

For an interval $[a,b]$, put
\[
  t=\frac{x-a}{b-a},
  \qquad
  B^{[a,b]}_{i,n}(x)=\binom ni t^i(1-t)^{n-i}
  \quad(0\le i\le n).
\]
The functions $B^{[a,b]}_{i,n}$ are nonnegative on $[a,b]$ and sum to
one.  Hence an expansion
\[
  P(x)=\sum_{i=0}^n \beta_i B^{[a,b]}_{i,n}(x)
\]
with all $\beta_i>0$ proves $P(x)>0$ throughout the interval.  The
coefficients can be checked directly as follows.  If
\[
  P(a+(b-a)t)=\sum_{j=0}^n c_jt^j,
\]
then
\begin{equation}\label{eq:power-to-bernstein}
  \beta_i=
  \sum_{j=0}^i c_j\frac{\binom ij}{\binom nj}
  \qquad(0\le i\le n).
\end{equation}
Indeed,
\[
  t^j=\frac{1}{\binom nj}
  \sum_{i=j}^n\binom ij B_{i,n}(t),
\]
and substitution gives \eqref{eq:power-to-bernstein}.

\begin{lemma}[The degree-nine certificate]\label{lem:bernstein-P9}
  The polynomial in \eqref{eq:plt23-P9-small}, namely
  \begin{align*}
    P_9(s)={}&
    \frac{10395}{128}s^9
    -\frac{333333}{2048}s^8
    +\frac{13635}{128}s^7
    +\frac{51005}{2048}s^6\\
    &-\frac{18765}{128}s^5
    +\frac{264969}{2048}s^4
    -\frac{4995}{64}s^3
    +\frac{11913}{256}s^2
    -\frac{135}{8}s+3,
  \end{align*}
  is positive on $[0,1/2]$.
\end{lemma}

\begin{proof}
  Its degree-nine Bernstein coefficients on $[0,1/2]$, in increasing order
  of the index, are
  \begin{align*}
    &3,
    &&\frac{33}{16},
    &&\frac{17795}{12288},
    &&\frac{29843}{28672},
    &&\frac{361761}{458752},\\[1mm]
    &\frac{918263}{1376256},
    &&\frac{7142309}{11010048},
    &&\frac{1096445}{1572864},
    &&\frac{390495}{524288},
    &&\frac{361511}{524288}.
  \end{align*}
  Every numerator and denominator is positive.  The Bernstein
  partition-of-unity argument therefore proves the claim.  Formula
  \eqref{eq:power-to-bernstein} gives a short exact-arithmetic verification of
  the displayed list.
\end{proof}

\begin{lemma}[The degree-ten certificate]\label{lem:bernstein-Q10}
  The polynomial $Q_{10}$ in \eqref{eq:plt23-Q10} is positive on
  $[9/20,5/7]$.
\end{lemma}

\begin{proof}
  Its degree-ten Bernstein coefficients on $[9/20,5/7]$, in increasing
  order of the index, are
  \begin{align*}
    &\frac{3243737479716939}{5120000000000},
    &&\frac{24439395309190297}{35840000000000},
    &&\frac{2256629681468091}{3136000000000},\\[1mm]
    &\frac{653354162075429}{878080000000},
    &&\frac{3033296182962901}{4033680000000},
    &&\frac{625950152734243}{847072800000},\\[1mm]
    &\frac{1151767076423}{1647086000},
    &&\frac{6182915802683}{9882516000},
    &&\frac{7481469276}{14706125},\\[1mm]
    &\frac{67938971678}{201768035},
    &&\frac{26748047904}{282475249}.
  \end{align*}
  Again every coefficient is positive.  Since the Bernstein basis is
  nonnegative and sums to one, $Q_{10}>0$ on the stated interval.
  The conversion formula \eqref{eq:power-to-bernstein} verifies the list using
  only rational arithmetic.
\end{proof}

\begin{remark}[Why this is a proof certificate]
  A Bernstein-coefficient list is not a numerical sample.  Once the exact
  identity with the Bernstein expansion is checked, positivity at every point
  of the whole interval follows from the nonnegativity of the basis functions.
  In the present proof no subdivision is used: one interval suffices for each
  of the two polynomials.
\end{remark}


\begin{thebibliography}{99}

  \bibitem{Anshelevich2009}
  M.~Anshelevich,
  \newblock Appell polynomials and their relatives III: conditionally free theory,
  \newblock \emph{Illinois Journal of Mathematics} \textbf{53} (2009), 39--66;
  \newblock arXiv:0803.4279.

  \bibitem{deBoor2005}
  C.~de~Boor,
  \newblock Divided differences,
  \newblock \emph{Surveys in Approximation Theory} \textbf{1} (2005), 46--69.

  \bibitem{DamanikPushnitskiSimon2008}
  D.~Damanik, A.~Pushnitski, and B.~Simon,
  \newblock The analytic theory of matrix orthogonal polynomials,
  \newblock \emph{Surveys in Approximation Theory} \textbf{4} (2008), 1--85.

  \bibitem{Farouki2012}
  R.~T. Farouki,
  \newblock The Bernstein polynomial basis: a centennial retrospective,
  \newblock \emph{Computer Aided Geometric Design} \textbf{29} (2012),
  379--419.

  \bibitem{Flajolet1980}
  P.~Flajolet,
  \newblock Combinatorial aspects of continued fractions,
  \newblock \emph{Discrete Mathematics} \textbf{32} (1980), 125--161.

  \bibitem{Gaunt2016}
  R.~E. Gaunt,
  \newblock Products of normal, beta and gamma random variables: Stein
  operators and distributional theory,
  \newblock \emph{Brazilian Journal of Probability and Statistics}
  \textbf{32} (2018), 437--466;
  \newblock arXiv:1507.07696.

  \bibitem{Goodman1959}
  A.~W. Goodman,
  \newblock On sets of acquaintances and strangers at any party,
  \newblock \emph{American Mathematical Monthly} \textbf{66} (1959),
  778--783.

  \bibitem{Huber1964}
  P.~J. Huber,
  \newblock Robust estimation of a location parameter,
  \newblock \emph{Annals of Mathematical Statistics} \textbf{35} (1964),
  73--101.

  \bibitem{Lehner2003}
  F.~Lehner,
  \newblock Cumulants, lattice paths, and orthogonal polynomials,
  \newblock \emph{Discrete Mathematics} \textbf{270} (2003), 177--191;
  \newblock arXiv:math/0110031.

  \bibitem{Lovasz2012}
  L.~Lov\'asz,
  \newblock \emph{Large Networks and Graph Limits},
  \newblock American Mathematical Society Colloquium Publications, vol.~60,
  American Mathematical Society, Providence, RI, 2012.

  \bibitem{PowersReznick2000}
  V.~Powers and B.~Reznick,
  \newblock Polynomials that are positive on an interval,
  \newblock \emph{Transactions of the American Mathematical Society}
  \textbf{352} (2000), 4677--4692.

  \bibitem{Sion1958}
  M.~Sion,
  \newblock On general minimax theorems,
  \newblock \emph{Pacific Journal of Mathematics} \textbf{8} (1958),
  171--176.

  \bibitem{Szego1975}
  G.~Szeg\H{o},
  \newblock \emph{Orthogonal Polynomials}, fourth edition,
  \newblock American Mathematical Society Colloquium Publications, vol.~23,
  American Mathematical Society, Providence, RI, 1975.

\end{thebibliography}

\end{document}
